%MIT License
%
%Copyright (c) 2018 Chen Wang [https://chenwang.site]
%
%Permission is hereby granted, free of charge, to any person obtaining a copy
%of this software and associated documentation files (the "Software"), to deal
%in the Software without restriction, including without limitation the rights
%to use, copy, modify, merge, publish, distribute, sublicense, and/or sell
%copies of the Software, and to permit persons to whom the Software is
%furnished to do so, subject to the following conditions:
%
%The above copyright notice and this permission notice shall be included in all
%copies or substantial portions of the Software.
%
%THE SOFTWARE IS PROVIDED "AS IS", WITHOUT WARRANTY OF ANY KIND, EXPRESS OR
%IMPLIED, INCLUDING BUT NOT LIMITED TO THE WARRANTIES OF MERCHANTABILITY,
%FITNESS FOR A PARTICULAR PURPOSE AND NONINFRINGEMENT. IN NO EVENT SHALL THE
%AUTHORS OR COPYRIGHT HOLDERS BE LIABLE FOR ANY CLAIM, DAMAGES OR OTHER
%LIABILITY, WHETHER IN AN ACTION OF CONTRACT, TORT OR OTHERWISE, ARISING FROM,
%OUT OF OR IN CONNECTION WITH THE SOFTWARE OR THE USE OR OTHER DEALINGS IN THE
%SOFTWARE.
\documentclass[12pt,a4paper]{Thesis} % Paper size, default font size and one-sided paper

\graphicspath{%
	{./Pictures/}%
	{./Figures/}%
}
\DeclareMathOperator{\Tr}{Tr}
\let\savedegree\degree
\let\degree\relax
\let\savedegree\ref
\let\rem\relax
\usepackage{adjustbox}
\usepackage{amsmath}
\usepackage[amssymb]{SIunits}
\usepackage{amssymb}
\usepackage{multirow}
\usepackage{wrapfig}
\usepackage{enumitem}
\usepackage{subcaption}
\usepackage{mathtools}
\usepackage{lipsum}
\usepackage{algorithm}
\usepackage{algpseudocode}
\usepackage[square, numbers, comma, sort&compress]{natbib} % Use the natbib reference package - read up on this to edit the reference style; if you want text (e.g. Smith et al., 2012) for the in-text references (instead of numbers), remove 'numbers'

\title{\ttitle} % Defines the thesis title - don't touch this

%extra packages
\usepackage{float}
\usepackage{color,soul}               % highlighting text
\usepackage{enumerate}
\usepackage{Styles/mydefs}

%---------------------------------------------------------------------------
% Additional packages and macros required by the compiled research chapters
% (Chapters 3-5), ported from the original conference manuscripts.
%---------------------------------------------------------------------------
\usepackage{xcolor}        % coloured table entries (best/second/third best)
\usepackage{array}         % extended column specifications
\usepackage{booktabs}      % \toprule, \midrule, \bottomrule, \cmidrule
\usepackage{colortbl}      % cell shading in tables
\usepackage{makecell}      % multi-line table cells
\usepackage{bm}            % bold math symbols
\usepackage{threeparttable}% table notes

% Ranking colours reused across the experimental tables.
\newcommand{\fst}[1]{\textcolor[RGB]{255,0,185}{#1}}   % best
\newcommand{\snd}[1]{\textcolor[RGB]{0,153,76}{#1}}    % second best
\newcommand{\trd}[1]{\textcolor[RGB]{0,128,255}{#1}}   % third best

% Math operators used in the methods.
\DeclareMathOperator*{\argmax}{arg\,max}
\DeclareMathOperator*{\argmin}{arg\,min}
\DeclareMathOperator{\corr}{corr}
\DeclareMathOperator{\Var}{Var}
\DeclareMathOperator{\Cov}{Cov}


\begin{document}

\frontmatter % Use roman page numbering style (i, ii, iii, iv...) for the pre-content pages

\setstretch{1.3} % Line spacing of 1.3
% Define the page headers using the Fr package and set up for one-sided printing

\pagestyle{fancy} % Finally, use the "fancy" page style to implement the FancyHdr headers
\fancyhead{}   %clear all fields
\fancyhead[LO]{\sl{\leftmark}}
\fancyhead[RE]{\sl{\rightmark}}
\fancyhead[LE,RO]{\thepage}


%----------------------------------------------------------------------------------------
%	TITLE PAGE
%----------------------------------------------------------------------------------------

%%% Use \maketitleforreview instead of \maketitle, if you want a PLAIN TITLE PAGE

\maketitle
% \maketitleforreview


% ----------------------------------------------------------------------------------------
% 	DECLARATION PAGE
% ----------------------------------------------------------------------------------------

\thesisdeclareOriginality{I hereby certify that the work embodied in this thesis is the result of original research, is free of plagiarised materials, and has not been submitted for a higher degree to any other University or Institution.}{Aug. 2018}{Styles/signature.png}

\thesisdeclareSupervisor{I have reviewed the content and presentation style of this thesis and declare it is free of plagiarism and of sufficient grammatical clarity to be examined.  To the best of my knowledge, the research and writing are those of the candidate except as acknowledged in the Author Attribution Statement. I confirm that the investigations were conducted in accord with the ethics policies and integrity standards of Nanyang Technological University and that the research data are presented honestly and without prejudice.}{Mar. 2019}{Styles/signature.png}

\thesisdeclareAuthorship{This thesis contains material from three papers in which I am listed as an author. The works reproduced in the technical chapters were carried out during my doctoral candidature.}
{The technical chapters of this thesis are based on the following papers, with placeholder author lists to be completed by the candidate:

	Chapter~\ref{ch:earnhft} is based on {\color{blue}M.~Qin et al., ``EarnHFT: Efficient Hierarchical Reinforcement Learning for High Frequency Trading,'' in \emph{Proceedings of the AAAI Conference on Artificial Intelligence}, 2024.}

	Chapter~\ref{ch:fineft} is based on {\color{blue}M.~Qin et al., ``FineFT: Efficient and Risk-Aware Ensemble Reinforcement Learning for Futures Trading,'' manuscript submitted to \emph{ACM SIGKDD}, 2026.}

	Chapter~\ref{ch:alphaqd} is based on {\color{blue}M.~Qin et al., ``AlphaQD: Quality-Diversity Reinforcement Learning for Formulaic Alpha Mining,'' manuscript submitted to \emph{ACM SIGKDD}.}

	The contributions of the candidate and co-authors are as follows:
	\begin{itemize}[topsep=1pt,itemsep=1pt,partopsep=1pt, parsep=1pt]
		\item I am the leading author of all three papers. I formulated the research problems, designed the methods, implemented the algorithms and the high-fidelity trading environments, conducted all experiments, and wrote the manuscripts.
		\item My supervisor provided overall research direction and revised the manuscripts.
		\item Co-authors contributed through technical discussions, partial baseline implementation, and proofreading of the manuscripts.
	\end{itemize}
}{[Month, Year]}{Styles/signature.png}

% ----------------------------------------------------------------------------------------
% 	ACKNOWLEDGEMENTS
% ----------------------------------------------------------------------------------------

\setstretch{1.3} % Reset the line-spacing to 1.3 for body text (if it has changed)

\acknowledgements{\addtocontents{toc}{\vspace{0.8em}} % Add a gap in the Contents, for aesthetics

I wish to express my deepest gratitude to my supervisor for the guidance, patience, and intellectual freedom that made this research possible, and for teaching me to look for the simplest structure behind a difficult problem. I am grateful to my collaborators and co-authors on the works that constitute this thesis, whose discussions sharpened every idea presented here, and to my colleagues in the laboratory for the day-to-day exchange that sustains research. I thank the thesis examination committee for their time and their careful reading. Finally, I thank my family for their unwavering support throughout my doctoral studies.

% \begin{flushright}
% \emph{Molei Qin}
% \end{flushright}
}

%----------------------------------------------------------------------------------------
%	QUOTATION PAGE
%----------------------------------------------------------------------------------------

\pagestyle{empty} % No headers or footers for the following pages
\emph{``If I had one hour to save the world, I would spend 55 minutes defining the problem and only five minutes finding the solution."}

\begin{flushright}
---Einstein, Albert
\end{flushright}
\null\vfill % Add some space to move the quote down the page a bit

\begin{center}
\large{To my dear family}
\end{center}
\vfill\vfill\null % Add some space at the bottom to position the quote just right
\cleardoublepage % Start a new page and make the next page a right-hand (odd-numbered) page, producing a blank page if necessary

\pagestyle{fancy} % Return the page headers back to the "fancy" style
\fancyhead{}   %clear all fields
\fancyhead[LO]{\sl{\leftmark}}
\fancyhead[RE]{\sl{\rightmark}}
\fancyhead[LE,RO]{\thepage}

%----------------------------------------------------------------------------------------
%	ABSTRACT PAGE
%----------------------------------------------------------------------------------------

\addtotoc{Abstract} % Add the "Abstract" page entry to the Contents

\abstract{\addtocontents{toc}{\vspace{0.8em}}}% Add a gap in the Contents, for aesthetics

% 1. The Problem 2. The Issues 3. The Approach 4. The Results 5. The Experiments 6. The Applications

Quantitative trading converts market data into trading decisions and is, at its core, a sequential decision-making problem under uncertainty. Reinforcement learning (RL) is therefore a natural modelling tool, yet financial markets violate the assumptions on which standard RL rests: they are non-stationary, their reward signals carry a low signal-to-noise ratio, the decision horizons can be extreme, and a single mistake can be financially catastrophic. This thesis studies what RL can and cannot do across the breadth of the quantitative-trading stack. It is organised around three tasks that span the trading-frequency spectrum and progress from trade execution to signal generation: second-level high-frequency trading of a single spot asset, minute-level trading of leveraged cryptocurrency futures, and daily-frequency formulaic alpha mining over a full cross-section of equities. For each task, the thesis isolates the specific obstacle that prevents off-the-shelf RL from succeeding and develops a method that removes it.

The first contribution, \textbf{EarnHFT}, addresses high-frequency trading, where the trajectory length (of the order of one million steps per month) cripples data efficiency and rapid regime changes break any single policy. EarnHFT is a three-stage hierarchical RL framework that computes an optimal action-value teacher by dynamic programming to accelerate and stabilise second-level training, assembles a pool of trend-specialised agents, and trains a minute-level router that switches among them. It exceeds the strongest baseline by at least 30\% in profitability across four cryptocurrency markets. The second contribution, \textbf{FineFT}, moves to leveraged futures, where leverage amplifies the aleatoric noise of the reward and where an agent confident outside its competence can be liquidated. FineFT is a three-stage ensemble framework that trains specialists by a selective-update rule driven by an ensemble temporal-difference matrix, learns the capability boundary of each specialist with per-dynamic variational autoencoders, and routes between specialists and a conservative fallback by out-of-distribution detection. It reduces risk by more than 40\% while remaining the most profitable learned method. The third contribution, \textbf{AlphaQD}, turns from execution to signal generation: it diagnoses the boosting-style reward of pool-based RL alpha miners as a source of reward non-stationarity and of an absent generalisation guarantee, and replaces it with a quality-diversity formulation in which a GFlowNet is trained against an archive-independent reward while a MAP-Elites archive, keyed on the normalised daily information-coefficient series as a return-source descriptor, organises diversity and selection. On full-market Chinese A-share and U.S.\ benchmarks it attains the leading signal-quality cluster on both markets with the lowest seed-to-seed variance and turnover among learned pool-style methods.

Read together, the three works trace a single methodological pattern --- decompose a non-stationary market into stationary sub-problems, learn a set of specialists, and compose them through a selector --- and a clear lesson about where the selector belongs: coupling the act of selection into the learning target re-introduces the non-stationarity it was meant to cure, so the selector is progressively decoupled from the policy reward across the three chapters, from a learned online router, to a generative risk-aware router, to a fully structural archive. The thesis thus offers both a set of state-of-the-art methods for distinct trading tasks and a unifying perspective on how reinforcement learning should be structured for non-stationary financial markets.



%----------------------------------------------------------------------------------------
%	LIST OF CONTENTS/FIGURES/TABLES PAGES
%----------------------------------------------------------------------------------------

\pagestyle{fancy} % The page style headers have been "empty" all this time, now use the "fancy" headers as defined before to bring them back

\tableofcontents % Write out the Table of Contents

\listoffigures % Write out the List of Figures

\listoftables % Write out the List of Tables


%----------------------------------------------------------------------------------------
%	ABBREVIATIONS
%----------------------------------------------------------------------------------------

%\clearpage % Start a new page
%
%\setstretch{1.5} % Set the line spacing to 1.5, this makes the following tables easier to read
%
%\lhead{\emph{Abbreviations}} % Set the left side page header to "Abbreviations"
%\listofsymbols{ll} % Include a list of Abbreviations (a table of two columns)
%{
%\textbf{LAH} & \textbf{L}ist \textbf{A}bbreviations \textbf{H}ere \\
%%\textbf{Acronym} & \textbf{W}hat (it) \textbf{S}tands \textbf{F}or \\
%}

%----------------------------------------------------------------------------------------
%	PHYSICAL CONSTANTS/OTHER DEFINITIONS
%----------------------------------------------------------------------------------------

%\clearpage % Start a new page
%
%\lhead{\emph{Physical Constants}} % Set the left side page header to "Physical Constants"
%
%\listofconstants{lrcl} % Include a list of Physical Constants (a four column table)
%{
%Speed of Light & $c$ & $=$ & $2.997\ 924\ 58\times10^{8}\ \mbox{ms}^{-\mbox{s}}$ (exact)\\
%% Constant Name & Symbol & = & Constant Value (with units) \\
%}

%----------------------------------------------------------------------------------------
%	SYMBOLS
%----------------------------------------------------------------------------------------

\cleardoublepage  % Start a new page (no need if the previous page is list of figures)

\listofnomenclature{ll} % Include a list of Symbols (a three column table)
{
\multicolumn{2}{l}{\LARGE{\textbf{Symbols}}}\\[0.4cm]
$\mathcal{S},\ \mathcal{A}$        & the state space and the action space of an MDP\\
$\pi_\theta$                       & a policy parameterised by $\theta$\\
$Q(s,a)$                           & the action-value (Q) function\\
$Q^*(s,a)$                         & the optimal action-value function\\
$\gamma$                           & the discount factor\\
$r_t$                              & the reward received at time $t$\\
$V_t$                              & the net value (margin balance) of the account at time $t$\\
$H_t,\ L_t$                        & the position and the leverage at time $t$\\
$\mathcal{H}(\cdot)$               & the Huber loss\\
$\mathrm{KL}(\cdot\,\|\,\cdot)$    & the Kullback--Leibler divergence\\
$\alpha$                           & a formulaic alpha (symbolic expression)\\
$\mathrm{IC}_t(\alpha)$            & the daily cross-sectional information coefficient at time $t$\\
$\mathrm{ICIR}(\alpha)$            & the information ratio of the IC series\\
$\mathbf{r}(\alpha)$               & the centred, unit-normalised IC series (return-source descriptor)\\
$\beta$                            & a preference (EarnHFT) or reward-temperature (AlphaQD) coefficient\\[0.6cm]

\multicolumn{2}{l}{\LARGE{\textbf{Acronyms}}}\\[0.4cm]
RL          & Reinforcement Learning\\
HFT         & High-Frequency Trading\\
MDP         & Markov Decision Process\\
DP-MDP      & Dynamic Parametric Markov Decision Process\\
HRL         & Hierarchical Reinforcement Learning\\
DQN         & Deep Q-Network\\
DDQN        & Double Deep Q-Network\\
PPO         & Proximal Policy Optimisation\\
TD          & Temporal Difference\\
ETD         & Ensemble Temporal Difference\\
MoE         & Mixture of Experts\\
VAE         & Variational Autoencoder\\
OOD         & Out-of-Distribution\\
ID          & In-Distribution\\
QD          & Quality-Diversity\\
GFlowNet    & Generative Flow Network\\
TB          & Trajectory Balance\\
RGCN        & Relational Graph Convolutional Network\\
MAP-Elites  & Multi-dimensional Archive of Phenotypic Elites\\
PFS         & Perturbation Fidelity Score\\
PCA         & Principal Component Analysis\\
RPN         & Reverse Polish Notation\\
GP          & Genetic Programming\\
DTW         & Dynamic Time Warping\\
LOB         & Limit Order Book\\
OHLC        & Open-High-Low-Close\\
IC          & Information Coefficient\\
ICIR        & Information-Coefficient Information Ratio\\
MACD        & Moving Average Convergence Divergence\\
IV          & Imbalance Volume\\
TR          & Total Return\\
ASR         & Annualised Sharpe Ratio\\
ACR         & Annualised Calmar Ratio\\
ASoR        & Annualised Sortino Ratio\\
MDD         & Maximum Drawdown\\
AVOL        & Annualised Volatility\\
Crypto      & Cryptocurrency\\
}

%----------------------------------------------------------------------------------------
%	THESIS CONTENT - CHAPTERS
%----------------------------------------------------------------------------------------

\mainmatter       % Begin numeric (1,2,3...) page numbering
\pagenumbering{arabic}
\setstretch{1.3}  % Return the line spacing back to 1.3

% define the headings for the body of the thesis
\fancyhead{}   %clear all fields
\fancyhead[LO]{\sl{\leftmark}}
\fancyhead[RE]{\sl{\rightmark}}
\fancyhead[LE,RO]{\thepage}


% Include the chapters of the thesis as separate files from the Chapters folder
% Uncomment the lines as you write the chapters

\part{Preliminaries}
%!TEX root=../mythesis.tex
% Chapter 1

\chapter{Introduction} % Main chapter title
\chaptermark{Introduction}
\label{ch:introduction} % For referencing the chapter elsewhere, use \cref{ch:introduction}

%----------------------------------------------------------------------------------------
\section{Quantitative Trading and Reinforcement Learning}
\label{sec:intro_qt_rl}

Quantitative trading is the discipline of converting market data into trading decisions through explicit, data-driven rules. A quantitative strategy ingests a stream of prices, volumes, and order-flow signals, and at each decision point it must choose what to hold and how to adjust that holding. Stated this way, quantitative trading is, at its core, a problem of \emph{sequential decision-making under uncertainty}: the present action changes the portfolio, the portfolio interacts with a stochastic and adversarial market, and the consequences are realised only over a horizon of subsequent steps. The natural formal language for such problems is reinforcement learning (RL), in which an agent learns a policy that maps states to actions so as to maximise a cumulative reward.

Reinforcement learning is attractive for trading for reasons that distinguish it from the supervised learning that dominates much of quantitative finance. Supervised approaches predict a label---a future return, a price direction---under the assumption that training samples are independent and identically distributed (IID), and they optimise a one-step predictive loss that is agnostic to how the prediction is used. Trading violates both premises. Transaction costs and inventory constraints make the profitability of an action depend on the current position, so consecutive decisions are coupled rather than independent; and the quantity a trader actually cares about is the long-horizon, risk-adjusted profit of a \emph{sequence} of decisions, not the accuracy of any single forecast. Reinforcement learning addresses both directly: it optimises a return that accumulates over the trajectory, and it models the path dependence induced by costs and positions as part of the environment dynamics. Encouraged by RL's successes in games and robotics, a substantial body of work has applied it across the quantitative-trading stack, from order execution and intraday trading to portfolio management~\cite{sun2022deepscalper,briola2021deep,zhu2022quantitative}.

The premise of this thesis is therefore that reinforcement learning is the right tool for quantitative trading; its subject is the gap between that premise and practice. Financial markets violate, sometimes severely, the assumptions on which standard RL algorithms were designed and validated. The contribution of this thesis is to identify, for several distinct and economically important trading tasks, exactly which assumption breaks, and to design a method that restores the broken assumption rather than papering over its symptoms.

%----------------------------------------------------------------------------------------
\section{Why Reinforcement Learning is Hard in Financial Markets}
\label{sec:intro_challenges}

The difficulties that financial markets pose for reinforcement learning recur across tasks, although each task stresses a different subset of them. It is useful to enumerate them once, because the technical chapters can then be read as attacks on specific items in this list.

\begin{itemize}
\item \textbf{Non-stationarity.} The defining property of financial markets is that their statistics drift. The relationship between observable features and future returns changes over time as market regimes shift between trends, ranges, and crises. This breaks the stationary-MDP assumption underpinning convergence guarantees for value- and policy-based RL: a policy that is optimal on the training distribution can be actively harmful once the regime changes, because the environment the agent now faces is not the one it was trained on.

\item \textbf{Low signal-to-noise ratio.} The predictable component of asset returns is small relative to the noise. The reward an RL agent receives is therefore a faint signal buried in stochastic price movement, which slows learning, inflates the variance of gradient estimates, and makes it easy to overfit idiosyncratic patterns of the training window that do not generalise.

\item \textbf{Extreme decision horizons.} At high frequency, the number of decision steps over an evaluation period is enormous. A second-level agent assessed over a few weeks faces a horizon on the order of one million steps, against which the few-thousand-step horizons of standard RL benchmarks~\cite{mnih2013playing} are small. Long horizons demand far more data to converge~\cite{zhang2023towards} and make both training and evaluation computationally heavy.

\item \textbf{Risk and irreversibility.} A trading loss is real and, under leverage, can be catastrophic: a single excursion into an unfamiliar market state can liquidate an account. Unlike a game, where a poor episode merely lowers the average score, trading has hard, asymmetric downside that an expected-return objective alone does not capture.

\item \textbf{Reward design and coupling.} What constitutes a good reward is itself non-trivial. When the objective concerns a \emph{set} of strategies or signals---for example, the marginal contribution of a new factor to an existing pool---it is tempting to fold the pool into the reward. Doing so couples the learning target to a quantity that itself evolves during training, silently re-introducing non-stationarity through the reward.

\item \textbf{Generalisation.} The training window is finite and the future is non-stationary, so a discovered strategy or signal must hold up out of sample. This places a premium on methods whose effective complexity is small relative to the search space, and on selection rules that do not merely fit the training window.
\end{itemize}

No single method removes all of these obstacles at once, nor should one expect it to: the obstacles trade off against one another, and the binding constraint differs by task. This thesis instead studies three concrete tasks, isolates the obstacle that is decisive for each, and removes it.

%----------------------------------------------------------------------------------------
\section{Tasks, Challenges, and Solutions}
\label{sec:intro_tasks}

The three studies that constitute this thesis are organised along the spectrum of trading frequency and, with it, a progression from \emph{trade execution} to \emph{signal generation}. \tref{tab:intro_overview} summarises them; the remainder of this section introduces each in turn.

\begin{table}[!htb]
\centering
\caption[Overview of the three studies.]{The three trading tasks studied in this thesis, the decisive obstacle each presents to standard reinforcement learning, and the method developed in response. The tasks span the frequency spectrum and progress from execution to signal generation.}
\label{tab:intro_overview}
\renewcommand{\arraystretch}{1.3}
\resizebox{\textwidth}{!}{
\begin{tabular}{p{2.1cm} p{2.6cm} p{4.6cm} p{4.8cm}}
\toprule
\textbf{Chapter} & \textbf{Task} & \textbf{Decisive challenge} & \textbf{Method} \\
\midrule
\cref{ch:earnhft} (EarnHFT) & Second-level spot high-frequency trading & Extreme horizon ($\sim 10^6$ steps) cripples data efficiency; regime changes break any single policy & Three-stage \emph{hierarchical} RL: dynamic-programming Q-teacher, a pool of trend specialists, and a minute-level router \\
\cref{ch:fineft} (FineFT) & Minute-level leveraged futures trading & Leverage amplifies reward stochasticity; agents are unaware of their competence boundary and are liquidated on unseen states & Three-stage \emph{ensemble} RL: selective-update specialists, per-dynamic VAE boundaries, and risk-aware routing with a conservative fallback \\
\cref{ch:alphaqd} (AlphaQD) & Daily formulaic alpha mining over a full cross-section & Boosting-style rewards couple the target to an evolving pool, causing non-stationarity and offering no generalisation handle & Three-stage \emph{quality-diversity} RL: GFlowNet exploration under an archive-independent reward, MAP-Elites archiving, and top-$k$ cell export \\
\bottomrule
\end{tabular}}
\end{table}

\tref{tab:intro_spine} makes the thesis spine explicit by tracing, across the three studies, what is decomposed, what is specialised, how the selector is built, and how tightly it is coupled to the policy's learning target.

\begin{table}[!htb]
\centering
\caption[Thesis spine: decomposition, specialist, selector.]{The thesis spine across the three studies. The decomposition criterion changes with the task, but the response to non-stationarity is the same---decompose, specialise, and select---and the selector is progressively decoupled from the policy's learning target.}
\label{tab:intro_spine}
\renewcommand{\arraystretch}{1.3}
\resizebox{\textwidth}{!}{
\begin{tabular}{p{2.0cm} p{3.2cm} p{3.0cm} p{4.0cm} p{3.5cm}}
\toprule
\textbf{Chapter} & \textbf{Decomposition criterion} & \textbf{Specialist} & \textbf{Selector mechanism} & \textbf{Coupling to policy reward} \\
\midrule
\cref{ch:earnhft} (EarnHFT) & Market trend (slope-based segmentation) & Second-level DDQN agent per trend & Minute-level DDQN router trained by RL on historical data and frozen for deployment & \emph{Tightly coupled}: the selector is itself a reward-maximising policy \\
\cref{ch:fineft} (FineFT) & Latent dynamic of a dynamic-parametric MDP & Selectively-updated learner per dynamic & Heuristic router driven by per-dynamic VAE losses and a conservative fallback & \emph{Partly decoupled}: selection is a generative uncertainty estimate, not a learned policy \\
\cref{ch:alphaqd} (AlphaQD) & Return-source behaviour (normalised daily IC series) & Per-cell elite of a MAP-Elites archive & Static top-$k$ cell export with a shared ridge combiner & \emph{Fully decoupled}: the generator's reward never reads the archive, and the archive never edits the reward \\
\bottomrule
\end{tabular}}
\end{table}

\paragraph{Task I: High-frequency trading (\cref{ch:earnhft}).} The first task is second-level high-frequency trading (HFT) of a single cryptocurrency, where the agent repeatedly adjusts its position to profit from short-lived price fluctuations. Two obstacles dominate. First, the horizon is extreme: evaluating a second-level agent over weeks entails roughly one million steps per month, so naive RL is hopelessly data-inefficient. Second, the market is non-stationary at a scale that defeats any single policy, because strategies that profit in a rising market lose in a falling one. EarnHFT resolves both through hierarchy. At high frequency a single trader's orders do not move prices, so future prices are exogenous to the policy. EarnHFT exploits this structural feature: it computes, by dynamic programming over future prices, an optimal action-value teacher that accelerates and stabilises second-level training. It then trains a \emph{pool} of agents, each specialised to a market trend, and a minute-level \emph{router} that selects among them as conditions change. EarnHFT exceeds the strongest baseline by at least 30\% in profitability across four markets.

\paragraph{Task II: Leveraged futures (\cref{ch:fineft}).} The second task moves from spot to leveraged futures, introducing long and short positions and high leverage. Leverage changes the nature of the difficulty in two ways. It amplifies the aleatoric noise of the reward, so that nearly identical states under identical actions can yield wildly different outcomes, undermining convergence; and it raises the cost of error to the point of liquidation, which exposes a deficiency that low-leverage methods can ignore---the absence of any \emph{self-awareness of competence}. An agent that acts confidently outside the market states it was trained on can be wiped out by a single black-swan move. FineFT addresses both. A selective-update rule assigns each transition to the learner that already explains it best, together with its neighbours, training a diverse ensemble efficiently. A per-dynamic variational autoencoder then learns the \emph{capability boundary} of each specialist. At deployment, FineFT routes between the specialists and a conservative fallback by detecting when the current market state is out of distribution. The method reduces risk---measured by maximum drawdown---by more than 40\% while remaining the most profitable learned baseline.

\paragraph{Task III: Formulaic alpha mining (\cref{ch:alphaqd}).} The third task departs from execution altogether. Rather than deciding how to trade, it asks an agent to \emph{generate} the predictive signals---short, interpretable, symbolic formulas, or ``alphas''---that a downstream portfolio is built upon. Here the binding obstacle is one of reward design. The dominant line of RL alpha miners scores a new formula by its marginal contribution to a running pool, which is a boosting-style objective: the reward changes every time the pool changes, so the policy is trained against a moving target, and the resulting ensemble carries no a-priori generalisation guarantee. AlphaQD diagnoses this coupling and removes it. The generator is a graph-based GFlowNet trained against a reward that depends only on the candidate and a fixed training window, never on the archive. Diversity is relocated into a separate quality-diversity archive. Its behaviour descriptor is the normalised daily information-coefficient series of each formula, an empirical proxy for the formula's return source. On full-market Chinese A-share and U.S.\ benchmarks AlphaQD attains the leading signal-quality cluster on both markets with the lowest seed-to-seed variance and turnover among learned pool-style methods.

%----------------------------------------------------------------------------------------
\section{A Recurring Methodological Pattern}
\label{sec:intro_pattern}

Although the three studies were motivated by different tasks and developed independently, they share a methodological skeleton that is worth naming, because it organises the thesis. In each case the response to non-stationarity is the same three-move pattern:

\begin{enumerate}
\item \textbf{Decompose} the non-stationary market into sub-problems that are approximately stationary---by market trend in EarnHFT, by latent dynamic in FineFT, by return-source behaviour in AlphaQD.
\item \textbf{Specialise} a learner to each sub-problem, producing a diverse pool of experts rather than a single monolithic policy.
\item \textbf{Select} or compose among the specialists at deployment through a selector that responds to the prevailing conditions.
\end{enumerate}

The thesis does not impose this pattern as dogma; it emerges because decomposition is the most direct way to convert a non-stationary problem into a family of problems on which RL's stationary-MDP assumption approximately holds. What does deserve emphasis is a lesson that surfaces only when the three selectors are compared. The selector is the component that decides which specialist acts, and the natural temptation is to \emph{learn} it jointly with the specialists, or to fold the selection criterion into the learning reward. Both couplings re-introduce, through the back door, the very non-stationarity the decomposition was meant to remove: a learned router changes the effective environment of the specialists as it learns, and a selection-aware reward changes the target as the pool evolves. Across the three chapters the selector is therefore progressively \emph{decoupled} from the policy reward. In EarnHFT the router is itself trained by reinforcement learning against the trading reward, offline on historical data and frozen before deployment. FineFT replaces this with a generative, risk-aware heuristic router that requires no policy learning at all. AlphaQD goes further: the selection rule is a structural archive whose membership never enters the generator's reward. This decoupling is the connective thread of the thesis and the single most transferable insight it offers. Stated as a single claim:
\begin{quote}
\itshape In a decompose--specialise--select design for a non-stationary environment, the selector should be kept out of the policy's learning target.
\end{quote}

%----------------------------------------------------------------------------------------
\section{Contributions}
\label{sec:intro_contributions}

The contributions of this thesis are as follows.

\begin{itemize}
\item \textbf{EarnHFT} (\cref{ch:earnhft}): the first, to our knowledge, hierarchical reinforcement-learning framework for single-asset high-frequency trading. Its technical core is an optimal action-value teacher computed by dynamic programming over future prices, which is shown both empirically and theoretically to accelerate and stabilise second-level training; together with trend-conditioned specialist training and a minute-level router, it delivers a 30\% profitability margin over the strongest baseline.

\item \textbf{FineFT} (\cref{ch:fineft}): a risk-aware ensemble framework for leveraged futures. It introduces a selective-update rule, driven by an ensemble temporal-difference matrix, that trains diverse specialists efficiently and segments the market automatically; and it is the first method to use variational autoencoders for out-of-distribution detection in RL for quantitative trading, turning competence-awareness into an explicit, generative risk control that reduces risk by more than 40\%.

\item \textbf{AlphaQD} (\cref{ch:alphaqd}): a quality-diversity reinforcement-learning method for formulaic alpha mining. It diagnoses the boosting-style reward of pool-based miners as the joint source of reward non-stationarity and absent generalisation control, and replaces it with an archive-independent reward and a MAP-Elites archive keyed on a novel return-source behaviour descriptor. The resulting selection rule admits a finite-sample ensemble generalisation bound whose complexity scales with the number of occupied archive cells rather than the cardinality of the formula space.

\item \textbf{A unifying perspective} (\cref{ch:introduction}, \cref{ch:conclusion}): the recognition that decompose--specialise--select is a general response to market non-stationarity, and the accompanying lesson that the selector should be decoupled from the policy reward, supported by three increasingly decoupled instantiations.
\end{itemize}

%----------------------------------------------------------------------------------------
\section{Organisation of the Thesis}
\label{sec:intro_outline}

The thesis is organised into four parts. \textbf{Part~I} (Preliminaries) comprises this introduction, \cref{ch:preliminaries}, which establishes the shared financial and reinforcement-learning background of the technical chapters, and \cref{ch:literature_review}, which surveys related work organised by task and synthesises the open challenges that motivate them. \textbf{Part~II} (Reinforcement Learning for Trading) contains the two trading studies: \cref{ch:earnhft} presents EarnHFT for high-frequency trading, and \cref{ch:fineft} presents FineFT for leveraged futures. \textbf{Part~III} (Reinforcement Learning for Signal Generation) contains \cref{ch:alphaqd}, which presents AlphaQD for formulaic alpha mining. \textbf{Part~IV} (Epilogue) consists of~\cref{ch:conclusion}, which revisits the unifying perspective, states the limitations of each study, and outlines directions for future work. Supporting proofs and extended derivations are collected in the appendix.
 % Introduction
%!TEX root=../mythesis.tex
% Chapter 2 -- Preliminaries

\chapter{Preliminaries} % Main chapter title
\chaptermark{Preliminaries}
\label{ch:preliminaries}

This chapter introduces the shared vocabulary used throughout the thesis. \sref{sec:bg_finance} introduces the financial instruments, market microstructure, and evaluation metrics common to all three studies; the metric definitions given there are referenced rather than repeated in later chapters. \sref{sec:bg_rl} is a brief primer on the reinforcement-learning machinery---from Markov decision processes to GFlowNets and quality-diversity search---used across the thesis. The survey of related work and the open challenges that motivate the technical chapters are deferred to \cref{ch:literature_review}.

%----------------------------------------------------------------------------------------
\section{Financial Markets and Quantitative Trading}
\label{sec:bg_finance}

\subsection{Instruments and Market Microstructure}
\label{sec:bg_instruments}

The three studies operate on three increasingly complex trading settings. EarnHFT (\cref{ch:earnhft}) trades a single \emph{spot} cryptocurrency, where the asset is bought and sold outright and only long positions are taken. FineFT (\cref{ch:fineft}) trades \emph{perpetual futures}, standardised contracts that permit both long and short positions, magnify exposure through leverage, and exchange a periodic funding payment between the two sides of the contract; depending on the sign of the funding rate and the direction of the position, a trader may either pay or receive funding. AlphaQD (\cref{ch:alphaqd}) operates on the \emph{cross-section} of an entire equity market, where the object of interest is the relative ranking of thousands of assets on each day rather than the timing of a single instrument. The following concepts recur throughout.

\paragraph{Limit order book (LOB).} The limit order book aggregates all resting buy and sell orders by price level and is the standard description of market microstructure~\cite{madhavan2000market}. A $d$-level LOB at time $t$ is denoted $B_t=(p_t^{b_1},p_t^{a_1},q_t^{b_1},q_t^{a_1},\dots,p_t^{b_d},p_t^{a_d},q_t^{b_d},q_t^{a_d})$, where $p_t^{b_i}$ and $p_t^{a_i}$ are the level-$i$ bid and ask prices and $q_t^{b_i},q_t^{a_i}$ are the corresponding quantities (the symbol $m$ is reserved for the annualisation constant of \sref{sec:bg_metrics}). \fref{fig:bg_lob} illustrates a snapshot. A \emph{market order} executes immediately against the resting orders on the opposite side, filling at the best quote first; only when its size exceeds the depth available at the best level does it \emph{walk the book}, consuming successively deeper levels, so that its average execution price is worse than the best quote---an effect known as slippage that any high-fidelity simulator must model.

\begin{figure}[!htb]
  \centering
  \begin{tabular}{lrr}
    \toprule
    Level & Price & Quantity \\
    \midrule
    Ask 5 & 105 & 90 \\
    Ask 4 & 104 & 62 \\
    Ask 3 & 103 & 45 \\
    Ask 2 & 102 & 28 \\
    Ask 1 (best ask) & 101 & 10 \\
    \midrule
    \multicolumn{3}{c}{mid-price $=100.5$} \\
    \midrule
    Bid 1 (best bid) & 100 & 12 \\
    Bid 2 & 99 & 30 \\
    Bid 3 & 98 & 48 \\
    Bid 4 & 97 & 66 \\
    Bid 5 & 96 & 85 \\
    \bottomrule
  \end{tabular}
  \caption[A limit order book snapshot.]{A schematic five-level limit order book snapshot with illustrative integer prices and quantities. Resting asks (sell orders) sit above the mid-price and resting bids (buy orders) below it, each aggregated by price level. A market order first consumes the depth at the best quote; if its size exceeds that depth, it walks the book \emph{away} from the mid-price into successively deeper levels, so larger orders execute at progressively worse average prices.}
  \label{fig:bg_lob}
\end{figure}

\paragraph{OHLC and trades.} Executed transactions are summarised over a time interval by the open, high, low, and close prices, denoted $x_t=(p_t^o,p_t^h,p_t^l,p_t^c)$. The raw stream of executed orders, or trades, records price, quantity, and side. For leveraged contracts, a \emph{mark price} $M_t$ is used as the reference fair value at which the position is valued and against which liquidation is assessed.

\paragraph{Technical indicators.} Both microstructure and price-history features are summarised by technical indicators---deterministic functions of recent LOB and trade data designed to expose latent patterns, denoted $y_t=\phi(\cdot)$. Examples are the imbalance volume (IV)~\cite{chordia2002order}, which measures buying versus selling pressure in the book, and the moving-average convergence divergence (MACD)~\cite{appel2005technical}, which contrasts trends over different windows. These indicators are inputs to the RL state in the two trading studies (\cref{ch:earnhft} and \cref{ch:fineft}); the alpha-mining study of \cref{ch:alphaqd} instead composes its own signals directly from raw price--volume fields. As standalone rules, however, technical indicators are sensitive to hyperparameters and produce false signals in volatile markets.

\paragraph{Position, leverage, and value.} A trader's position $H_t$ is the signed quantity held; under leverage $L_t$, the same capital supports a larger position and amplifies both gains and losses. For spot trading, the account value (net value) is the cash balance plus the position marked at the current price, net of realised costs such as commissions. For leveraged futures, the corresponding quantity is the \emph{margin balance}: the wallet balance plus the unrealised profit-and-loss of open positions, net of commissions and funding payments---the full notional value of the position does not enter the account value, which is precisely what leverage means.

\subsection{The Trading-Frequency Spectrum}
\label{sec:bg_frequency}

Trading strategies are conventionally classified by the frequency at which they act, and the binding difficulty changes along this spectrum. \emph{High-frequency} strategies act at the second or sub-second scale and profit from fleeting microstructure imbalances; their defining computational challenge is the sheer length of the decision horizon. \emph{Medium-frequency} strategies act at the minute-to-hour scale. \emph{Low-frequency} strategies act daily or slower; here the difficulty is not the horizon but the discovery of robust predictive signals from a low signal-to-noise cross-section. Acting frequency should not be conflated with other, logically independent dimensions of a trading problem: leverage, funding payments, and overnight exposure are properties of the instrument traded and of the holding period, not of the sampling rate itself, although longer holding periods tend to make them more consequential. In this thesis these dimensions become material in the leveraged perpetual-futures setting of \cref{ch:fineft}. Note also that the 24/7 operation of cryptocurrency markets removes the closed-market gap between trading sessions, but not the risk of holding a position overnight. This thesis traverses the entire frequency spectrum, from second-level execution in \cref{ch:earnhft} to daily signal generation in \cref{ch:alphaqd}.

\subsection{Evaluation Metrics}
\label{sec:bg_metrics}

Strategies are evaluated at two layers. The \emph{portfolio layer} measures the realised profit-and-loss of an executed strategy; the \emph{signal layer}, used for alpha mining, measures the predictive quality of a signal independently of any portfolio construction. Let $V_1,\dots,V_{n+1}$ denote the account value at the $n{+}1$ decision points of the trading period, and let $\mathbf{ret}=(ret_1,\dots,ret_n)$ with $ret_i=(V_{i+1}-V_i)/V_i$ denote the per-step \emph{simple} (not log) returns. Throughout, $m$ is the number of decision steps in a year at the strategy's acting frequency (the annualisation constant), the risk-free rate is taken to be zero, and $\mathbb{E}[\cdot]$ and $\sigma[\cdot]$ denote the sample mean and sample standard deviation over the return series.

\paragraph{Portfolio-layer metrics.} The following six metrics are used throughout \cref{ch:earnhft} and \cref{ch:fineft}, and the return/risk pair recurs in \cref{ch:alphaqd}.
\begin{itemize}
\item \textbf{Total Return (TR)}: the overall return over the trading period, $TR=(V_{n+1}-V_1)/V_1$, where $V_1$ and $V_{n+1}$ are the initial and final account values.
\item \textbf{Annualised Volatility (AVOL)}: $\mathrm{AVOL}=\sigma[\mathbf{ret}]\,\sqrt{m}$, a measure of risk.
\item \textbf{Maximum Drawdown (MDD)}~\cite{magdon2004maximum}: the largest peak-to-trough decline of the value curve, $\mathrm{MDD}=\max_{1\le i\le j\le n+1}\,(V_i-V_j)/V_i$, capturing the worst case. MDD is reported as a positive fraction throughout this thesis, so smaller is better (marked $\downarrow$ in the result tables).
\item \textbf{Annualised Sharpe Ratio (ASR)}~\cite{sharpe1994sharpe}: $\mathrm{ASR}=\sqrt{m}\,\frac{\mathbb{E}[\mathbf{ret}]}{\sigma[\mathbf{ret}]}$, return per unit of total risk (with zero risk-free rate); the $\sqrt{m}$ factor multiplies the whole ratio, so ASR equals the annualised mean return $m\,\mathbb{E}[\mathbf{ret}]$ divided by AVOL.
\item \textbf{Annualised Calmar Ratio (ACR)}: $\mathrm{ACR}=m\,\mathbb{E}[\mathbf{ret}]/\mathrm{MDD}$, return per unit of worst-case loss.
\item \textbf{Annualised Sortino Ratio (ASoR)}: $\mathrm{ASoR}=\sqrt{m}\,\frac{\mathbb{E}[\mathbf{ret}]}{DD}$, where the downside deviation $DD$ is the standard deviation of the negative per-step returns (downside threshold zero) and the $\sqrt{m}$ factor is placed as in the ASR; it penalises only harmful volatility.
\end{itemize}

\paragraph{Signal-layer metrics.} For a cross-sectional signal $\alpha$ and an $h$-day forward return $r^{(h)}$, the daily \emph{information coefficient} is the cross-sectional correlation $\mathrm{IC}_t(\alpha)=\corr(\alpha_t,r^{(h)}_t)$. Its time average and stability are combined into the \emph{information ratio} $\mathrm{ICIR}(\alpha)=\overline{\mathrm{IC}(\alpha)}/\hat\sigma(\mathrm{IC}(\alpha))$, the canonical scalar quality of a single alpha in the active-management literature~\cite{grinold2000active} and standard engineering practice in formulaic-alpha libraries~\cite{kakushadze2016101}, since it rewards predictive strength and temporal stability simultaneously. The rank-based variant, \emph{RankIC}, replaces Pearson with Spearman correlation and is robust to outliers.

%----------------------------------------------------------------------------------------
\section{Reinforcement Learning Foundations}
\label{sec:bg_rl}

This section is a brief primer, confined to the reinforcement-learning machinery the technical chapters actually use; a reader fluent in RL may skip directly to the task survey of \cref{ch:literature_review}, which carries the weight of the literature review.

\subsection{Markov Decision Processes}
\label{sec:bg_mdp}

A Markov decision process (MDP) is the tuple $(\mathcal{S},\mathcal{A},P,r,\gamma,T)$, where $\mathcal{S}$ is the state space, $\mathcal{A}$ the action space, $P:\mathcal{S}\times\mathcal{A}\times\mathcal{S}\to[0,1]$ the transition function, $r$ the reward function, $\gamma\in(0,1]$ the discount factor, and $T$ the horizon. At each step the agent observes $s_t$, takes an action $a_t$ under a policy $\pi_\theta:\mathcal{S}\times\mathcal{A}\to[0,1]$, and receives a reward $r_t$ and a next state $s_{t+1}$. The objective is to find parameters maximising the expected discounted return, $\theta^\star=\argmax_\theta \mathbb{E}_{s_0\sim\rho_0,\,\pi_\theta}\big[\sum_{t=0}^{T}\gamma^t r_t\big]$, where $\rho_0$ is the initial-state distribution; the optimal policy is then $\pi^\star=\pi_{\theta^\star}$. The convergence guarantees of standard RL algorithms presuppose that $P$ and $r$ are \emph{stationary}; the central difficulty of this thesis is that financial environments are not, a fact made explicit in \cref{ch:fineft} by the dynamic-parametric MDP formulation that augments the tuple with a slowly evolving latent dynamic.

\subsection{Value-Based Methods}
\label{sec:bg_value}

Value-based methods learn action-value functions. For a policy $\pi$, $Q^\pi(s,a)$ is the expected return of taking $a$ in $s$ and following $\pi$ thereafter; the \emph{optimal} action-value function $Q^\star(s,a)=\max_\pi Q^\pi(s,a)$ is the expected return of taking $a$ in $s$ and acting optimally thereafter. Q-learning~\cite{watkins1992qlearning} learns an estimate $Q$ of $Q^\star$ by updating it towards the temporal-difference (TD) target $r+\gamma\max_{a'}Q(s',a')$. The Deep Q-Network (DQN)~\cite{mnih2015human} parameterises $Q$ with a neural network and stabilises learning through experience replay and a target network. Double DQN (DDQN)~\cite{van2016deep} decouples action selection from evaluation to reduce the overestimation bias of the $\max$ operator. EarnHFT and FineFT both build on DDQN, augmenting its TD target with supervision from a dynamic-programming teacher, a design related in spirit to deep Q-learning from demonstrations~\cite{hester2018dqfd}; the convergence of this augmented operator is established in \cref{ch:earnhft}.

\subsection{Policy-Based Methods}
\label{sec:bg_policy}

Policy-based methods optimise the policy directly. Proximal Policy Optimisation (PPO)~\cite{schulman2017proximal} maximises a clipped surrogate objective with importance sampling, trading some bias for stability and sample reuse. In trading, policy-based methods are prone to collapsing onto trivial policies---holding the current position---when transaction costs make any action locally unprofitable, a failure mode documented empirically in \cref{ch:earnhft}.

\subsection{Distributional and Ensemble Reinforcement Learning}
\label{sec:bg_dist_ensemble}

Two families address the high stochasticity of financial rewards. \emph{Distributional} RL learns the full return distribution rather than its mean; quantile-regression DQN~\cite{dabney2018distributional} and its risk-averse variants represent the return by a set of quantiles and can act on the lower tail to avoid worst cases. \emph{Ensemble} RL trains several learners and aggregates them: averaged DQN reduces target variance, while SUNRISE~\cite{lee2021sunrise} treats inter-learner disagreement as a measure of uncertainty and down-weights uncertain samples. FineFT belongs to the ensemble family but departs from prior work by updating learners \emph{selectively} rather than uniformly, which is what allows specialists to emerge.

\subsection{Hierarchical Reinforcement Learning}
\label{sec:bg_hrl}

Hierarchical RL decomposes a long-horizon task into a hierarchy of sub-tasks, typically a high-level policy that selects among low-level policies or temporally extended options~\cite{sutton1999between,pateria2021hierarchical}. In trading, this naturally maps onto a separation of time scales: a slow policy that reacts to macro conditions and fast policies that execute. Prior hierarchical trading frameworks address portfolio management and order execution~\cite{wang2021commission} or select among strategies by market situation~\cite{niu2022metatrader}. EarnHFT brings this idea to single-asset high-frequency trading, where the high-level router operates at the minute scale and the low-level specialists at the second scale.

\subsection{Generative Flow Networks}
\label{sec:bg_gflownet}

Generative flow networks (GFlowNets)~\cite{bengio2021gflownet} learn a stochastic policy that constructs compositional objects---here, symbolic expressions---step by step, with the aim that the probability of producing a complete object be proportional to a terminal reward. This proportionality holds exactly only at the global optimum of the training objective; at the optimum of the trajectory-balance objective~\cite{malkin2022trajectorybalance} used here, the sampler is reward-proportional, and in practice a trained GFlowNet samples approximately proportionally to reward rather than collapsing onto the single maximiser, so it covers many high-reward modes of the search space under a fixed budget. This mode-covering property is precisely what an alpha-mining search requires, and AlphaQD adopts a relational-graph GFlowNet as its generator.

\subsection{Quality-Diversity Search and MAP-Elites}
\label{sec:bg_qd}

Quality-diversity (QD) algorithms~\cite{pugh2016quality} seek not a single optimum but an archive of high-quality solutions that are diverse along a chosen \emph{behaviour descriptor}. MAP-Elites~\cite{mouret2015mapelites} discretises the descriptor space into cells and retains the best solution(s) per cell, yielding a structured population that covers the behaviour space. QD has not previously been applied to formulaic alpha mining; AlphaQD's contribution is to identify a meaningful behaviour descriptor for this domain---the normalised daily IC series, an empirical proxy for a formula's return source---and to use the archive as the mechanism that organises diversity and selection outside the learning reward.

\subsection{Variational Autoencoders and Out-of-Distribution Detection}
\label{sec:bg_vae}

A variational autoencoder (VAE)~\cite{kingma2013auto} learns a probabilistic encoder that maps inputs to a regularised latent distribution and a decoder that reconstructs them. Inputs resembling the training distribution tend to reconstruct well, so the VAE reconstruction loss is widely used as a heuristic reconstruction-based score for out-of-distribution (OOD) detection; it is a practical signal rather than a calibrated probabilistic measure. FineFT trains a VAE per market dynamic and uses its loss to detect when the current market state lies outside a specialist's competence, converting OOD detection into an explicit risk control.
 % Background and Literature Review
\part{Reinforcement Learning for Trade Execution}
%!TEX root=../mythesis.tex
% Chapter 3 -- Literature Review

\chapter{Literature Review} % Main chapter title
\chaptermark{Literature Review}
\label{ch:literature_review}

This chapter surveys the prior work on which the three technical chapters build. With the financial and reinforcement-learning preliminaries of \cref{ch:preliminaries} in hand, \sref{sec:bg_rl4qt} reviews reinforcement learning for quantitative trading organised by task---single-asset and intraday trading, high-frequency trading, portfolio management, order execution, strategy routing, and formulaic alpha mining---before turning to the cross-cutting question of uncertainty and non-stationarity and closing with further related work. \sref{sec:bg_gaps} then synthesises the open challenges that motivate \cref{ch:earnhft}, \cref{ch:fineft}, and \cref{ch:alphaqd}.

%----------------------------------------------------------------------------------------
\section[Reinforcement Learning for Quantitative Trading]{Reinforcement Learning for\\Quantitative Trading}
\label{sec:bg_rl4qt}

Reinforcement learning has been applied across the entire quantitative-trading stack, and we survey that body of work organised by task. For each task we note the dominant design and the assumption it leaves unaddressed, which the open challenges of \sref{sec:bg_gaps} then collect. Open frameworks and benchmarks have standardised much of this research---FinRL~\cite{liu2021finrl}, TradeMaster~\cite{sun2024trademaster}, Qlib~\cite{yang2020qlib}, and the PRUDEX evaluation compass~\cite{sun2023prudex}---and the studies in this thesis follow their evaluation conventions where applicable.

\subsection{Single-Asset and Intraday Trading}
\label{sec:bg_intraday}

The earliest and largest body of RL-for-trading work targets a single asset at the daily-to-intraday scale, where the agent maps a window of price and technical features to a discrete trading action. Foundational studies established deep RL as a viable trader on this task~\cite{zhang2019deep,theate2021application,li2019deep}, and later work improved robustness and sample efficiency through imitation of profitable behaviour~\cite{liu2020adaptive,ding2018investor}, multimodal or convolutional state encoders~\cite{shin2019deep}, reward shaping for sparse-reward environments~\cite{takara4411793deep}, and multi-objective formulations for intraday trading of stock-index futures~\cite{si2017multi}. DeepScalper~\cite{sun2022deepscalper} adds a hindsight bonus and an auxiliary price-prediction task to improve generalisation in intraday trading and remains a strong representative of the single-agent design. CDQNRP~\cite{zhu2022quantitative}, sometimes reused as a baseline in HFT studies, evaluates a perturbed convolutional DQN at daily frequency and belongs to this adjacent single-asset group rather than to HFT proper. The common thread is a single, monolithic policy: effective at its target frequency, but not by itself designed to address the horizon growth of second-level trading or the regime non-stationarity that erodes one policy over long deployments.

\subsection{High-Frequency Trading}
\label{sec:bg_hft}

For high-frequency trading specifically, DRA~\cite{briola2021deep} couples an LSTM~\cite{hochreiter1997long} state encoder with PPO, and asynchronous deep double-duelling Q-learning has been applied to rapid trading-signal decisions~\cite{nagy2023asynchronous}. A common non-learned alternative is rule-based technical analysis, built on signals such as MACD~\cite{appel2005technical} and the five-level order-book volume-imbalance rule described in \sref{sec:bg_instruments}. EarnHFT addresses second-level position control through hierarchy; FineFT is a separate five-minute leveraged-futures study rather than an HFT method.

\subsection{Portfolio Management and Asset Allocation}
\label{sec:bg_portfolio}

Beyond single-asset timing, a substantial line applies RL to portfolio management, where the action is a continuous allocation across many assets. Early frameworks cast allocation as a deterministic policy-gradient problem over price tensors~\cite{jiang2017deep,jiang2017cryptocurrency}; subsequent work added explicit risk--return balancing~\cite{wang2021deeptrader}, ensemble strategies that switch among actor--critic learners~\cite{yang2020deep}, hierarchical decomposition of selection, allocation, and commission-aware trade execution~\cite{gao2021framework,zha2022hierarchical,wang2021commission}, and maskable representations that respect tradability constraints~\cite{zhang2024reinforcement}. Margin Trader~\cite{gu2023margin} extends the setting to leverage and short selling---the regime that \cref{ch:fineft} studies. Portfolio management is not the subject of this thesis, but it shares the central difficulty of non-stationarity and motivates much of the ensemble and hierarchical machinery reused here.

\subsection{Order Execution}
\label{sec:bg_execution}

A complementary task is optimal order execution---acquiring or liquidating a large position while minimising market impact. Recent RL work on execution stresses generalisation across instruments and regimes~\cite{zhang2023towards}, and the adjacent microstructure task of automated market making has been tackled with imitative RL over predictive state representations~\cite{niu2023imm}. Execution is the task in which EarnHFT's historical-replay assumption---that simulated actions do not alter the recorded price path (\sref{sec:earnhft_stage1})---is most visibly an approximation, and the literature here is correspondingly attentive to market impact.

\subsection{Hierarchical and Routing Approaches}
\label{sec:bg_routing}

A growing line of work does not commit to one policy but routes among several. MetaTrader~\cite{niu2022metatrader} learns a router that picks the most suitable strategy for the current market, and mixture-of-experts designs train diversified traders and combine them~\cite{sun2023mastering}. Winnow~\cite{littlestone1988winnow} is originally a mistake-driven linear-threshold learning algorithm; the trading experiments later use a performance-based adaptation under that name, so it is not treated here as a generic theorem for expert weighting. EarnHFT~\cite{qin2024earnhft} and MacroHFT~\cite{zong2024macrohft} construct hierarchical MDPs for high-frequency trading and train the high-level router \emph{by reinforcement learning on historical data}; in EarnHFT the hindsight Q supervision applies only to the second-level specialists, and the router itself is trained without a Q-teacher (\cref{ch:earnhft}). In both systems the router is fitted offline and then frozen for deployment rather than updated online.

\subsection{Formulaic Alpha Mining}
\label{sec:bg_alpha}

Formulaic alpha mining seeks short symbolic expressions whose cross-sectional values predict returns. Hand-crafted catalogues such as Alpha101~\cite{kakushadze2016101} gave way to automated search. Genetic programming~\cite{stephens2015gplearn} and its hierarchical and warm-started variants~\cite{zhang2020autoalpha,ren2024warmstart} evolve expression trees by mutation and crossover. The broader quality-diversity lineage seeks a repertoire rather than one optimum~\cite{pugh2016quality}; MAP-Elites discretises a behaviour space and retains high-quality occupants of its cells~\cite{mouret2015mapelites} (see also \sref{sec:bg_qd}). AutoAlpha is especially relevant because its PCA-QD component is an earlier use of these ideas in alpha mining; AlphaQD's novelty is therefore the separate return-source MAP-Elites archive and its separation from the generator reward, not QD in the domain per se. AlphaEvolve~\cite{cui2021alphaevolve} is another learning framework for discovering novel alphas and is reviewed here even though it is absent from the unified experimental table. Deep symbolic optimisation (DSO)~\cite{petersen2021dso} trains a recurrent generator against a standalone expression score, providing a counterexample to the claim that every RL-style symbolic generator must use a pool-coupled reward. The pool-based reinforcement-learning line includes AlphaGen~\cite{yu2023alphagen}, which casts mining as a PPO problem whose terminal reward is the marginal contribution of a new alpha to a running pool. AlphaQCM~\cite{zhu2025alphaqcm} adds a distributional value head and a variance bonus; RiskMiner~\cite{ren2024riskminer} pairs risk-seeking Monte-Carlo tree search with a pool reward; Alpha\textsuperscript{2}~\cite{xu2024alpha2} searches logical formulas; AlphaForge~\cite{shi2025alphaforge} replaces the RL search with a generator--predictor surrogate and a dynamic re-combination stage; and AlphaSAGE~\cite{chen2025alphasage} encodes expressions as typed graphs and trains a GFlowNet against a dense, multi-faceted reward. RiskMiner and Alpha\textsuperscript{2} are contextual methods rather than reported baselines here; their omission limits the coverage of the empirical comparison. A parallel and very recent line replaces the search policy with a large language model~\cite{wang2025alphagpt,tang2025alphaagent,li2024fama}. As \cref{ch:alphaqd} details, AlphaQD addresses the moving target in this influential pool-based family; it does not claim that every symbolic RL objective is pool-coupled.

\subsection{Uncertainty and Non-Stationary RL}
\label{sec:bg_uncertainty}

The leverage setting of \cref{ch:fineft} foregrounds two kinds of uncertainty. \emph{Aleatoric} uncertainty is the irreducible noise of the environment. Leverage does not change the underlying transition uncertainty, but it magnifies that noise in account returns and rewards; one risk-aware quantile construction evaluated in automated driving~\cite{bernhard2019addressing} and adapted to the trading setting in \cref{ch:fineft} targets the resulting lower-tail risk. Ensemble value methods address a different axis: Averaged-DQN reduces target variance~\cite{anschel2017averaged}, while SUNRISE uses ensemble disagreement in its sample weighting and action selection~\cite{lee2021sunrise}. \emph{Epistemic} uncertainty is the model's ignorance about states it has not seen and motivates OOD detection. A VAE~\cite{kingma2013auto} can supply a reconstruction-based distribution-familiarity score, although deep generative scores can be misleading under distribution shift~\cite{zhang2021understanding}; the score is a heuristic rather than calibrated epistemic uncertainty (see \sref{sec:bg_vae}). Modelling slowly drifting environments as a dynamic-parameter MDP~\cite{xie2020deep} provides the formal setting in which a short window can be treated as approximately stationary. FineFT combines selective ensemble updates with one such VAE score per inferred dynamic.

\subsection{Further Related Work}
\label{sec:bg_further}

Beyond the task-specific lines above, four strands of prior work help situate the contributions of this thesis. Moody and Saffell~\cite{moody2001learning} provide a foundational treatment of reinforcement learning for trading, training a recurrent trader by direct reinforcement on risk-adjusted performance rather than through an intermediate value function. Nevmyvaka, Feng, and Kearns~\cite{nevmyvaka2006reinforcement} give the classical large-scale study of reinforcement learning for optimised trade execution on limit-order-book data, the setting in which the exogenous-price approximation exploited by \cref{ch:earnhft} is most commonly invoked. The Q-teacher of \cref{ch:earnhft} is related to deep Q-learning from demonstrations~\cite{hester2018dqfd}, which shapes value learning with a supervised loss on expert data; the Q-teacher differs in that dynamic programming computes sample-path hindsight continuation values under the exogenous-replay assumption rather than using externally supplied demonstrations. These are trajectory-conditioned training signals, not optimal action values of the unknown stochastic market. Finally, Harvey, Liu, and Zhu~\cite{harvey2016cross} document how multiple testing inflates apparent performance in cross-sectional return research---a winner's-curse concern that applies directly to alpha mining, where many candidate formulas are evaluated on the same data (\cref{ch:alphaqd}).

%----------------------------------------------------------------------------------------
\section{Open Challenges}
\label{sec:bg_gaps}

This survey exposes three gaps that the technical chapters close.

\begin{itemize}
\item \textbf{Gap I (high-frequency efficiency and adaptivity).} Prior single-policy methods do not jointly address second-level horizon length and changing trends. A method is needed that improves learning efficiency at second scale and can route among differentiated policies. This motivates EarnHFT (\cref{ch:earnhft}).

\item \textbf{Gap II (risk and distribution awareness under leverage).} Methods designed for low-leverage markets do not necessarily control the magnified P\&L consequences of leverage or flag unfamiliar states, leaving them exposed during abrupt shifts. This motivates FineFT (\cref{ch:fineft}).

\item \textbf{Gap III (reward stationarity and selection in alpha mining).} A dominant family of pool-based RL alpha miners couples the reward to an evolving pool, training the policy against a moving target and complicating archive-conditional selection analysis. This motivates AlphaQD (\cref{ch:alphaqd}).
\end{itemize}

A thread common to all three gaps is the placement of the aggregation mechanism. EarnHFT learns it, FineFT derives it heuristically from generative uncertainty estimates, and AlphaQD removes it from the reward entirely---a progression that the remainder of the thesis develops.
 % EarnHFT: Hierarchical RL for High-Frequency Trading
%!TEX root=../mythesis.tex
% Chapter 3 -- EarnHFT

\chapter{EarnHFT: Hierarchical Reinforcement Learning for High-Frequency Trading} % Main chapter title
\chaptermark{EarnHFT for High-Frequency Trading}
\label{ch:earnhft}

%----------------------------------------------------------------------------------------
\section{Introduction}
\label{sec:earnhft_intro}

High-frequency trading (HFT) uses computer algorithms to place and cancel orders at very short time scales, and accounts for a majority of volume in modern financial markets~\cite{ALMEIDA2023100785}. It is especially prevalent in the cryptocurrency market, whose round-the-clock trading removes overnight risk and whose dramatic price fluctuations create frequent short-lived opportunities. A profitable HFT strategy can outperform a low-frequency one, which is why it is pursued aggressively by quantitative traders.

Reinforcement learning has produced strong results on low-frequency tasks such as daily stock or futures trading~\cite{sun2022deepscalper,briola2021deep}, but these methods rarely survive the move to HFT, for two reasons identified in \sref{sec:bg_gaps} as Gap~I.

\begin{itemize}
\item \textbf{The horizon is extreme, so RL is data-inefficient.} A second-level agent must be evaluated over dozens of days, giving a horizon of roughly one million steps---for instance, $60\times 60\times 24\times 12 = 1{,}036{,}800$ seconds over twelve days---against the few-thousand-step horizon of Atari benchmarks~\cite{mnih2013playing}. Such horizons demand far more data to converge~\cite{zhang2023towards}, straining computation.
\item \textbf{The market drifts, so a single agent cannot maintain performance.} Standard RL assumes the training and testing environments coincide. In cryptocurrency, the market trend changes substantially between them; an agent trained under one trend incurs large losses when the trend reverses.
\end{itemize}

This chapter presents EarnHFT, an \textbf{E}fficient hier\textbf{A}rchical \textbf{R}einforcement lear\textbf{N}ing method for \textbf{H}igh \textbf{F}requency \textbf{T}rading, whose three stages are summarised in \fref{fig:earnhft_pipeline}. The design follows the decompose--specialise--select pattern of \sref{sec:intro_pattern}.

\begin{enumerate}
\item \textbf{Stage~I (Q-teacher).} We exploit a structural property of execution---that at second scale a single trader's orders do not move prices, so future prices are independent of the policy---to compute, by dynamic programming, the optimal action value for every state. Used as a regulariser and a demonstrator, this ``Q-teacher'' accelerates and stabilises second-level training.
\item \textbf{Stage~II (specialist pool).} We train hundreds of agents under different market-trend preferences, segment and label the market without hyperparameter tuning, and select a small pool of the most profitable specialists.
\item \textbf{Stage~III (router).} We train a minute-level router that dynamically picks a second-level specialist from the pool to maintain performance as the market changes.
\end{enumerate}

Across four cryptocurrency markets spanning bull, bear, and sideways conditions, EarnHFT outperforms six state-of-the-art baselines on six financial metrics, exceeding the runner-up by at least 30\% in profitability.

\begin{figure}[!htb]
\centering
\includegraphics[width=\textwidth]{Chapter4/earnhft_pipeline}
\caption[The EarnHFT framework.]{Overview of the three stages of EarnHFT. Stage~I computes a Q-teacher to enhance the performance and training efficiency of second-level agents. Stage~II trains diverse agents under various market trends and selects a small fraction by profitability to form an agent pool. Stage~III trains a minute-level router that dynamically picks a second-level agent from the pool.}
\label{fig:earnhft_pipeline}
\end{figure}

%----------------------------------------------------------------------------------------
\section{Problem Formulation}
\label{sec:earnhft_formulation}

The financial concepts---limit order book, OHLC, technical indicators, market order, position, and net value---are defined in \sref{sec:bg_instruments}. We restrict attention to a single asset with long-only positions, and the objective is to maximise net value $V_t=V_{ct}+H_t\, p_t^{b_1}$, where $V_{ct}$ is cash and $H_t$ the position valued at the best bid. The execution price of a market order walks the book, $E_t(M)=\sum_i (p_t^i\min(q_t^i,R_{i-1}))(1+\sigma)$, where $R_{i-1}$ is the residual quantity after level $i$ and $\sigma$ the commission rate.

\subsection{A Hierarchical Markov Decision Process}
\label{sec:earnhft_hmdp}

The two challenges of \sref{sec:earnhft_intro} have a common structural cause: micro-level information drifts, so no single agent profits under all trends, while macro-level information, an aggregation of the micro level, carries the slow trend signal. We therefore decompose HFT into a hierarchical MDP $(\mathrm{MDP}_h,\mathrm{MDP}_l)$, in which a low-level MDP at second scale models micro-level dynamics and execution, and a high-level MDP at minute scale models macro-level trend and strategy switching:
\begin{align*}
\mathrm{MDP}_h &= (S_h,A_h,P_h,R_h,\gamma_h,T_h), &
\mathrm{MDP}_l &= (S_l,A_l,P_l,R_l,\gamma_l,T_l).
\end{align*}

\paragraph{Low level (second scale).} The state $S_{lt}$ comprises a latent representation $y_{lt}$ of 54 micro-level technical indicators computed from a length-60 rolling window of second-level OHLC and LOB snapshots, together with the current position $H_t$. The action $a_{lt}$ is a target position chosen from a finite pool $\{0,\frac{H}{|A|-1},\dots,H\}$; if it differs from the current position, a market order is executed immediately. The reward $r_{lt}=H_{t+1}p_{t+1}^{b_1}-(H_t p_t^{b_1}+E_t(H_{t+1}-H_t))$ is the per-second change in net value.

\paragraph{High level (minute scale).} The state $S_{ht}$ comprises a latent representation $y_{ht}$ of 19 macro-level indicators from a length-60 rolling window of minute-level OHLC, together with $H_t$. The action $a_{ht}$ selects an agent from the pre-trained pool $A_h$. The reward $r_{hT}=\sum_{t=T}^{T+\tau}r_{lt}$ is the net-value change accrued by the selected low-level agent over one minute. In summary, every minute the high-level router picks a low-level agent, which then adjusts its position every second; the goal is to find a set of specialists and a router that jointly maximise total profit.

%----------------------------------------------------------------------------------------
\section{Method}
\label{sec:earnhft_method}

\subsection{Stage I: Efficient Reinforcement Learning with a Q-Teacher}
\label{sec:earnhft_stage1}

A long trajectory normally inflates the cost of RL. In the low-level MDP, however, prices are not influenced by the policy, so future prices are known in hindsight on historical data. This lets us construct the optimal action value $Q^\star$ by dynamic programming~\cite{sutton2018reinforcement} and use it to teach the agent, both as a value supervisor and as a demonstrator.

\paragraph{Optimal value supervisor.} Because the position set is finite, the optimal action value can be computed backwards: from the last step, the value of each (position, action) pair is the immediate reward plus the best continuation value, as in Algorithm~\ref{alg:earnhft_qtable}. We add to the DDQN loss a supervision term that is the Kullback--Leibler divergence between the agent's action values and the optimal action values at the same state:
\begin{equation}
\label{eq:earnhft_loss}
L(\theta_i) = L_{td} + \alpha\,\mathrm{KL}\!\left(Q_t(\chi,p,\cdot;\theta_i)\,\|\,Q^\star(\chi,p,\cdot)\right),
\end{equation}
where $L_{td}=\left(r+\gamma\max Q_t(\chi',a,\cdot;\theta_i')-Q_t(\chi,p,a;\theta_i)\right)^2$ is the DDQN temporal-difference error, $\chi$ the latent representation, $p$ the position, and $\alpha$ a coefficient decaying over time. The supervision term lets the agent learn the relative advantage of all actions at a state without exploring them, accelerating training.

\begin{algorithm}[!htb]
\caption{Construction of the Optimal Action Value}
\label{alg:earnhft_qtable}
\begin{algorithmic}[1]
\State \textbf{Input:} time series $\mathcal{D}$ of length $N$, commission rate $\delta$, action space $A$
\State \textbf{Output:} table $Q^\star$ of optimal action values indexed by time, position, action
\State Initialise $Q^\star$ with shape $(N,|A|,|A|)$ and all elements $0$.
\For{$t \gets N-1$ \textbf{to} $1$}
  \For{$p \gets 1$ \textbf{to} $|A|$}
    \For{$a \gets 1$ \textbf{to} $|A|$}
      \State $Q^\star[t,p,a] \gets \max_{a'} Q^\star[t+1,a,a'] + a\,p_{t+1}^{b_1} - (p\,p_{t}^{b_1} + E_t(p-a))$
    \EndFor
  \EndFor
\EndFor
\State \textbf{return} $Q^\star$
\end{algorithmic}
\end{algorithm}

\paragraph{Optimal actor.} Supervision alone is insufficient, because the optimal action value is conditioned on the current position: once the agent deviates from the optimal policy, the supervision target shifts, compounding the deviation. We therefore additionally train on transitions generated by the optimal policy itself---the action with the highest optimal value---which supplies demonstration experience and prevents the agent from settling in a local trap. Algorithm~\ref{alg:earnhft_efficient} collects experience from both the $\epsilon$-greedy policy and the optimal actor, then updates the network using both TD errors and the KL supervision.

\begin{algorithm}[!htb]
\caption{Efficient Reinforcement Learning with a Q-Teacher}
\label{alg:earnhft_efficient}
\begin{algorithmic}[1]
\State \textbf{Input:} time series $\mathcal{D}$ of length $N$, commission rate $\delta$, action space $A$
\State \textbf{Output:} network parameters $\theta$
\State Initialise replay buffer $R$, network $Q_\theta$, target $Q_{\theta'}$; build $Q^\star$ via Algorithm~\ref{alg:earnhft_qtable}; initialise environment.
\For{$t = 1$ \textbf{to} $N-1$} \Comment{exploration rollout}
  \State Choose $a_\epsilon$ by the $\epsilon$-greedy policy; store $(s,a_\epsilon,r,s',Q^\star)$ in $R$.
\EndFor
\For{$t = 1$ \textbf{to} $N-1$} \Comment{optimal-actor rollout}
  \State Choose $a_o = \argmax_{a} Q^\star[t,p,a]$; store $(s,a_o,r,s',Q^\star)$ in $R$.
\EndFor
\State Sample transitions; compute $L$ by \eqref{eq:earnhft_loss}; descend on $\theta$ and set $\theta'=\tau\theta+(1-\tau)\theta'$.
\State \textbf{return} $Q_\theta$
\end{algorithmic}
\end{algorithm}

The Q-teacher does not compromise correctness. Treating the supervised update as an operator $H$ on action values, one can show $H$ is a contraction in the sup-norm whose fixed point is $Q^\star$, so the learned value still converges to the optimum. Following the convergence argument for Q-learning~\cite{jaakkola1993convergence}, it suffices to verify these two properties of $H$. The operator combines the usual Bellman backup with the supervision term and is
\begin{equation}
\label{eq:earnhft_operator}
Hq(x,a) = \lambda\sum_{y\in X}T(y\mid x,a)\bigl(r(x,a,y)+\gamma\max_b q(y,b)\bigr) + (1-\lambda)\Bigl(\tfrac{1}{|A|}\textstyle\sum_{a\in A}q(x,a)+Re(x,a)\Bigr),
\end{equation}
where the supervision residual is $Re(x,a) = q^\star(x,a) - \tfrac{1}{|A|}\sum_{a\in A}q^\star(x,a)$ and $\lambda\in(0,1]$ weights the Bellman backup against the supervision term.

\begin{proposition}[$Q^\star$ is a fixed point of $H$]
\label{prop:earnhft_fixed}
The optimal action value $q^\star$ satisfies $Hq^\star = q^\star$.
\end{proposition}
\begin{proof}
Substituting $q^\star$ into \eqref{eq:earnhft_operator} and using the Bellman optimality equation $q^\star(x,a)=\sum_y T(y\mid x,a)(r(x,a,y)+\gamma\max_b q^\star(y,b))$ for the first term and the definition of $Re$ for the second,
\begin{align*}
Hq^\star(x,a) &= \lambda\, q^\star(x,a) + (1-\lambda)\Bigl(\tfrac{1}{|A|}\textstyle\sum_a q^\star(x,a) + q^\star(x,a) - \tfrac{1}{|A|}\textstyle\sum_a q^\star(x,a)\Bigr) \\
&= \lambda\, q^\star(x,a) + (1-\lambda)\, q^\star(x,a) = q^\star(x,a). \qedhere
\end{align*}
\end{proof}

\begin{proposition}[Contraction in the sup-norm]
\label{prop:earnhft_contraction}
For any action values $q_1,q_2$, $\|Hq_1-Hq_2\|_\infty \le k\,\|q_1-q_2\|_\infty$ with $k=1-(1-\lambda)\gamma<1$.
\end{proposition}
\begin{proof}
The supervision residual $Re(x,a)$ depends only on $q^\star$ and not on the argument of $H$, so it cancels in the difference $Hq_1-Hq_2$. Bounding the $\max$ operator by $|\max_b q_1(y,b)-\max_b q_2(y,b)|\le\max_b|q_1(y,b)-q_2(y,b)|$,
\begin{align*}
\|Hq_1-Hq_2\|_\infty
&\le \lambda\gamma\,\|q_1-q_2\|_\infty + (1-\lambda)\,\|q_1-q_2\|_\infty \\
&= \bigl(1-(1-\lambda)\gamma\bigr)\|q_1-q_2\|_\infty.
\end{align*}
Since $\gamma\in(0,1]$ and $\lambda\in(0,1]$, the factor $k=1-(1-\lambda)\gamma<1$.
\end{proof}

\begin{theorem}[Convergence of the supervised update]
\label{thm:earnhft_convergence}
Under a finite MDP, the action value learned with the optimal-value-supervisor update of \eqref{eq:earnhft_loss} converges to the optimal action value $Q^\star$.
\end{theorem}
\begin{proof}
Propositions~\ref{prop:earnhft_fixed} and~\ref{prop:earnhft_contraction} show that $H$ is a sup-norm contraction with unique fixed point $Q^\star$. By the stochastic-approximation argument of~\cite{jaakkola1993convergence}, the iterates therefore converge almost surely to $Q^\star$.
\end{proof}

\subsection{Stage II: Construction of the Agent Pool}
\label{sec:earnhft_stage2}

Because micro-level information changes rapidly, no single agent maintains performance over a long period; preliminary experiments confirm that training across conflicting trends is incompatible. We therefore decompose the market into trends and develop a specialist for each.

\paragraph{Generating diverse agents.} Prior approaches obtain diversity as a by-product of random seeds or hyperparameter search~\cite{sun2023mastering}, which is unstructured. We instead train agents with explicit \emph{preferences} over market trends. The training series (over three million steps) is cut into chunks of length $L$, each a continuous trend; a preference parameter $\beta$ assigns each chunk with buy-and-hold return rate $r$ a sampling priority
\begin{equation}
\label{eq:earnhft_pes}
f(x) =
\begin{cases}
\dfrac{e^{\beta r}}{\mathrm{pdf}(r)} & \text{if } Q_{\frac{\theta}{2}}(R) \le r \le Q_{1-\frac{\theta}{2}}(R),\\[2mm]
e^{\beta r} & \text{otherwise},
\end{cases}
\end{equation}
where $Q_{\frac{\theta}{2}}(R)$ is the $\frac{\theta}{2}$-quantile of return rates and $\mathrm{pdf}$ is a kernel density estimate,
\begin{equation}
\label{eq:earnhft_kernel}
\mathrm{pdf}(x)=\frac{1}{nh}\sum_{r\in R} K\!\left(\frac{x-r}{h}\right),
\end{equation}
with bandwidth $h$ searched around Silverman's rule~\cite{silverman1984spline} and $K$ a standard normal kernel. Dividing by the density removes the influence of the dataset's own distribution, while $e^{\beta r}$ tilts sampling towards a preferred trend; clipping the extremes prevents a few violent chunks from poisoning the agent everywhere. Different $\beta$ values yield agents specialised to different trends.

\paragraph{Market segmentation and labelling.} To evaluate specialists per trend we need to label the market. Algorithm~\ref{alg:earnhft_segment} denoises the series, splits it at extrema, merges adjacent segments whose dynamic-time-warping~\cite{muda2010voice} and slope differences are small, and labels segments by slope quantiles. Unlike prior methods, it requires no per-dataset hyperparameter tuning.

\begin{algorithm}[!htb]
\caption{Market Segmentation and Labelling}
\label{alg:earnhft_segment}
\begin{algorithmic}[1]
\State \textbf{Input:} time series $\mathcal{D}$ of length $N$; risk threshold $\theta$; number of labels $M$
\State \textbf{Output:} a trend label for every point of $\mathcal{D}$
\State $\mathcal{D}' \gets$ denoise high-frequency noise of $\mathcal{D}$.
\State Divide $\mathcal{D}'$ at its extrema into segments $S$.
\State Merge adjacent segments in $S$ whose DTW and slope differences are small, until $S$ is stable.
\State Compute thresholds $H = Q_{1-\frac{\theta}{2}}(R)$, $L = Q_{\frac{\theta}{2}}(R)$ on segment slopes $R$.
\State Compute per-label slope bounds from the quantiles and thresholds; label each segment.
\State \textbf{return} the label of each segment.
\end{algorithmic}
\end{algorithm}

\paragraph{Agent selection.} Including all trained agents would inflate the router's action space, so we select a small pool. Each point of the validation set is labelled by Algorithm~\ref{alg:earnhft_segment}; agents are evaluated across trends and initial positions, and the most profitable agent under each (trend, initial-position) pair is retained, forming a two-dimensional pool indexed by $(m,n)$ for $m$ trends and $n$ initial positions.

\subsection{Stage III: Dynamic Routing}
\label{sec:earnhft_stage3}

We train the router for the high-level MDP with DDQN~\cite{van2016deep}. Operating at minute scale shortens the trajectory by 98.33\% relative to the second scale, and we further restrict the router, at each step, to agents whose initial position matches the current position, reducing the choice to $m$ agents. We deliberately do \emph{not} compute a Q-teacher for the router: the horizon is already short, and the high-level action is a strategy rather than a target position, which would make the optimal value expensive to compute. With the computation for self-exploration reduced and that for the teacher increased, plain DDQN is the more efficient choice at the high level.

%----------------------------------------------------------------------------------------
\section{Experimental Setup}
\label{sec:earnhft_setup}

\paragraph{Datasets.} We evaluate on four cryptocurrencies covering mainstream and niche assets and bull, bear, and sideways conditions, summarised in \tref{tab:earnhft_dataset}. On each, the last nine days are used for testing, the preceding nine for validation, and the remainder for training. Low-level agents are trained on the training set; the validation set is segmented and labelled for model selection; the router is then trained on the training set and selected on the full validation set before testing.

\begin{table}[!htb]
\setlength\tabcolsep{6pt}
\centering
\caption[EarnHFT dataset statistics.]{The four cryptocurrency datasets, with market dynamics, number of seconds, and chronological span. Dates are in YY/MM/DD format.}
\label{tab:earnhft_dataset}
\begin{tabular}{lcccc}
\toprule
\textbf{Dataset} & \textbf{Dynamics} & \textbf{Seconds} & \textbf{From} & \textbf{To} \\
\midrule
BTC/TUSD & Sideways & 4{,}057{,}140 & 23/03/30 & 23/05/15 \\
BTC/USDT & Sideways & 3{,}884{,}400 & 22/09/01 & 22/10/15 \\
ETH/USDT & Bear     & 3{,}970{,}800 & 22/05/01 & 22/06/15 \\
GALA/USDT & Bull    & 3{,}970{,}740 & 22/07/01 & 22/08/15 \\
\bottomrule
\end{tabular}
\end{table}

\paragraph{Metrics.} We report the six financial metrics defined in \sref{sec:bg_metrics}: TR (profit); ASR, ACR, ASoR (risk-adjusted profit); and AVOL, MDD (risk).

\paragraph{Baselines.} We compare against six methods: two policy-based RL algorithms, PPO~\cite{schulman2017proximal} and DRA~\cite{briola2021deep}; two value-based RL algorithms, DQN~\cite{mnih2015human} and CDQNRP~\cite{zhu2022quantitative}; and two rule-based methods, MACD~\cite{krug2022enforcing} and IV~\cite{chordia2002order}. Where official implementations exist we reuse their hyperparameters; otherwise we reimplement faithfully from the original papers.

\paragraph{Training.} All experiments run on a single RTX~4090 GPU. The commission rate is $0$ for BTC/TUSD and $0.02\%$ for the others, following Binance. We train with $\beta\in\{-90,-10,30,100\}$ for 50 epochs each, producing 200 agents, optimised with Adam. The full set of experiments completes in about ten hours.

%----------------------------------------------------------------------------------------
\section{Results and Analysis}
\label{sec:earnhft_results}

\subsection{Comparison with Baselines}
\label{sec:earnhft_comparison}

\tref{tab:earnhft_main} reports the main comparison. EarnHFT achieves the highest profit (TR) on all four datasets and the highest risk-adjusted profit on three. Value-based methods (DQN, CDQNRP) perform acceptably only when the validation and test sets are similar under a stable trend. Policy-based methods (PPO, DRA) tend to collapse onto a do-nothing policy---returning a target equal to the current position to avoid commission---even at a learning rate of $10^{-7}$, and fail badly in the bear market. Rule-based methods are highly sensitive to their take-profit and stop-loss thresholds and earn only moderate profit in volatile markets. EarnHFT, by contrast, is a deliberately aggressive trader because the optimal-value supervisor and optimal actor convey profit-oriented experience; it consequently trails on some risk metrics while dominating on profit. \fref{fig:earnhft_trading} shows a representative episode: EarnHFT opens and closes a position within thirty seconds, profiting from a movement that appears as a mere pullback at minute scale.

\begin{table}[!htb]
\centering
\caption[EarnHFT main results.]{Performance comparison on four cryptocurrency markets against six baselines (two policy-based, two value-based, two rule-based RL/rule methods). Arrows indicate the preferred direction; \fst{pink} marks the best entry per market and metric. TR, AVOL, and MDD are in percent.}
\label{tab:earnhft_main}
\renewcommand{\arraystretch}{1.15}
\resizebox{\textwidth}{!}{
\begin{tabular}{llcccccc}
\toprule
Market & Model & TR$\uparrow$ & ASR$\uparrow$ & ACR$\uparrow$ & ASoR$\uparrow$ & AVOL$\downarrow$ & MDD$\downarrow$ \\
\midrule
\multirow{7}{*}{BTCU}
 & DRA    & $-4.56$ & $-4.28$ & $-19.57$ & $-4.65$ & $42.25$ & $9.24$ \\
 & PPO    & $-3.61$ & $-5.25$ & $-22.74$ & $-5.71$ & $27.76$ & $6.41$ \\
 & CDQNRP & $-2.83$ & $-2.91$ & $-14.85$ & $-3.31$ & $37.61$ & $7.38$ \\
 & DQN    & $-3.48$ & $-12.37$ & $-35.01$ & $-11.86$ & \fst{$11.57$} & $4.09$ \\
 & MACD   & $-6.07$ & $-10.11$ & $-25.17$ & $-8.05$ & $24.84$ & $9.98$ \\
 & IV     & $-2.99$ & $-3.78$ & $-14.24$ & $-3.27$ & $31.35$ & $8.32$ \\
 & EarnHFT & \fst{$0.72$} & \fst{$1.22$} & \fst{$10.77$} & \fst{$0.93$} & $27.08$ & \fst{$3.07$} \\
\midrule
\multirow{7}{*}{BTCT}
 & DRA    & $-2.65$ & $-4.82$ & $-17.48$ & $-4.77$ & $21.18$ & $5.84$ \\
 & PPO    & $-0.60$ & $-14.74$ & $-35.80$ & $-0.14$ & $1.59$ & $0.65$ \\
 & CDQNRP & $-0.60$ & $-19.52$ & $-37.88$ & $-0.74$ & \fst{$1.20$} & \fst{$0.61$} \\
 & DQN    & $0.47$ & \fst{$4.21$} & \fst{$27.94$} & \fst{$1.14$} & $4.38$ & $0.66$ \\
 & MACD   & $-4.02$ & $-5.80$ & $-24.21$ & $-4.32$ & $26.87$ & $6.44$ \\
 & IV     & $-12.01$ & $-17.83$ & $-38.99$ & $-13.90$ & $27.68$ & $12.66$ \\
 & EarnHFT & \fst{$0.99$} & $1.34$ & $7.76$ & $1.01$ & $32.40$ & $5.61$ \\
\midrule
\multirow{7}{*}{ETH}
 & DRA    & $-33.37$ & $-9.06$ & $-32.23$ & $-9.20$ & $163.25$ & $45.88$ \\
 & PPO    & $-22.61$ & $-10.11$ & $-31.17$ & $-10.39$ & $96.12$ & $31.17$ \\
 & CDQNRP & $-6.82$ & $-24.41$ & $-40.19$ & $-3.11$ & \fst{$11.46$} & \fst{$6.96$} \\
 & DQN    & $-11.02$ & $-9.47$ & $-32.81$ & $-8.43$ & $47.76$ & $13.79$ \\
 & MACD   & $-4.29$ & $-1.78$ & $-8.71$ & $-1.19$ & $79.64$ & $16.35$ \\
 & IV     & $-27.42$ & $-12.27$ & $-36.01$ & $-9.00$ & $99.67$ & $33.96$ \\
 & EarnHFT & \fst{$4.52$} & \fst{$2.92$} & \fst{$14.30$} & \fst{$1.78$} & $67.92$ & $13.89$ \\
\midrule
\multirow{7}{*}{GALA}
 & DRA    & $10.56$ & $4.77$ & $41.63$ & $4.44$ & $92.43$ & $10.60$ \\
 & PPO    & $10.56$ & $4.77$ & $41.63$ & $4.44$ & $92.43$ & $10.60$ \\
 & CDQNRP & $5.22$ & $4.51$ & $39.42$ & $4.16$ & $47.27$ & $5.41$ \\
 & DQN    & $2.94$ & $3.55$ & $32.02$ & $2.66$ & \fst{$34.08$} & \fst{$3.78$} \\
 & MACD   & $2.37$ & $1.79$ & $11.45$ & $1.22$ & $62.89$ & $9.84$ \\
 & IV     & $13.95$ & $6.74$ & $55.67$ & $5.41$ & $81.79$ & $9.91$ \\
 & EarnHFT & \fst{$19.41$} & \fst{$9.77$} & \fst{$79.08$} & \fst{$7.79$} & $74.94$ & $9.26$ \\
\bottomrule
\end{tabular}}
\end{table}

\begin{figure}[!htb]
\centering
\includegraphics[width=0.6\textwidth]{Chapter4/earnhft_trading_example}
\caption[An EarnHFT trading episode.]{A representative EarnHFT trading episode on ETH. The agent opens and closes a position within roughly thirty seconds, capturing a movement that registers only as a pullback at the minute scale.}
\label{fig:earnhft_trading}
\end{figure}

\subsection{Effectiveness of the Hierarchical Framework}
\label{sec:earnhft_hierarchy}

\fref{fig:earnhft_hierarchy} compares the router with each specialist in its pool. Bull and rally specialists tend to buy and hold and excel in the bull market (GALA); sideways specialists trade little and hold; pullback and bear specialists tend to close positions and excel in the bear market (ETH). The router combines these strengths and delivers the best profit on all four datasets. \fref{fig:earnhft_selection} shows the router's selection distribution: high-volatility datasets (ETH, GALA) elicit a balanced distribution across trends because dynamics change frequently, whereas low-volatility datasets (BTCT, BTCU) concentrate on a few dynamics.

\begin{figure}[!htb]
\centering
\includegraphics[width=0.85\textwidth]{Chapter4/earnhft_hierarchical_comparison}
\caption[Router versus specialist pool.]{Comparison of the router against each specialist in its pool across the four datasets. Each specialist excels in the trend it was trained for; the router combines their strengths.}
\label{fig:earnhft_hierarchy}
\end{figure}

\begin{figure}[!htb]
\centering
\includegraphics[width=0.7\textwidth]{Chapter4/earnhft_trend_selection}
\caption[Router selection distribution.]{Distribution of the router's specialist selections across the four datasets. High-volatility markets elicit a more balanced distribution; low-volatility markets concentrate on fewer dynamics.}
\label{fig:earnhft_selection}
\end{figure}

\subsection{Effectiveness of the Optimal Action Value}
\label{sec:earnhft_ablation}

\tref{tab:earnhft_ablation} ablates the optimal value supervisor (OS) and optimal actor (OA) on ETH and GALA, measuring the steps to converge (CS), the converged reward sum (RS), and the average holding length (AHL). On GALA, adding OS to plain DDQN reduces the steps to converge to about 15\% while raising the return; OA further raises the return at the cost of more steps. On ETH, a bull market, OS does not reduce CS but substantially raises the return. OS is the more effective component because it supplies information about all actions and adapts as the policy changes, whereas OA only provides demonstrations.

\begin{table}[!htb]
\setlength\tabcolsep{5pt}
\centering
\caption[Ablation of the Q-teacher.]{Ablation of the optimal value supervisor (OS) and optimal actor (OA) on ETH and GALA. CS is steps to converge, RS the converged reward sum, AHL the average holding length.}
\label{tab:earnhft_ablation}
\begin{tabular}{cccccccc}
\toprule
\multirow{2}{*}{OA} & \multirow{2}{*}{OS} & \multicolumn{3}{c}{GALA} & \multicolumn{3}{c}{ETH} \\
\cmidrule(lr){3-5}\cmidrule(lr){6-8}
 & & CS$\downarrow$ & RS$\uparrow$ & AHL & CS$\downarrow$ & RS$\uparrow$ & AHL \\
\midrule
\checkmark & \checkmark & $78848$ & \fst{$4.43$} & $448$ & $102400$ & \fst{$12.32$} & $81.3$ \\
\checkmark &            & $102400$ & $0.24$ & $38.7$ & $102400$ & $-1.40$ & $4.15$ \\
           & \checkmark & \fst{$4608$} & $2.89$ & $147$ & \fst{$30720$} & $4.87$ & $35.8$ \\
           &            & $30720$ & $-0.01$ & $284$ & $30720$ & $-29.6$ & $39.1$ \\
\bottomrule
\end{tabular}
\end{table}

\subsection{Trading Behaviour and Computational Cost}
\label{sec:earnhft_behaviour}

The trading behaviour of EarnHFT varies with the underlying market in a manner consistent with the unit cost of each asset and the stability of its trend. On the two Bitcoin datasets (BTCT and BTCU), the high unit price makes each round trip relatively expensive in dollar terms, and the agent trades less frequently and holds positions only when the price excursion is comfortably larger than the round-trip commission. On ETH, where the trend reverses sharply within the test window, the agent trades much more often to capture the short pullbacks that the minute-level router exposes. On GALA, which trends upward through most of the test period without major reversals, the agent holds positions for long stretches and lets the trend do the work. In aggregate, EarnHFT trades less frequently in stable bull markets and more frequently in volatile or bear markets---behaviour that matches what a discretionary high-frequency trader would do, but emerges from the hierarchy without any hand-coded regime detector.

\paragraph{Computational cost.} All experiments run on a single NVIDIA RTX~4090 GPU. The full pipeline---building the Q-teacher table, training the $200$ specialist agents under the four preference values $\beta\in\{-90,-10,30,100\}$, segmenting and labelling the validation set, selecting the pool, and training the minute-level DDQN router---completes in approximately ten hours for each market, dominated by stage II specialist training. At inference time the only on-policy work is one specialist DDQN forward pass per second and one router forward pass per minute, both of which fit comfortably inside the second-level decision budget on a commodity GPU.

%----------------------------------------------------------------------------------------
\section{Chapter Summary}
\label{sec:earnhft_summary}

This chapter presented EarnHFT, a three-stage hierarchical RL framework that makes high-frequency trading tractable for RL. It addresses the extreme horizon with a dynamic-programming Q-teacher that, under the exogenous-price assumption of \sref{sec:earnhft_stage1} and a fixed historical training window, preserves the optimal Bellman fixed point of the supervised update while accelerating learning empirically; and it addresses regime non-stationarity by training a pool of trend specialists and a minute-level router that switches among them. EarnHFT outperforms six strong baselines across four markets, exceeding the runner-up by at least 30\% in profitability, and the ablations confirm the contribution of each component. In the language of the thesis, EarnHFT is the first instantiation of the decompose--specialise--select pattern, and its selector---a DDQN router trained by reinforcement learning against the trading reward---is the most tightly coupled of the three. The next chapter retains the pattern but confronts a setting, leveraged futures, where the router must additionally manage risk and the agent must become aware of its own competence boundary.
 % FineFT: Risk-Aware Ensemble RL for Futures Trading
\part{Reinforcement Learning for Signal Generation}
%!TEX root=../mythesis.tex
% Chapter 5 -- AlphaQD

\chapter{AlphaQD: Quality-Diversity Reinforcement Learning for Formulaic Alpha Mining} % Main chapter title
\chaptermark{AlphaQD for Alpha Mining}
\label{ch:alphaqd}

%----------------------------------------------------------------------------------------
\section{Introduction}
\label{sec:alphaqd_intro}

The two preceding chapters made trading \emph{decisions}: an agent observed the market and chose how to adjust its position. This chapter changes the task. Rather than acting in the market, the agent must \emph{generate} the predictive signals from which a portfolio is later built. The signals are \emph{formulaic alphas}: short symbolic expressions---trees of arithmetic, cross-sectional, and time-series operators over a small vocabulary of price and volume features---whose cross-sectional value at each timestamp is treated as a forecast of future returns~\cite{kakushadze2016101}. A portfolio is assembled from a \emph{set} of such alphas.

Formulaic alphas remain the dominant signal-generation paradigm in low-frequency quantitative investment for two reasons that monolithic neural predictors do not satisfy: each formula is readable and auditable to risk committees and regulators, and a small set of complementary formulas degrades more gracefully than a single black box when the regime shifts. Automating their discovery, however, is a hard exploration problem under three simultaneous pressures: \emph{quality} (the supervisory signal is the information coefficient of a finished formula on noisy data, with a low signal-to-noise ratio relative to the legal expression space); \emph{diversity} (what is deployed is a set covering different return sources, so a high-quality duplicate is worthless); and \emph{generalisation} (the training window is bounded and the test market non-stationary).

\paragraph{Diagnosis.} The reinforcement-learning line descended from AlphaGen~\cite{yu2023alphagen} scores each new alpha by its marginal contribution to a running pool. This is structurally a boosting objective, and it has two pathologies, the third gap of \sref{sec:bg_gaps}. First, the reward changes every time the pool changes, so the policy is trained against a moving target---a severe problem for exploration, where the policy never sees a stable landscape long enough to commit to coverage. Second, residual-driven boosting carries no a-priori bound on the generalisation of the resulting ensemble: with a legal expression space exceeding $10^{20}$, the late-round residual maximiser is dominated by the search-space noise floor, and the retained set inherits no concentration guarantee tied to its size. \fref{fig:alphaqd_reward_drift} contrasts the pool-dependent AlphaGen reward with the static reward AlphaQD uses.

\begin{figure}[!htb]
\centering
\includegraphics[width=0.8\textwidth]{Chapter5/alphaqd_reward_drift}
\caption[Reward drift under two reward functions.]{Per-candidate reward across pool-insertion steps under two reward functions: the AlphaGen-style pool-dependent reward (left) and the AlphaQD terminal reward (right). Rows are twelve probe alphas; columns are pool-insertion steps; colour encodes the reward (separate scale per panel). The AlphaGen reward shifts as the pool grows; the AlphaQD reward does not.}
\label{fig:alphaqd_reward_drift}
\end{figure}

\paragraph{Remedy.} AlphaQD removes the coupling structurally, in three stages (\fref{fig:alphaqd_pipeline}), following the decompose--specialise--select pattern with the selector fully detached from the reward.
\begin{enumerate}
\item \textbf{Efficient GFlowNet exploration.} A relational-graph GFlowNet samples symbolic expressions and is trained by trajectory balance against a terminal reward that depends \emph{only} on the candidate's training-window information ratio, so the learning target is invariant to later archive updates.
\item \textbf{Return-source MAP-Elites archiving.} Each candidate is back-tested into a daily IC series; its information ratio scores quality, and the normalised IC series is its behaviour descriptor---an empirical proxy for the formula's return source---placing it in a PCA-discretised MAP-Elites grid.
\item \textbf{Top-$k$ cell selection.} One high-quality representative per occupied cell is exported, and the top $K$ cells form the deployed pool.
\end{enumerate}
On full-market Chinese A-share and U.S.\ benchmarks against eleven baselines, AlphaQD lies in the leading signal-quality cluster on both markets, attains the highest annualised return and the strongest net Sharpe among learned formula-mining methods on A-shares, and shows the smallest seed-to-seed variance and the lowest turnover among learned pool-style methods on both markets.

\begin{figure}[!htb]
\centering
\includegraphics[width=\textwidth]{Chapter5/alphaqd_pipeline}
\caption[The AlphaQD pipeline.]{The three stages of AlphaQD. \emph{Left:} a relational-graph GFlowNet samples a legal symbolic expression from raw inputs and is trained by trajectory balance. \emph{Middle:} the candidate is back-tested into a daily IC series; the information ratio scores quality, the normalised IC series is the return-source descriptor, and the candidate enters a PCA-discretised MAP-Elites grid; the quality score is propagated back to the generator. \emph{Right:} the top $K$ occupied cells export their representatives as the deployed pool, evaluated by a shared ridge back-test.}
\label{fig:alphaqd_pipeline}
\end{figure}

%----------------------------------------------------------------------------------------
\section{Problem Formulation}
\label{sec:alphaqd_formulation}

Let $T$ be the number of trading days in the training window, $N$ the cross-sectional universe size, and $X\in\mathbb{R}^{T\times N\times d}$ the per-asset feature tensor (open, high, low, close, volume, VWAP, and short-window returns). A formulaic alpha is a function $\alpha:\mathbb{R}^{T\times N\times d}\to\mathbb{R}^{T\times N}$ realised as a reverse-Polish-notation (RPN) token sequence over a typed grammar $\mathcal{G}$, with an action mask rejecting any prefix that cannot complete into a legal formula. Given the $h$-day forward return $r^{(h)}_{t,i}=\mathrm{close}_{t+h,i}/\mathrm{close}_{t,i}-1$ (we use $h=20$), the daily information coefficient and information ratio are
\begin{equation}
\label{eq:alphaqd_ic}
\mathrm{IC}_t(\alpha) = \corr(\alpha_t, r^{(h)}_t),
\qquad
\mathrm{ICIR}(\alpha) = \overline{\mathrm{IC}(\alpha)}\,\big/\,\hat\sigma(\mathrm{IC}(\alpha)).
\end{equation}

\paragraph{Mining goal.} Let $\mathcal{A}$ be the set of legal formulas. The output is a bounded set $\mathcal{P}\subset\mathcal{A}$ maximising the predictive power of a downstream combination of its members on held-out data. The goal is therefore not the per-formula maximiser $\argmax_{\alpha\in\mathcal{A}}\mathrm{ICIR}(\alpha)$, which a single high-ICIR formula would solve, but the joint problem
\begin{equation}
\label{eq:alphaqd_goal}
\argmax_{\mathcal{P}\subset\mathcal{A}}\;\mathbb{E}_{\mathrm{test}}\!\left[\mathrm{IC}(c(\mathcal{P}))\right]
\quad\text{subject to}\quad |\mathcal{P}|\le K,
\end{equation}
where $c(\cdot)$ is a fixed downstream combiner. This imposes two structural requirements: each retained alpha must be of high quality, and members must capture different \emph{return sources} so their combination generalises. We use the centred, unit-normalised daily IC series as the empirical proxy for the return source carried by $\alpha$.

\paragraph{Contrast with AlphaGen.} AlphaGen approximates \eqref{eq:alphaqd_goal} by a boosting relaxation. A pool $\mathrm{Pool}_t$ is grown one element at a time, with the per-step terminal reward the correlation-penalised marginal contribution
\begin{equation}
\label{eq:alphaqd_alphagen}
R_t(\alpha) = \mathrm{ICIR}(\alpha) - \lambda\max_{\alpha_j\in\mathrm{Pool}_t}\bigl|\corr(\mathrm{IC}(\alpha),\mathrm{IC}(\alpha_j))\bigr|.
\end{equation}
Every pool update changes $R_t(\cdot)$, so the policy is trained against a moving target.

\paragraph{Static formulation.} We retain the per-formula MDP $(\mathcal{S},\mathcal{A},P,R,\gamma=1)$, in which $s_t$ is a partial RPN prefix and $a_t$ the next token under the legal-action mask, but we replace \eqref{eq:alphaqd_alphagen} with a terminal reward that depends only on the candidate and the fixed training window. Diversity is moved out of the reward and into a separate selection rule over a behaviour descriptor of $\alpha$. The reward function then does not change as the archive grows, restoring stationarity; and the selection rule operates over a bounded number of behaviour cells, yielding a generalisation-complexity term that scales with the size of the retained set rather than with $|\mathcal{A}|$ (\sref{sec:alphaqd_theory}).

%----------------------------------------------------------------------------------------
\section{Method}
\label{sec:alphaqd_method}

The three stages exchange candidates and scalar scores but never a pool-dependent reshape of the reward. Algorithm~\ref{alg:alphaqd} ties them together.

\subsection{Efficient GFlowNet Exploration}
\label{sec:alphaqd_generator}

\paragraph{Search surface.} A state $s_t$ is the partial RPN prefix and an action $a_t$ the next token; a legal-action mask, derived from stack depth, type consistency, window-parameter constraints, and a terminator rule, confines the search to syntactically and dimensionally valid expressions. Each prefix is materialised as a typed directed graph---nodes tagged by operator, operand, constant, and window type, edges encoding operator-to-argument relations---and encoded by a two-layer relational graph convolutional network~\cite{schlichtkrull2018rgcn} of hidden width $128$. The pooled node embeddings feed two heads, $\log P_F(\cdot\mid s_t)$ over legal next tokens and $\log P_B(\cdot\mid s_{t+1})$ over legal predecessors, and a scalar $\log Z$.

\paragraph{Terminal reward.} For a completed expression $\alpha$, AlphaQD evaluates its factor values on the training window, cross-sectionally winsorises and rank-normalises by date, computes $\mathrm{IC}(\alpha)$ and $\mathrm{ICIR}(\alpha)$ via \eqref{eq:alphaqd_ic}, and sets
\begin{equation}
\label{eq:alphaqd_reward}
R(\alpha) = \max\!\bigl\{\exp(\beta_q\,\mathrm{ICIR}(\alpha)),\,\exp(-10)\bigr\},
\end{equation}
with quality temperature $\beta_q$ and an exponential floor for unevaluable candidates. Equation~\eqref{eq:alphaqd_reward} reads neither the archive nor any pool statistic---not coverage, novelty, cell identity, incumbent quality, or acceptance status---so for a fixed training window and fixed $\beta_q$ the reward of a candidate is invariant to later archive updates.

\paragraph{Trajectory-balance training.} The generator is trained by trajectory balance~\cite{malkin2022trajectorybalance,bengio2021gflownet}: for a terminal trajectory $\tau$ with formula $\alpha(\tau)$,
\begin{equation}
\label{eq:alphaqd_tb}
\mathcal{L}_{\mathrm{TB}}(\tau) = \left(\log Z + \sum_{t=0}^{T-1}\log P_F(a_t\mid s_t) - \sum_{t=0}^{T-1}\log P_B(s_t\mid s_{t+1}) - \log R(\alpha(\tau))\right)^2.
\end{equation}
At the optimum, the probability of sampling a terminal expression under $P_F$ is proportional to $R(\alpha)$, so the policy assigns mass to every high-reward mode rather than collapsing onto the single maximiser; under a fixed budget, this mode-covering sampling is the ``efficient'' part of the stage~\cite{bengio2021gflownet}. Failed evaluations receive the floor reward and still contribute a finite gradient, so the policy learns to avoid them without a separate invalid-expression branch.

\subsection{Return-Source MAP-Elites Archiving}
\label{sec:alphaqd_archive}

\paragraph{Quality and behaviour descriptor.} The IC series $\mathrm{IC}(\alpha)\in\mathbb{R}^T$ supplies both the quality score $\mathrm{ICIR}(\alpha)$ and a behaviour descriptor that summarises \emph{when} the signal is strong. AlphaQD centres and unit-normalises the IC series,
\begin{equation}
\label{eq:alphaqd_shape}
\mathbf{r}(\alpha) = \frac{\mathbf{x}(\alpha)-\overline{\mathbf{x}(\alpha)}\mathbf{1}}{\|\mathbf{x}(\alpha)-\overline{\mathbf{x}(\alpha)}\mathbf{1}\|_2+\varepsilon},
\qquad \mathbf{x}(\alpha)=\mathrm{IC}(\alpha),
\end{equation}
removing level and scale (absorbed into ICIR) and keeping only the relative temporal pattern. By construction $\mathbf{r}(\alpha)$ is invariant to any positive affine transform of the IC series, so it isolates which dates the formula explains, not how strong its edge is; we treat it as the empirical proxy for the formula's return source. The shape is projected by an incremental PCA basis $W$ (canonical dimension $d_b=2$) to $\mathbf{z}(\alpha)=W^\top\mathbf{r}(\alpha)$, then discretised by per-axis quantile bin edges into a $B\times B$ grid ($B=10$, so $100$ cells), giving the cell index $c(\alpha)$.

\paragraph{Per-cell ranking and update.} Within each cell, elites are ranked by two terms: predictive quality (the ICIR) and a \emph{perturbation fidelity score} (PFS) that jitters each integer window parameter by $\pm 10\%$ over five perturbations and measures the normalised agreement between the perturbed and unperturbed IC series, capturing whether quality survives small parameter changes. A candidate with non-positive ICIR is not offered for insertion. Otherwise, if its cell holds fewer than $\kappa$ elites (canonical $\kappa=3$) it enters directly; if the cell is full it enters only if it Pareto-dominates an existing elite in the $(\mathrm{ICIR},\mathrm{PFS})$ plane or strictly improves the cell's combined score, evicting one elite to maintain the cap. The bounded quota prevents dense regions from absorbing the archive, and the Pareto front retains both the strongest and the most robust formulas per return source.

\subsection{Top-$k$ Cell Selection}
\label{sec:alphaqd_export}

For each occupied cell, the representative is the elite with the highest training-window ICIR (ties broken by PFS), which keeps representatives comparable across cells. Let $K^\star=\min(10,\max(5\kappa,5))$ ($K^\star=10$ under $\kappa=3$, matching the pool-style baselines). AlphaQD ranks occupied cells by representative ICIR and exports the top $K^\star$ as the pool $\mathcal{P}^\star$, which is fed to a shared ridge combiner~\cite{hoerl1970ridge} fitted on the training window. Because the per-cell representative inherits the positive-ICIR admission rule and the cells partition formulas by behaviour, the exported pool satisfies the structural premises (positive standalone edges, bounded pairwise correlation) under which the ensemble bound of \sref{sec:alphaqd_theory} applies.

\begin{algorithm}[!htb]
\caption{AlphaQD: three-stage training and export}
\label{alg:alphaqd}
\begin{algorithmic}[1]
\State Initialise the RGCN encoder, $\log P_F$, $\log P_B$, $\log Z$, an empty MAP-Elites archive, and an empty PCA buffer.
\For{episode $= 1,\dots,E$}
  \State \textbf{Stage 1:} sample a batch of terminal trajectories under $P_F$ and the legal-action mask.
  \For{each terminal trajectory $\tau$ with formula $\alpha$}
    \State Evaluate $\alpha$ on the training window; compute $\mathrm{IC}(\alpha)$, $\mathrm{ICIR}(\alpha)$, and $R(\alpha)$ via \eqref{eq:alphaqd_reward}.
    \State \textbf{Stage 2:} update the PCA buffer with $\mathbf{r}(\alpha)$ from \eqref{eq:alphaqd_shape}; refit $W$.
    \If{$\mathrm{ICIR}(\alpha) > 0$}
      \State Compute PFS via $\pm 10\%$ window-parameter perturbations.
      \State Insert $\alpha$ into cell $c(\alpha)=\mathrm{discretize}(W^\top\mathbf{r}(\alpha))$ under the cell update rule.
    \EndIf
    \State Accumulate $\mathcal{L}_{\mathrm{TB}}(\tau)$ from \eqref{eq:alphaqd_tb}.
  \EndFor
  \State Update $\{\log P_F,\log P_B,\log Z\}$ on the accumulated loss.
\EndFor
\State \textbf{Stage 3:} take the highest-ICIR elite per cell as its representative; rank cells by representative ICIR; export the top $K^\star$ as $\mathcal{P}^\star$ and fit the shared ridge combiner.
\end{algorithmic}
\end{algorithm}

%----------------------------------------------------------------------------------------
\section{Generalisation Properties}
\label{sec:alphaqd_theory}

The staging is not merely convenient; it is generalisation-favouring. This section states the central results; the full statements, assumptions, and proofs are collected in \aref{app:alphaqd_theory}. Two recurring assumptions are a sign-symmetric warm-up shape distribution (A1) and strictly positive ridge weights on retained orbit pairs (A2).

The first result formalises why decoupling quality from the shape descriptor is sound.
\begin{lemma}[Affine invariance of the shape descriptor]
\label{lem:alphaqd_affine}
For any candidate $\alpha$ and any positive affine transform $\mathbf{y}=a\,\mathbf{x}(\alpha)+b\mathbf{1}$ with $a>0$, the normalised shape vector built from $\mathbf{y}$ equals $\mathbf{r}(\alpha)$.
\end{lemma}
Thus the descriptor encodes only which dates were strong, with mean and scale absorbed into the quality score, so two alphas with the same temporal pattern but different overall edge occupy the same cell rather than leaking into different ones.

The second result is the structural reason the static reward survives even archive-aware exploration.
\begin{theorem}[Stationary conditional reward]
\label{thm:alphaqd_stationary}
Let the per-alpha reward be $R_q(\alpha)=\exp(\beta_q\,\mathrm{ICIR}(\alpha))$. For any fixed target orbit $\mathcal{O}^\star$ and any kernel $K$ independent of the archive state and training step, the conditional reward $R_{\mathcal{O}^\star}(\alpha)=R_q(\alpha)\,K(\mathcal{O}(c(\alpha)),\mathcal{O}^\star)$ is a stationary terminal reward. A scheduler that varies $\mathcal{O}^\star$ across episodes changes the distribution over conditions but never the reward function attached to any single condition.
\end{theorem}

The third result contrasts the search complexity of AlphaQD with boosting.
\begin{theorem}[Search-complexity contrast]
\label{thm:alphaqd_complexity}
A $K$-step boosting procedure over a candidate set of size $M$ has hypothesis-class complexity scaling as $\mathcal{C}_{\mathrm{boost},K}=\sqrt{K\log M/n_{\mathrm{eff}}}$, while a quality-diversity rule that maps candidates to at most $G$ occupied cells and exports one representative per cell has complexity scaling as $\mathcal{C}_{\mathrm{QD},K}=\sqrt{K\log(G/K)/n_{\mathrm{eff}}}$. For $G\ll M$, $\mathcal{C}_{\mathrm{QD},K}\ll\mathcal{C}_{\mathrm{boost},K}$.
\end{theorem}
In the reported grid $G\le 100$ while $|\mathcal{A}|>10^{20}$, so the gap spans many orders of magnitude.

The final result lower-bounds the quality of the exported ensemble.
\begin{theorem}[Ensemble ICIR lower bound]
\label{thm:alphaqd_ensemble}
If the export retains $K$ positive-ICIR elites with daily IC variables $X_1,\dots,X_K$ satisfying $\mathbb{E}[X_i]\ge\mu_{\min}>0$, $\Var(X_i)\le\sigma^2$, and $\corr(X_i,X_j)\le\rho$, then the equal-weight ensemble $X_{\mathrm{ens}}=\frac{1}{K}\sum_i X_i$ obeys
\begin{equation}
\label{eq:alphaqd_ensemble}
\mathrm{ICIR}(X_{\mathrm{ens}}) \ge \frac{\mu_{\min}\sqrt{K}}{\sigma\sqrt{1+(K-1)\rho}}.
\end{equation}
\end{theorem}
The bound rests on two premises that AlphaQD supplies by construction: the positive-ICIR admission rule guarantees $\mu_i>0$, and the bounded per-cell quota with the PCA descriptor controls $\rho$ by spreading elites across cells (\sref{sec:alphaqd_diagnostics} reports an empirical pairwise correlation of $0.2478$). Boosting-style fitting preserves neither premise, so no analogous bound is available for it.

%----------------------------------------------------------------------------------------
\section{Experimental Setup}
\label{sec:alphaqd_setup}

\paragraph{Datasets.} We evaluate on two markets so the comparison is not anchored to one regime. The Chinese A-share panel (the canonical alpha-mining setting) is the primary benchmark: full-market BaoStock daily data, $7{,}207$ instruments over $3{,}963$ calendar days from 2010-01-04 to 2026-04-30. The U.S.\ panel is an independent check with different microstructure and base rates: a Tiingo Russell-3000-style universe of $2{,}571$ instruments over $4{,}109$ calendar days. All methods share the same split: 2010--2021 for training, 2022 for validation, and 2023-01-01 to 2026-04-30 for testing, with forward-return horizon $h=20$ days. The training window drives the search; the validation window fits only the shared ridge combiner; the test window is touched once per (method, seed).

\paragraph{Compared methods.} We compare against eleven baselines in four families: supervised (Alpha101~\cite{kakushadze2016101}, MLP~\cite{murtagh1991mlp}, XGBoost~\cite{chen2016xgboost}, LightGBM~\cite{ke2017lightgbm}); genetic programming (GP single-predictor; GP with mutual-IC filter; GP without filter~\cite{stephens2015gplearn}); and generative/RL (AlphaForge~\cite{shi2025alphaforge}, AlphaGen-PPO~\cite{yu2023alphagen}, AlphaQCM~\cite{zhu2025alphaqcm}, AlphaSAGE~\cite{chen2025alphasage}). Pool-style baselines export ten alphas each, matching AlphaQD's $K^\star=10$.

\paragraph{Unified protocol.} Every pool-style method, including AlphaQD, exports a fixed-size pool fed to the same closed-form ridge combiner with weights fitted on training and held fixed at test, so the table reports the quality of the discovered set rather than method-specific weighting. We report signal-layer IC, ICIR, and RankIC together with portfolio-layer annualised net Sharpe, annualised net return, maximum drawdown, and annualised turnover. The portfolio is long-only top-$50$ rebalanced every twenty trading days, executed at $\mathrm{VWAP}_{t+1}$ on A-shares and $\mathrm{Open}_{t+1}$ on the U.S.\ panel, with round-trip costs of $12.5$ and $3.5$ basis points respectively.

%----------------------------------------------------------------------------------------
\section{Results and Analysis}
\label{sec:alphaqd_results}

\subsection{Main Results}
\label{sec:alphaqd_main}

\tref{tab:alphaqd_main} reports both panels under the unified-ridge protocol, aggregated over seeds $\{0,1,2,12345\}$ where stochastic.

\begin{table}[!htb]
\centering
\caption[AlphaQD main results.]{Main results on the 2023-01-01 to 2026-04-30 test window under the unified-ridge protocol. Entries are mean $\pm$ population standard deviation over seeds $\{0,1,2,12345\}$ where stochastic. Portfolio metrics are net of round-trip costs on a long-only top-$50$ portfolio rebalanced every twenty trading days. \fst{First}, \snd{second}, and \trd{third} per column within each panel.}
\label{tab:alphaqd_main}
\renewcommand{\arraystretch}{1.15}
\footnotesize
\resizebox{\textwidth}{!}{
\begin{tabular}{llccccccc}
\toprule
Family & Method & IC$\uparrow$ & ICIR$\uparrow$ & RankIC$\uparrow$ & Sharpe$\uparrow$ & AnnRet(\%)$\uparrow$ & MDD(\%)$\uparrow$ & TO/yr$\downarrow$ \\
\midrule
\multicolumn{9}{c}{\textit{Panel A: Chinese A-share}} \\
\midrule
Formulaic & Alpha101 & $0.0487$ & $0.408$ & $0.082$ & \fst{$0.983$} & \snd{$19.9$} & \fst{$-26.1$} & $23.9$ \\
\multirow{3}{*}{Supervised} & MLP & $0.0079$ & $0.111$ & $0.017$ & $0.319$ & $9.1$ & $-49.7$ & $12.9$ \\
 & LightGBM & $0.0189$ & $0.262$ & $0.035$ & $0.309$ & $10.0$ & $-52.7$ & \snd{$12.0$} \\
 & XGBoost & $0.0199$ & $0.266$ & $0.036$ & $0.236$ & $6.9$ & $-50.6$ & \trd{$12.4$} \\
\multirow{3}{*}{Genetic Prog.} & GP (predictor) & $-0.0101$ & $-0.082$ & $-0.016$ & $0.292$ & $6.4$ & $-40.4$ & \fst{$3.8$} \\
 & GP w/ filter & \trd{$0.0652$} & $0.507$ & \snd{$0.105$} & $0.663$ & \trd{$18.6$} & $-39.1$ & $15.1$ \\
 & GP w/o filter & \fst{$0.0703$} & \fst{$0.581$} & \trd{$0.102$} & $0.501$ & $10.8$ & $-35.7$ & $20.2$ \\
\multirow{4}{*}{Generative / RL} & AlphaForge & $0.0636$ & \trd{$0.553$} & $0.086$ & $0.534$ & $14.4$ & $-37.0$ & $19.5$ \\
 & AlphaGen-PPO & $0.0435$ & $0.495$ & $0.049$ & $0.275$ & $4.1$ & $-43.3$ & $22.2$ \\
 & AlphaQCM (FQF) & $0.0206$ & $0.255$ & $0.020$ & $0.053$ & $1.7$ & $-46.8$ & $19.3$ \\
 & AlphaSAGE (GNN) & $0.0552$ & \snd{$0.555$} & $0.095$ & \trd{$0.730$} & $14.9$ & \snd{$-27.4$} & $18.3$ \\
Ours & \textbf{AlphaQD} & \snd{$0.0673$} & $0.522$ & \fst{$0.112$} & \snd{$0.911$} & \fst{$24.7$} & \trd{$-31.5$} & $12.5$ \\
\midrule
\multicolumn{9}{c}{\textit{Panel B: U.S.\ Russell-3000-style universe}} \\
\midrule
Formulaic & Alpha101 & $0.0188$ & $0.358$ & \snd{$-0.002$} & $1.511$ & $38.2$ & \snd{$-15.9$} & $23.0$ \\
\multirow{3}{*}{Supervised} & MLP & $0.0446$ & \fst{$0.604$} & \trd{$-0.002$} & $1.796$ & $99.8$ & \trd{$-16.3$} & $9.1$ \\
 & LightGBM & $0.0404$ & $0.454$ & $-0.009$ & $2.090$ & $119.1$ & $-21.6$ & $8.7$ \\
 & XGBoost & $0.0335$ & $0.395$ & $-0.006$ & $2.120$ & $129.2$ & $-23.4$ & $9.6$ \\
\multirow{3}{*}{Genetic Prog.} & GP (predictor) & $-0.0104$ & $-0.073$ & \fst{$0.004$} & $1.202$ & $39.4$ & \fst{$-15.2$} & \fst{$5.7$} \\
 & GP w/ filter & \snd{$0.0691$} & $0.512$ & $-0.005$ & \fst{$2.180$} & \fst{$141.6$} & $-26.7$ & $9.5$ \\
 & GP w/o filter & \fst{$0.0717$} & \snd{$0.568$} & $-0.013$ & \trd{$2.164$} & \snd{$141.3$} & $-24.7$ & $10.2$ \\
\multirow{4}{*}{Generative / RL} & AlphaForge & $0.0481$ & $0.430$ & $-0.008$ & $1.856$ & $107.0$ & $-23.4$ & $11.6$ \\
 & AlphaGen-PPO & $0.0542$ & $0.437$ & $-0.010$ & $2.056$ & $129.8$ & $-23.3$ & $13.8$ \\
 & AlphaQCM (FQF) & $0.0486$ & $0.453$ & $-0.002$ & $1.942$ & $103.8$ & $-21.1$ & \trd{$7.9$} \\
 & AlphaSAGE (GNN) & $0.0614$ & $0.539$ & $-0.007$ & $2.111$ & $131.3$ & $-24.2$ & $10.5$ \\
Ours & \textbf{AlphaQD} & \trd{$0.0655$} & \trd{$0.551$} & $-0.010$ & \snd{$2.177$} & \trd{$133.3$} & $-22.5$ & \snd{$6.7$} \\
\bottomrule
\end{tabular}}
\end{table}

\paragraph{A-share findings.} On the A-share panel, AlphaQD attains the highest annualised return and the highest RankIC of any method, the second-highest IC (behind GP-without-filter), and the second-highest Sharpe (behind only the fixed Alpha101 catalogue), while showing the smallest seed-to-seed standard deviation on Sharpe, annualised return, and RankIC within the top quality cluster. That cluster---GP-without-filter, AlphaForge, AlphaSAGE, and AlphaQD---has statistically indistinguishable mean ICIRs, but the three competitors do not convert signal quality into stable portfolio performance to the same extent: GP-without-filter and AlphaForge carry seed-level Sharpe deviations larger than their means, so any single deployment is essentially a coin flip; AlphaSAGE is stable but with annualised return roughly ten points below AlphaQD and turnover 46\% higher. Among learned pool-style baselines AlphaQD also has the lowest turnover. Its gains therefore come from converting signal quality into portfolio performance reliably, not from out-signalling the field. \fref{fig:alphaqd_nav} plots the cumulative net NAV of the five learned pool-style methods.

\begin{figure}[!htb]
\centering
\includegraphics[width=0.8\textwidth]{Chapter5/alphaqd_portfolio_nav}
\caption[Cumulative NAV on the A-share test window.]{Cumulative net NAV (start $=1.0$) on the A-share test window 2023-01-01 to 2026-04-30 for the five learned pool-style methods. Solid lines are the mean NAV across seeds $\{0,1,2,12345\}$ at each rebalance date; shaded bands span $\pm 1$ population standard deviation; the dashed line marks NAV $=1$. The portfolio is the long-only top-$50$ rebalanced every twenty trading days at $\mathrm{VWAP}_{t+1}$ with $12.5$ bps round-trip cost.}
\label{fig:alphaqd_nav}
\end{figure}

\paragraph{U.S.\ findings.} On the U.S.\ panel the signal-quality columns (IC, ICIR, RankIC) are the credible discriminators: RankIC clusters around zero and Sharpe-style numbers around two for all baselines, reflecting market beta in the 2023--2026 window rather than alpha-driven return. AlphaQD ranks third on both IC and ICIR while again showing the smallest seed deviation on every IC-family metric and the lowest turnover among learned-search methods. Critically, the IC magnitude on the U.S.\ panel (around $0.066$) matches the A-share magnitude (around $0.067$), so the method ordering is preserved across two dissimilar markets---a cross-market stability that the boosting-style baselines do not exhibit (for example, AlphaQCM moves from $0.021$ on A-shares to $0.049$ on the U.S.\ panel).

\subsection{Archive Diagnostics}
\label{sec:alphaqd_diagnostics}

On a representative training archive, AlphaQD occupies $97$ of $100$ behaviour cells, stores $282$ elites, and reaches a best in-archive ICIR of $0.7540$. The mean pairwise correlation among stored daily IC series is $0.2478$. Both statistics support the descriptor's role: near-complete occupancy indicates that the normalised IC shape spreads candidates across the grid rather than collapsing into one mode, and the modest correlation indicates that retained formulas explain genuinely different return sources rather than re-encoding one signal---precisely the controlled $\rho$ that \tref{tab:alphaqd_main}'s portfolio stability and Theorem~\ref{thm:alphaqd_ensemble} require. \fref{fig:alphaqd_archive} renders the archive.

\begin{figure}[!htb]
\centering
\includegraphics[width=0.8\textwidth]{Chapter5/alphaqd_archive_heatmap}
\caption[The AlphaQD MAP-Elites archive.]{The AlphaQD MAP-Elites archive on the $10\times 10$ behaviour-descriptor grid (vertical axis PC1 bin, horizontal axis PC2 bin). \emph{(a)} number of stored elites per cell; \emph{(b)} representative ICIR per cell, with the top-$10$ exported cells outlined.}
\label{fig:alphaqd_archive}
\end{figure}

\subsection{Ablations}
\label{sec:alphaqd_ablation}

A grid over behaviour-descriptor dimension $d_b\in\{1,2,3\}$, per-cell capacity $\kappa\in\{1,3,5\}$, and reward temperature $\beta_q\in\{0.5,1,2,5\}$ confirms the design choices. Marginal ICIR declines gently from $d_b=1$ to $d_b=3$, consistent with a higher-dimensional descriptor spreading candidates over more cells; the canonical $d_b=2$ sits in the favourable region while preserving an interpretable two-axis archive. Capacity $\kappa\ge 3$ is preferred over $\kappa=1$, supporting the in-cell Pareto front over a single elite. The reward temperature is non-monotone, peaking at $\beta_q=2$, consistent with the boundary between coverage-preserving and quality-concentrating sampling. Full ablation tables are given in \aref{app:alphaqd_theory}.

%----------------------------------------------------------------------------------------
\section{Summary}
\label{sec:alphaqd_summary}

This chapter presented AlphaQD, a quality-diversity RL method for formulaic alpha mining. Its central choice is structural: the running pool is removed from the generator's reward, and diversity, robustness, and selection are handled by a MAP-Elites archive whose behaviour descriptor is the normalised daily IC series---an empirical proxy for a formula's return source. The training-window reward is thereby invariant to archive updates, and the export rule yields a pool whose effective complexity scales with the number of occupied cells rather than the cardinality of the formula space. On full-market A-share and U.S.\ benchmarks AlphaQD lies in the leading signal cluster on both markets, attains the highest annualised return and net Sharpe among learned formula-mining methods on A-shares, and exhibits the smallest seed variance and lowest turnover among learned pool-style methods. In the arc of the thesis, AlphaQD is the endpoint of the selector-decoupling trajectory: where EarnHFT learns its router and FineFT derives one heuristically from generative uncertainty, AlphaQD removes the selector from the learning reward entirely, and reaps from this decoupling both reward stationarity and a finite-sample generalisation guarantee.
 % AlphaQD: Quality-Diversity RL for Formulaic Alpha Mining
\part{Epilogue}
%!TEX root=../mythesis.tex
% Chapter 6 -- AlphaQD

\chapter{AlphaQD: Quality-Diversity Reinforcement Learning for Formulaic Alpha Mining} % Main chapter title
\chaptermark{AlphaQD for Alpha Mining}
\label{ch:alphaqd}

%----------------------------------------------------------------------------------------
\section{Introduction}
\label{sec:alphaqd_intro}

The two preceding chapters made trading \emph{decisions}: an agent observed the market and chose how to adjust its position. This chapter changes the task.\footnote{This chapter is based on the AlphaQD manuscript, which is currently under review.} Rather than acting in the market, the agent must \emph{generate} the predictive signals from which a portfolio is later built. The signals are \emph{formulaic alphas}: short symbolic expressions whose cross-sectional value at each timestamp is treated as a forecast of future returns~\cite{kakushadze2016101}. Each alpha is a tree of arithmetic, cross-sectional, and time-series operators over a small vocabulary of price and volume features. A portfolio is assembled from a \emph{set} of such alphas.

Formulaic alphas remain widely used in low-frequency quantitative investment because they are readable and can be combined into complementary sets. Automating their discovery is, however, a hard exploration problem under three simultaneous pressures: \emph{quality} (the supervisory signal is the information coefficient of a finished formula on noisy data); \emph{diversity} (a near-duplicate adds little diversification value); and \emph{generalisation} (the training window is bounded and the test market is non-stationary).

\paragraph{Motivation.} Pool-based reinforcement-learning methods such as AlphaGen~\cite{yu2023alphagen} score a candidate by its contribution to an evolving pool. The resulting target changes with the pool, which may complicate optimisation. \fref{fig:alphaqd_reward_drift} illustrates this effect for the correlation-penalised surrogate in \eqref{eq:alphaqd_alphagen}; the theoretical comparison in \sref{sec:alphaqd_theory} concerns this surrogate class rather than the published AlphaGen system.

\begin{figure}[!htb]
\centering
\includegraphics[width=0.8\textwidth]{Chapter6/alphaqd_reward_drift}
\caption[Reward drift under two reward functions.]{Per-candidate reward across pool-insertion steps under two reward functions: the pool-dependent correlation-penalised reward of \eqref{eq:alphaqd_alphagen}, an AlphaGen-style objective (left), and the AlphaQD terminal reward (right). Rows are twelve probe alphas; columns are pool-insertion steps; colour encodes the reward (separate scale per panel). The pool-coupled reward shifts as the pool grows; the AlphaQD reward does not.}
\label{fig:alphaqd_reward_drift}
\end{figure}

\paragraph{Remedy.} AlphaQD removes this coupling in three stages (\fref{fig:alphaqd_pipeline}), following the decompose--specialise--aggregate pattern with the aggregation mechanism fully detached from the reward.
\begin{enumerate}
\item \textbf{GFlowNet exploration.} A relational-graph GFlowNet samples symbolic expressions and is trained by trajectory balance against a terminal reward that depends \emph{only} on the candidate's training-window information ratio, so the learning target is invariant to later archive updates.
\item \textbf{Return-source MAP-Elites archiving.} Each candidate is backtested into a daily IC series; its information ratio scores quality, and the normalised IC series is its behaviour descriptor---an empirical proxy for the formula's return source---placing it in a PCA-discretised MAP-Elites grid~\cite{mouret2015mapelites,pugh2016quality}.
\item \textbf{Top-$K$ cell selection.} One high-quality representative per occupied cell is exported, and the top $K$ cells form the deployed pool.
\end{enumerate}
Against eleven baselines on the reported Chinese A-share and U.S.\ panels, AlphaQD is competitive on the signal metrics. On A-shares it has the highest annualised return among the learned formula-mining methods; among the four generative/RL pool-style methods it has the lowest turnover on both panels.

\begin{figure}[!htb]
\centering
\includegraphics[width=\textwidth]{Chapter6/alphaqd_pipeline}
\caption[The AlphaQD pipeline.]{The three stages of AlphaQD. \emph{Left:} a relational-graph GFlowNet samples a legal symbolic expression and is trained by trajectory balance using a candidate-only reward. \emph{Middle:} the expression is backtested into a daily IC series and inserted into a PCA-discretised MAP-Elites grid. \emph{Right:} the top $K$ occupied cells export their representatives for evaluation with a shared ridge combiner.}
\label{fig:alphaqd_pipeline}
\end{figure}

%----------------------------------------------------------------------------------------
\section{Problem Formulation}
\label{sec:alphaqd_formulation}

Let $T$ be the number of trading days in the training window, $N$ the cross-sectional universe size, and $X\in\mathbb{R}^{T\times N\times d}$ the per-asset feature tensor (open, high, low, close, volume, VWAP, and short-window returns). A formulaic alpha is a function $\alpha:\mathbb{R}^{T\times N\times d}\to\mathbb{R}^{T\times N}$ realised as a reverse-Polish-notation (RPN) token sequence over a typed grammar $\mathcal{G}$, with an action mask rejecting any prefix that cannot complete into a legal formula. Given the $h$-day forward return $r^{(h)}_{t,i}=\mathrm{close}_{t+h,i}/\mathrm{close}_{t,i}-1$ (we use $h=20$), the daily information coefficient and information ratio are, restating the definitions of \sref{sec:bg_metrics} in this chapter's notation~\cite{grinold2000active},
\begin{equation}
\label{eq:alphaqd_ic}
\mathrm{IC}_t(\alpha) = \corr(\alpha_t, r^{(h)}_t),
\qquad
\mathrm{ICIR}(\alpha) = \overline{\mathrm{IC}(\alpha)}\,\big/\,\hat\sigma(\mathrm{IC}(\alpha)).
\end{equation}

\paragraph{Mining goal.} Let $\mathcal{F}$ be the set of legal formulas. The output is a bounded set $\mathcal{P}\subset\mathcal{F}$ whose members, once combined by a fixed downstream combiner, should predict returns on data not seen during mining. The goal is therefore not the per-formula maximiser $\argmax_{\alpha\in\mathcal{F}}\mathrm{ICIR}(\alpha)$, which a single high-ICIR formula would solve, but the joint problem
\begin{equation}
\label{eq:alphaqd_goal}
\argmax_{\mathcal{P}\subset\mathcal{F}}\;\mathbb{E}_{\mathrm{oos}}\!\left[\mathrm{IC}(g(\mathcal{P}))\right]
\quad\text{subject to}\quad |\mathcal{P}|\le K,
\end{equation}
where $g(\cdot)$ is the fixed combiner, $K$ is the pool-size budget, and the expectation is taken over out-of-sample data. The archive-cell map introduced later is denoted separately by $c(\alpha)$. Equation~\eqref{eq:alphaqd_goal} is an idealised target rather than an operational objective: no mining procedure can evaluate the out-of-sample expectation, so search must optimise training-window proxies, and the relation between those proxies and \eqref{eq:alphaqd_goal} is precisely the generalisation question that \sref{sec:alphaqd_theory} addresses. The target imposes two structural requirements: each retained alpha must be of high quality, and members must capture different \emph{return sources} so their combination generalises. We use the centred, normalised daily IC series as the empirical proxy for the return source carried by $\alpha$.

\paragraph{Contrast with AlphaGen.} AlphaGen approximates \eqref{eq:alphaqd_goal} by a boosting relaxation: a pool $\mathrm{Pool}_t$ is grown one element at a time, and the RL reward is the increment in the combined pool's predictive performance after the combination model is refitted with the candidate included. Throughout this chapter we work with a correlation-penalised \emph{variant} that captures the pool coupling of this objective in closed form,
\begin{equation}
\label{eq:alphaqd_alphagen}
R_t(\alpha) = \mathrm{ICIR}(\alpha) - \lambda\max_{\alpha_j\in\mathrm{Pool}_t}\bigl|\corr(\mathrm{IC}(\alpha),\mathrm{IC}(\alpha_j))\bigr|.
\end{equation}
with the correlation penalty defined as zero when $\mathrm{Pool}_t$ is empty. Equation~\eqref{eq:alphaqd_alphagen} is a surrogate for---not a restatement of---the original AlphaGen reward, and the analyses and instrumentation that follow (\fref{fig:alphaqd_reward_drift}, \sref{sec:alphaqd_theory}) address this variant and the pool-coupled reward class it represents rather than the published system. What matters here is the property the variant shares with the original: every pool update changes $R_t(\cdot)$, so the policy is trained against a moving target.

\paragraph{Static formulation.} We retain the per-formula MDP $(\mathcal{S},\mathcal{A},P,R,\gamma=1)$, in which $s_t$ is a partial RPN prefix and $a_t$ the next token under the legal-action mask, but we replace \eqref{eq:alphaqd_alphagen} with a terminal reward that depends only on the candidate and the fixed training window. Diversity is moved out of the reward and into a separate selection rule over a behaviour descriptor of $\alpha$. The reward does not change as the archive grows, removing archive-induced reward drift. The selection rule operates over a bounded number of behaviour cells, yielding a conditional finite-class comparison in \sref{sec:alphaqd_theory}.

%----------------------------------------------------------------------------------------
\section{Method}
\label{sec:alphaqd_method}

The three stages exchange candidates and scalar scores but never a pool-dependent reshape of the reward. Algorithm~\ref{alg:alphaqd} ties them together.

\subsection{GFlowNet Exploration}
\label{sec:alphaqd_generator}

\paragraph{Search surface.} A state $s_t$ is the partial RPN prefix and an action $a_t$ the next token; a legal-action mask, derived from stack depth, type consistency, window-parameter constraints, and a terminator rule, confines the search to syntactically and dimensionally valid expressions. Each prefix is materialised as a typed directed graph---nodes tagged by operator, operand, constant, and window type, edges encoding operator-to-argument relations---and encoded by a two-layer relational graph convolutional network~\cite{schlichtkrull2018rgcn} of hidden width $128$. The pooled node embeddings feed two heads, $\log P_F(\cdot\mid s_t)$ over legal next tokens and $\log P_B(\cdot\mid s_{t+1})$ over legal predecessors, and a scalar $\log Z$.

\paragraph{Terminal reward.} For a completed expression $\alpha$, AlphaQD evaluates its factor values on the training window and cross-sectionally standardises both the factor and the forward-return target to zero mean and unit population standard deviation on each date. Values that remain non-finite after this operation are replaced by zero. The implementation does not winsorise or rank-normalise the factor before computing Pearson IC; this differs from the winsorisation-and-rank-normalisation description in the original method write-up. It then computes $\mathrm{IC}(\alpha)$ and $\mathrm{ICIR}(\alpha)$ via \eqref{eq:alphaqd_ic}, and sets
\begin{equation}
\label{eq:alphaqd_reward}
R(\alpha) = \max\!\bigl\{\exp(\beta_q\,\mathrm{ICIR}(\alpha)),\,\exp(-10)\bigr\},
\end{equation}
with quality temperature $\beta_q$ ($\beta_q=0.5$ in the headline runs; \sref{sec:alphaqd_ablation}) and an exponential floor for an invalid expression, an evaluation failure, or a non-finite or degenerate ICIR. Individual non-finite entries in an otherwise usable daily IC series are zero-filled before the ICIR is formed. The floor also binds for any evaluable candidate with $\mathrm{ICIR}(\alpha)\le -10/\beta_q$, so below that level the generator does not distinguish syntactically valid but extremely poor formulas from unevaluable ones. Equation~\eqref{eq:alphaqd_reward} reads neither the archive nor any pool statistic---not coverage, novelty, cell identity, incumbent quality, or acceptance status---so for a fixed training window and fixed $\beta_q$ the reward of a candidate is invariant to later archive updates.

\paragraph{Trajectory balance.} The generator is trained by trajectory balance~\cite{malkin2022trajectorybalance,bengio2021gflownet}: for a terminal trajectory $\tau$ of token length $L$ (distinct from the trading-day count $T$) with formula $\alpha(\tau)$,
\begin{equation}
\label{eq:alphaqd_tb}
\begin{aligned}
\mathcal{L}_{\mathrm{TB}}(\tau)
=\Biggl(&\log Z+\sum_{t=0}^{L-1}\log P_F(a_t\mid s_t)\\[-0.2em]
&-\sum_{t=0}^{L-1}\log P_B(s_t\mid s_{t+1})
-\log R(\alpha(\tau))\Biggr)^2.
\end{aligned}
\end{equation}
For a finite acyclic construction graph with positive rewards, sufficient model capacity, and an exact trajectory-balance optimum, the terminal sampling probability is proportional to $R(\alpha)$~\cite{bengio2021gflownet}. The trained model only approximates this objective. Failed evaluations receive the floor reward and contribute a finite training signal.

\subsection{Return-Source MAP-Elites Archiving}
\label{sec:alphaqd_archive}

\paragraph{Quality and behaviour descriptor.} The IC series $\mathrm{IC}(\alpha)\in\mathbb{R}^T$ supplies both the quality score $\mathrm{ICIR}(\alpha)$ and a behaviour descriptor that summarises \emph{when} the signal is strong. AlphaQD centres the IC series and normalises it to unit length up to the $\varepsilon$-regularisation in the denominator,
\begin{equation}
\label{eq:alphaqd_shape}
\mathbf{r}(\alpha) = \frac{\mathbf{x}(\alpha)-\overline{\mathbf{x}(\alpha)}\mathbf{1}}{\|\mathbf{x}(\alpha)-\overline{\mathbf{x}(\alpha)}\mathbf{1}\|_2+\varepsilon},
\qquad \mathbf{x}(\alpha)=\mathrm{IC}(\alpha),
\end{equation}
removing level and scale (absorbed into ICIR) and keeping only the relative temporal pattern. By construction $\mathbf{r}(\alpha)$ is invariant to any positive affine transform of the IC series---exactly in the $\varepsilon\to 0$ limit, and up to a negligible $O(\varepsilon)$ term for the machine-precision $\varepsilon$ used in practice (Lemma~\ref{lem:alphaqd_affine})---so it isolates which dates the formula explains, not how strong its edge is; we treat it as the empirical proxy for the formula's return source. The shape is projected by an incremental PCA basis $W$ (canonical dimension $d_b=2$) to $\mathbf{z}(\alpha)=W^\top\mathbf{r}(\alpha)$, then discretised by per-axis quantile bin edges into a $B\times B$ grid ($B=10$, so $100$ cells), giving the cell index $c(\alpha)$. The initial basis and edges are fitted after the first $2{,}000$ \emph{finite candidates}, not after $2{,}000$ episodes. Those candidates are kept pending and are offered to the newly created cells retrospectively after the fit. Thereafter only IC series of \emph{accepted elites} enter a FIFO buffer of at most $1{,}024$ series. In the current AIFML experiment grid, every $K_{\mathrm{pca}}=100$ completed episodes the implementation checks whether at least $32$ accepted series are available; if so, it applies \texttt{IncrementalPCA.partial\_fit} to the buffer, moves each quantile edge towards its new buffer quantile with EMA decay $0.99$, and reprojects every stored elite into the updated cells. When reprojection makes a cell exceed capacity, only the $\kappa$ elites with largest $\mathrm{ICIR}+\mathrm{PFS}$ are retained. Thus both coordinates and membership can change throughout a $2{,}000$-episode run; the descriptor map is not frozen after warm-up.

\paragraph{Per-cell score and update.} For a positive-ICIR candidate, PFS is the mean \emph{ICIR-retention ratio} over the successful perturbations,
\[
\mathrm{PFS}(\alpha)=\operatorname{clip}_{[0,2]}\!\left(\frac{1}{|\mathcal S_\alpha|}\sum_{\widetilde\alpha\in\mathcal S_\alpha}\frac{\mathrm{ICIR}(\widetilde\alpha)}{\mathrm{ICIR}(\alpha)}\right).
\]
The implementation attempts five perturbations. Each rolling-window node is independently replaced by another legal window. It first searches within the configured $\pm10\%$ band, but the legal set is $\{10,20,30,40,50\}$; when that band contains no alternative it samples from the two nearest legal alternatives. Failed or non-finite perturbations are omitted, and PFS is set to $1$ when none succeeds, including formulas with no rolling-window node. Values are clipped to $[0,2]$.

PFS is computed lazily. If a cell is full, a cheap pre-screen first requires the candidate's ICIR to be at least the smallest incumbent ICIR up to $10^{-6}$. A non-positive-ICIR candidate is rejected from the archive in the current configuration and does not update the post-warm-up PCA buffer. A non-full cell accepts an eligible candidate without a dominance check. In a full cell, a candidate that dominates one or more incumbents evicts the dominated incumbent with smallest $\mathrm{ICIR}+\mathrm{PFS}$; a candidate dominated by any incumbent is rejected; otherwise it replaces the incumbent with the smallest sum only when its own sum is strictly larger. This is a bounded score-based cell, not a maintained Pareto front: insertion into a non-full cell can retain dominated members, and one dominance replacement need not remove every dominated member. The quota still prevents a dense region from occupying more than $\kappa$ archive slots per cell.

\subsection{\texorpdfstring{Top-$K$ Cell Selection}{Top-K Cell Selection}}
\label{sec:alphaqd_export}

For each occupied cell, the representative is the elite with the highest training-window ICIR; the export code has no explicit PFS tie-break, so an exact ICIR tie retains iteration order. Let $K^\star=\min(10,5\kappa)$, i.e.\ $5\kappa$ capped at $10$: $K^\star=10$ under the canonical $\kappa=3$, matching the pool-style baselines, but $K^\star=5$ under $\kappa=1$, so the rule couples the exported pool size to the per-cell capacity---a coupling that matters when interpreting the $\kappa$ axis of the sensitivity grid in \sref{sec:alphaqd_ablation}. AlphaQD ranks occupied cells by representative ICIR and exports the top $K^\star$ as the pool $\mathcal{P}^\star$. The main comparison then uses the shared closed-form ridge combiner~\cite{hoerl1970ridge},
\[
\widehat{\mathbf w}=(X^\top X+I)^{-1}X^\top\mathbf y,
\]
fitted to tradable validation-split asset--date rows with finite targets. There is no intercept or feature standardisation; missing factor values are replaced by zero. The regularisation coefficient is fixed at $\lambda=1.0$, and the resulting weights are held fixed on the test set. Positive-ICIR-proportional weights are written with the native AlphaQD artifact, but they are not the weights used by the unified-ridge headline table. Because the per-cell representative inherits the positive-ICIR admission rule and the cells partition formulas by behaviour, the export is designed to encourage positive standalone training-window edge and lower redundancy; whether the exported pool meets the formal premises of \sref{sec:alphaqd_theory} out of sample is not guaranteed (see the discussion after Theorem~\ref{thm:alphaqd_ensemble}).

\begin{algorithm}[!htb]
\caption{AlphaQD: three-stage training and export}
\label{alg:alphaqd}
\begin{algorithmic}[1]
\State Initialise the RGCN encoder, $\log P_F$, $\log P_B$, $\log Z$, an empty MAP-Elites archive, and an empty PCA buffer.
\For{episode $= 1,\dots,E$}
  \State \textbf{Stage 1:} sample a batch of terminal trajectories under $P_F$ and the legal-action mask.
  \For{each terminal trajectory $\tau$ with formula $\alpha$}
    \State Evaluate $\alpha$ on the training window; compute $\mathrm{IC}(\alpha)$, $\mathrm{ICIR}(\alpha)$, and $R(\alpha)$ via \eqref{eq:alphaqd_reward}.
    \State \textbf{Stage 2:} update the descriptor and bounded archive as follows.
    \If{the first $2{,}000$ finite candidates have not yet been collected}
      \State Append $\mathbf{r}(\alpha)$ to the warm-up set and retain $\alpha$ as pending; at the threshold, fit $W$ and the initial quantile edges and offer all pending positive-ICIR candidates to their cells.
    \ElsIf{$\mathrm{ICIR}(\alpha)>0$ and the cell ICIR pre-screen passes}
      \State Compute PFS from successful legal-window perturbations and apply the bounded cell update rule.
      \If{$\alpha$ is accepted}
        \State Append its IC series to the $1{,}024$-element FIFO PCA buffer.
      \EndIf
    \EndIf
    \State Accumulate $\mathcal{L}_{\mathrm{TB}}(\tau)$ from \eqref{eq:alphaqd_tb}.
  \EndFor
  \State Update $\{\log P_F,\log P_B,\log Z\}$ on the accumulated loss.
  \If{the episode count is a multiple of $K_{\mathrm{pca}}=100$ and the FIFO buffer has at least $32$ series}
    \State Incrementally update $W$, move the quantile edges by EMA, reproject all elites, and trim over-capacity cells by $\mathrm{ICIR}+\mathrm{PFS}$.
  \EndIf
\EndFor
\State \textbf{Stage 3:} take the highest-ICIR elite per cell as its representative; rank cells by representative ICIR; export the top $K^\star$ as $\mathcal{P}^\star$ and fit the shared $\lambda=1$ ridge combiner on validation rows.
\end{algorithmic}
\end{algorithm}

%----------------------------------------------------------------------------------------
\section{Conditional Analytical Properties}
\label{sec:alphaqd_theory}

This section gives conditional analytical results, with proofs inline. The deployed descriptor map is data-dependent and changes at the periodic PCA updates described in \sref{sec:alphaqd_archive}; none of the results below treats that moving partition as fixed unless the premise says so explicitly. One piece of notation is needed by Theorem~\ref{thm:alphaqd_stationary}: for a cell index $c=(b_1,\dots,b_{d_b})$ with $b_j\in\{0,\dots,B-1\}$, the \emph{anti-cell} is $\mathrm{anti}(c)=(B-1-b_1,\dots,B-1-b_{d_b})$, the \emph{orbit} of $c$ is the unordered pair $\mathcal{O}(c)=\{c,\mathrm{anti}(c)\}$, and, by a mild abuse of notation, $\mathcal{O}$ also denotes the set of all orbits. A formula and its sign-flip have oppositely signed IC shapes, but empirical quantile edges---especially after EMA updates---need not be symmetric, so antipodal cell membership is an idealised condition rather than a deployed invariant. Throughout, the discussion after each theorem distinguishes such fixed-map statements from the implementation.

\subsection{Decoupling Quality from Shape}
\label{sec:alphaqd_decouple}

The first result formalises why decoupling quality from the shape descriptor is sound.

\begin{lemma}[Affine invariance of the shape descriptor]
\label{lem:alphaqd_affine}
For a candidate $\alpha$, write $\mathbf{v}=\mathbf{x}(\alpha)-\overline{\mathbf{x}(\alpha)}\mathbf{1}$ and suppose $\|\mathbf{v}\|_2>0$. For any positive affine transform $\mathbf{y}=a\,\mathbf{x}(\alpha)+b\mathbf{1}$ with $a>0$, let $\mathbf{r}_{\mathbf{y}}$ denote the shape vector \eqref{eq:alphaqd_shape} built from $\mathbf{y}$. Then $\mathbf{r}_{\mathbf{y}}=\mathbf{r}(\alpha)$ exactly in the limit $\varepsilon\to 0$, and for $\varepsilon>0$,
\[
\bigl\|\mathbf{r}_{\mathbf{y}}-\mathbf{r}(\alpha)\bigr\|_2
=\frac{\varepsilon\,|a-1|\,\|\mathbf{v}\|_2}{(a\|\mathbf{v}\|_2+\varepsilon)(\|\mathbf{v}\|_2+\varepsilon)}
=O\!\bigl(\varepsilon/\|\mathbf{v}\|_2\bigr).
\]
\end{lemma}
\begin{proof}
The mean satisfies $\overline{\mathbf{y}}=a\,\overline{\mathbf{x}(\alpha)}+b$, so $\mathbf{y}-\overline{\mathbf{y}}\mathbf{1}=a\mathbf{v}$ and $\mathbf{r}_{\mathbf{y}}=a\mathbf{v}/(a\|\mathbf{v}\|_2+\varepsilon)$. Hence
\[
\mathbf{r}_{\mathbf{y}}-\mathbf{r}(\alpha)
=\Bigl(\frac{a}{a\|\mathbf{v}\|_2+\varepsilon}-\frac{1}{\|\mathbf{v}\|_2+\varepsilon}\Bigr)\mathbf{v}
=\frac{\varepsilon\,(a-1)}{(a\|\mathbf{v}\|_2+\varepsilon)(\|\mathbf{v}\|_2+\varepsilon)}\,\mathbf{v},
\]
whose norm is the stated quantity. It vanishes as $\varepsilon\to 0$ (and for $a=1$); with $\varepsilon=0$ the normalisation cancels $a>0$ exactly.
\end{proof}

Thus, up to an error of order $\varepsilon$---which is fixed at machine-precision scale while $\|\mathbf{v}\|_2$ is bounded away from zero for any non-constant IC series---the descriptor encodes only which dates were strong; the mean and scale are absorbed into the quality score. For the same reason the descriptor is only approximately of unit norm, $\|\mathbf{r}(\alpha)\|_2=\|\mathbf{v}\|_2/(\|\mathbf{v}\|_2+\varepsilon)<1$, and ``unit-normalised'' is used in this $O(\varepsilon)$ sense throughout. Two alphas with the same temporal pattern but different overall edge occupy the same cell rather than leaking into different ones; a zero-variance IC series ($\|\mathbf{v}\|_2=0$) has no meaningful shape and is handled by the reward floor of \eqref{eq:alphaqd_reward} rather than by the descriptor.

\subsection{Reward Stationarity}
\label{sec:alphaqd_stationary}

The second result is the structural reason the static reward survives even archive-aware exploration extensions.

\begin{theorem}[Stationary conditional reward under a fixed descriptor map]
\label{thm:alphaqd_stationary}
Let the per-alpha reward be $R_q(\alpha)=\exp(\beta_q\,\mathrm{ICIR}(\alpha))$, and assume for this conditional extension that the PCA basis $W$ and the bin edges entering $c(\cdot)$ are fixed. For any fixed target orbit $\mathcal{O}^\star$ and any kernel $k:\mathcal{O}\times\mathcal{O}\to\mathbb{R}_{\ge 0}$ independent of the archive state and training step, the conditional reward $R_{\mathcal{O}^\star}(\alpha)=R_q(\alpha)\,k(\mathcal{O}(c(\alpha)),\mathcal{O}^\star)$ is a stationary terminal reward.
\end{theorem}
\begin{proof}
For each fixed $\mathcal{O}^\star$, $R_{\mathcal{O}^\star}(\cdot)$ depends on $\alpha$ only through $\mathrm{ICIR}(\alpha)$ on the fixed training window and through the fixed PCA basis and bin edges entering $c(\cdot)$. The kernel $k$ is fixed and $\mathcal{O}^\star$ is a parameter, not a function of training-time state. The reward therefore has no time-dependent component. A scheduler that varies $\mathcal{O}^\star$ across episodes changes only the per-episode prior over conditions, never the conditional reward itself.
\end{proof}

The fixed-map premise is \emph{not} supplied by the deployed schedule: the current code updates the PCA basis and edges every $100$ episodes when its buffer is large enough. The implementation uses the unconditional reward ($k\equiv 1$), whose stationarity is independent of the descriptor map and therefore survives those updates. The conditional result applies only to a hypothetical descriptor-biased extension with a fixed map; it is not a theorem about the deployed rolling-PCA partition.

\subsection{Search Complexity and a Boosting Counterexample}
\label{sec:alphaqd_complexity}

The next two results contrast AlphaQD's quality-diversity ensemble with boosting-style residual fitting in the abstract form used by the pool-coupled rewards of \sref{sec:alphaqd_formulation}; they concern this class of procedures, not the specific published AlphaGen system.

\begin{theorem}[Residual-edge noise floor]
\label{thm:alphaqd_boost}
Let $\mathcal{F}=\{\alpha_1,\dots,\alpha_M\}$ be a finite candidate set. After $t-1$ boosting rounds with residual $r_{t-1}$, define the population residual edge and its training-window estimate by
\[
\gamma_t(\alpha)=\mathbb{E}[\langle\alpha,r_{t-1}\rangle],
\qquad
\hat\gamma_t(\alpha)=\gamma_t(\alpha)+\xi_{t,\alpha}.
\]
Assume that each deviation is sub-Gaussian: for a scale parameter $\sigma$ and effective sample size $n_{\mathrm{eff}}$,
\[
\mathbb{P}(|\xi_{t,\alpha}|>u)
\le 2\exp\!\left(-\frac{n_{\mathrm{eff}}u^2}{2\sigma^2}\right).
\]
Let $\alpha_t=\argmax_{\alpha\in\mathcal{F}}\hat\gamma_t(\alpha)$, $\Delta_t=\sup_\alpha\gamma_t(\alpha)$, and $\eta_M=\sigma\sqrt{2\log(2M/\delta)/n_{\mathrm{eff}}}$. Then, with probability at least $1-\delta$, $\gamma_t(\alpha_t)\ge\Delta_t-2\eta_M$. When $\Delta_t\le 2\eta_M$, the bound is non-positive.
\end{theorem}
\begin{proof}
Uniform sub-Gaussian concentration on the finite class $\mathcal{F}$ gives $\sup_\alpha|\hat\gamma_t(\alpha)-\gamma_t(\alpha)|\le\eta_M$ with probability at least $1-\delta$. Writing $\alpha_t^\star=\argmax_\alpha\gamma_t(\alpha)$, $\gamma_t(\alpha_t)\ge\hat\gamma_t(\alpha_t)-\eta_M\ge\hat\gamma_t(\alpha_t^\star)-\eta_M\ge\gamma_t(\alpha_t^\star)-2\eta_M=\Delta_t-2\eta_M$, where the middle step uses the empirical optimality of $\alpha_t$.
\end{proof}

With $|\mathcal{F}|>10^{20}$ and $n_{\mathrm{eff}}$ bounded by the training-window days, $\eta_M$ is large relative to plausible late-round residual edges, so the guarantee $\Delta_t-2\eta_M$ turns non-positive and hence vacuous: the empirical residual winner retains no assurance of positive population edge. This is a statement about the emptiness of the guarantee, not an upper bound on the winner's edge; establishing that late-round winners are in fact noise-dominated would require a matching upper bound or an explicit winner's-curse argument, which we do not attempt.

\begin{theorem}[Search-complexity contrast]
\label{thm:alphaqd_complexity}
A $K$-step boosting procedure over $\mathcal{F}$ with $|\mathcal{F}|=M$ (here $K$ plays the same retained-set-size role as in \eqref{eq:alphaqd_goal}) has an ordered-sequence hypothesis class of size at most $M^K$, with finite-class complexity $\mathcal{C}_{\mathrm{boost},K}=\sqrt{K\log M/n_{\mathrm{eff}}}$. A quality-diversity rule mapping candidates to at most $G$ occupied cells and exporting one representative per cell has $K$-subset class of size at most $\binom{G}{K}$, with complexity $\mathcal{C}_{\mathrm{QD},K}=\sqrt{K\log(eG/K)/n_{\mathrm{eff}}}$. For $G\ll M$, $\mathcal{C}_{\mathrm{QD},K}\ll\mathcal{C}_{\mathrm{boost},K}$.
\end{theorem}
\begin{proof}
The finite-class generalisation bound scales as $\sqrt{\log|\mathcal{H}|/n_{\mathrm{eff}}}$. Substituting $|\mathcal{H}_{\mathrm{boost},K}|\le M^K$ gives $\log|\mathcal{H}_{\mathrm{boost},K}|\le K\log M$. Substituting $|\mathcal{H}_{\mathrm{QD},K}|\le\binom{G}{K}$ and $\log\binom{G}{K}\le K\log(eG/K)$ gives the QD scaling.
\end{proof}

In the reported grid $G\le 100$ while $|\mathcal{F}|>10^{20}$: the class sizes differ by many orders of magnitude, though the complexity terms themselves differ only by a modest factor (about $3.7$ under $M=10^{20}$, $G=100$, $K=10$). Two caveats govern the reading. First, the comparison is conditional on the realised archive: the count $\binom{G}{K}$ treats the per-cell representatives, the PCA basis, and the bin edges as fixed, whereas all three are data-dependent products of a search that itself screened a number of candidates comparable to what boosting sees, and the ridge weights add a continuous layer on top. Theorem~\ref{thm:alphaqd_complexity} therefore bounds the complexity of the final selection stage given the archive, not the effective complexity of the full pipeline. Second, the two class-size counts are upper bounds on differently structured procedures, so the contrast is indicative rather than a like-for-like risk comparison.

\subsection{Per-Cell Selection Bias}
\label{sec:alphaqd_selection}

When a cell aggregates several candidates and the elite is selected by maximum empirical ICIR, the chosen elite's empirical ICIR is upward-biased relative to its population value---the formula-mining analogue of a multiple-testing winner's curse~\cite{harvey2016cross}. The bounded per-cell quota controls the final, within-cell component of this bias.

\begin{theorem}[Per-cell selection bias under a bounded quota]
\label{thm:alphaqd_selection}
Fix an occupied cell holding at most $\kappa$ admitted elites at selection time, each with population quality $\vartheta_i$ and training-window estimate $\hat\vartheta_i=\vartheta_i+\varepsilon_i$, where each $\varepsilon_i$ is sub-Gaussian with $\mathbb{P}(\varepsilon_i>u)\le\exp(-n_{\mathrm{eff}}u^2/(2\sigma^2))$ for a scale parameter $\sigma$ and effective sample size $n_{\mathrm{eff}}$. With probability at least $1-\delta$, the selection bias of the per-cell representative on raw ICIR satisfies $\max_{i\le\kappa}(\hat\vartheta_i-\vartheta_i)\le\sigma\sqrt{2\log(\kappa/\delta)/n_{\mathrm{eff}}}$.
\end{theorem}
\begin{proof}
A union bound on the $\kappa$ sub-Gaussian tails gives $\mathbb{P}(\max_{i\le\kappa}\varepsilon_i>u)\le\kappa\exp(-n_{\mathrm{eff}}u^2/(2\sigma^2))$; setting the right-hand side to $\delta$ yields $u=\sigma\sqrt{2\log(\kappa/\delta)/n_{\mathrm{eff}}}$.
\end{proof}

Under the canonical $\kappa=3$, $\delta=0.05$ the bound is $\approx 2.86\,\sigma/\sqrt{n_{\mathrm{eff}}}$. Two caveats bound its reach. First, the deployed ranking uses raw ICIR with PFS only as a tie-break (\sref{sec:alphaqd_archive}); the five PFS perturbations share the same training window and are neither independent nor same-mean estimates of the elite's quality, so no $\sqrt{1+m}$-style tightening of the bound is claimed here---any such tightening would be heuristic. Second, and more importantly, the theorem conditions on the $\kappa$ stored elites as a fixed set, whereas they are themselves the survivors of an empirical-score tournament over every candidate the cell ever received. The quota therefore bounds the bias of the final within-cell $\argmax$, but it does not remove the selection bias accumulated during the search---the search-stage multiple-testing winner's curse~\cite{harvey2016cross} lies outside the theorem's scope.

\subsection{Ensemble ICIR Lower Bound}
\label{sec:alphaqd_ensemble}

The final result is an abstract lemma about equal-weight ensembles; its relation to the deployed pool is discussed below.

\begin{theorem}[Ensemble ICIR lower bound]
\label{thm:alphaqd_ensemble}
Let $X_1,\dots,X_K$ be daily IC variables satisfying $\mathbb{E}[X_i]\ge\mu_{\min}>0$, $\Var(X_i)\le\sigma^2$, and $\corr(X_i,X_j)\le\rho_{\max}$ for all $i\ne j$, for some $\rho_{\max}\in[0,1]$. Then the equal-weight ensemble $X_{\mathrm{ens}}=\frac{1}{K}\sum_i X_i$ obeys
\begin{equation}
\label{eq:alphaqd_ensemble}
\mathrm{ICIR}(X_{\mathrm{ens}}) \ge \frac{\mu_{\min}\sqrt{K}}{\sigma\sqrt{1+(K-1)\rho_{\max}}}.
\end{equation}
\end{theorem}
\begin{proof}
By linearity, $\mathbb{E}[X_{\mathrm{ens}}]\ge\mu_{\min}$. For the variance,
\begin{align*}
\Var(X_{\mathrm{ens}})
&=\frac{1}{K^2}\sum_i\Var(X_i)
  +\frac{1}{K^2}\sum_{i\ne j}\Cov(X_i,X_j)\\
&\le\frac{\sigma^2}{K^2}\bigl[K+K(K-1)\rho_{\max}\bigr]\\
&=\frac{\sigma^2\bigl(1+(K-1)\rho_{\max}\bigr)}{K}.
\end{align*}
so $\mathrm{ICIR}(X_{\mathrm{ens}})=\mathbb{E}[X_{\mathrm{ens}}]/\sqrt{\Var(X_{\mathrm{ens}})}\ge\mu_{\min}\sqrt{K}/(\sigma\sqrt{1+(K-1)\rho_{\max}})$.
\end{proof}

The proof also uses $\rho_{\max}\ge 0$: with a negative correlation cap the covariance inequality would reverse. Theorem~\ref{thm:alphaqd_ensemble} is a statement about any family of random variables satisfying its premises; three gaps separate it from a guarantee for the deployed pool. First, the admission rule enforces a positive \emph{training-window} ICIR, which makes the premise $\mathbb{E}[X_i]\ge\mu_{\min}>0$ plausible for the exported elites but does not imply it out of sample. Second, the premise requires a uniform cap $\rho_{\max}$ on all pairwise correlations of the exported components. Spreading elites across PCA cells encourages lower redundancy but does not enforce such a cap, and the archive figure reports no pairwise-correlation statistic. Third, the bound concerns the equal-weight average of the components' daily IC variables, whereas the deployed signal is a ridge-weighted combination of factor values whose IC is not a linear average of the component ICs. Theorem~\ref{thm:alphaqd_boost} shows that boosting-style residual fitting need not preserve even the first premise, so the analogous bound is not available for it by this route; for AlphaQD too, however, the result is motivation for the export design rather than a finite-sample guarantee for the deployed pool.

%----------------------------------------------------------------------------------------
\section{Experimental Setup}
\label{sec:alphaqd_setup}

\paragraph{Datasets.} We evaluate on two markets so the comparison is not anchored to one regime. The implementation is maintained in the AIFML project repository.\footnote{\url{https://github.com/qinmoelei/AIFML}} The Chinese A-share panel (the canonical alpha-mining setting) is the primary benchmark: the experiment record reports BaoStock daily data for $7{,}207$ instruments over $3{,}963$ trading days from 2010-01-04 to 2026-04-30, with front-adjusted OHLCV and $\mathrm{VWAP}=\mathrm{amount}/\mathrm{volume}$. The U.S.\ panel is an independent check with different microstructure and base rates: the record reports a Tiingo Russell-3000-style panel of $2{,}571$ instruments over $4{,}109$ trading days, using Tiingo's adjusted price fields. Tradability masks apply at execution: on A-shares an asset must have finite prices, positive volume, normal trading status, non-ST status, and an open-to-previous-close gap within $\pm 9.7\%$ (approximating price-limit constraints); on the U.S.\ panel, finite prices, positive volume, and a minimum close of one dollar (excluding penny names).

The datasets and preprocessing follow the experimental setting established by AlphaGen~\cite{yu2023alphagen}. Every method within a market receives the same instruments, dates, labels, execution rule, transaction costs, and portfolio construction protocol, enabling a controlled comparison of the alpha-mining methods.

\paragraph{Split and usage.} All methods share the same three-way chronological split: 2010--2021 for training, 2022 for validation, and 2023-01-01 to 2026-04-30 for testing, with forward-return horizon $h=20$ trading days. Following the AlphaGen protocol, the training set is used to search for or learn formulaic alphas. Once a method exports its alpha pool, the validation set is used to fit the shared ridge combiner that maps the pool to a scalar signal. The fitted linear weights are then fixed, and the test set is used only for final performance evaluation. This separation gives every pool-style method the same downstream combiner and isolates the quality of the generated alpha pool from method-specific weighting choices.

\paragraph{Random seeds.} Stochastic methods are evaluated using seeds 0, 1, 2, and 12345. We report the mean and population standard deviation across these runs; deterministic methods have zero run variance.

\paragraph{Compared methods.} We compare against eleven baselines in four families: formulaic (the fixed Alpha101 catalogue~\cite{kakushadze2016101}); supervised (MLP~\cite{murtagh1991mlp}, XGBoost~\cite{chen2016xgboost}, LightGBM~\cite{ke2017lightgbm}); genetic programming (GP single-predictor; GP with mutual-IC filter; GP without filter~\cite{stephens2015gplearn}); and generative/RL (AlphaForge~\cite{shi2025alphaforge}, AlphaGen-PPO~\cite{yu2023alphagen}, AlphaQCM~\cite{zhu2025alphaqcm}, AlphaSAGE~\cite{chen2025alphasage}). Pool-style baselines export ten alphas each, matching AlphaQD's $K^\star=10$.

\paragraph{Shared protocol.} Every pool-style method, including AlphaQD, exports a fixed-size pool. That pool is fed to the same closed-form $\lambda=1$ ridge combiner described in \sref{sec:alphaqd_export}, with weights fitted on validation data and held fixed on the test set. We report signal-layer IC, ICIR, and RankIC~\cite{grinold2000active} together with portfolio-layer annualised net Sharpe~\cite{sharpe1994sharpe}, annualised net return (the arithmetic mean daily net return scaled by $252$), maximum drawdown~\cite{magdon2004maximum} (reported as a positive magnitude; smaller is better), and annualised turnover. The portfolio is long-only top-$50$ rebalanced every twenty trading days, executed at $\mathrm{VWAP}_{t+1}$ on A-shares and $\mathrm{Open}_{t+1}$ on the U.S.\ panel, with round-trip costs of $12.5$ and $3.5$ basis points respectively. Applying this common protocol to every method makes the within-panel ranks the primary comparison.

%----------------------------------------------------------------------------------------
\section{Results and Analysis}
\label{sec:alphaqd_results}

\subsection{Main Results}
\label{sec:alphaqd_main}

\tref{tab:alphaqd_main} reports both panels under the unified-ridge protocol, aggregated over seeds $\{0,1,2,12345\}$ where stochastic.

\begin{table}[!htb]
\centering
\caption[AlphaQD main results.]{Main results on the 2023-01-01 to 2026-04-30 test window under the unified-ridge protocol. Entries are mean $\pm$ population standard deviation over seeds $\{0,1,2,12345\}$ where stochastic; deterministic methods carry zero seed variance. Portfolio metrics are net of round-trip costs on a long-only top-$50$ portfolio rebalanced every twenty trading days; maximum drawdown is reported as a positive magnitude (smaller is better). Per column within each panel, the \fst{best}, \snd{second-best}, and \trd{third-best} entries are marked in bold, underline, and dotted underline respectively (colour retained for screen reading).}
\label{tab:alphaqd_main}
\renewcommand{\arraystretch}{1.15}
\footnotesize
\resizebox{\textwidth}{!}{
\begin{tabular}{llccccccc}
\toprule
Family & Method & IC$\uparrow$ & ICIR$\uparrow$ & RankIC$\uparrow$ & Sharpe$\uparrow$ & AnnRet(\%)$\uparrow$ & MDD(\%)$\downarrow$ & TO/yr$\downarrow$ \\
\midrule
\multicolumn{9}{c}{\textit{Panel A: Chinese A-share}} \\
\midrule
Formulaic & Alpha101 & $0.0487\pm0.0000$ & $0.408\pm0.000$ & $0.082\pm0.000$ & \fst{$0.983\pm0.000$} & \snd{$19.9\pm0.0$} & \fst{$26.1\pm0.0$} & $23.9$ \\
\multirow{3}{*}{Supervised} & MLP & $0.0079\pm0.0017$ & $0.111\pm0.018$ & $0.017\pm0.003$ & $0.319\pm0.038$ & $9.1\pm1.3$ & $49.7\pm1.0$ & $12.9$ \\
 & LightGBM & $0.0189\pm0.0002$ & $0.262\pm0.006$ & $0.035\pm0.000$ & $0.309\pm0.022$ & $10.0\pm0.7$ & $52.7\pm1.4$ & \snd{$12.0$} \\
 & XGBoost & $0.0199\pm0.0000$ & $0.266\pm0.000$ & $0.036\pm0.000$ & $0.236\pm0.000$ & $6.9\pm0.0$ & $50.6\pm0.0$ & \trd{$12.4$} \\
\multirow{3}{*}{Genetic Prog.} & GP (predictor) & $-0.0101\pm0.0132$ & $-0.082\pm0.118$ & $-0.016\pm0.024$ & $0.292\pm0.309$ & $6.4\pm6.1$ & $40.4\pm16.4$ & \fst{$3.8$} \\
 & GP w/ filter & \trd{$0.0652\pm0.0080$} & $0.507\pm0.070$ & \snd{$0.105\pm0.010$} & $0.663\pm0.082$ & \trd{$18.6\pm3.9$} & $39.1\pm5.2$ & $15.1$ \\
 & GP w/o filter & \fst{$0.0703\pm0.0157$} & \fst{$0.581\pm0.064$} & \trd{$0.102\pm0.025$} & $0.501\pm0.573$ & $10.8\pm12.9$ & $35.7\pm9.4$ & $20.2$ \\
\multirow{4}{*}{Generative / RL} & AlphaForge & $0.0636\pm0.0099$ & \trd{$0.553\pm0.030$} & $0.086\pm0.016$ & $0.534\pm0.493$ & $14.4\pm13.4$ & $37.0\pm15.7$ & $19.5$ \\
 & AlphaGen-PPO & $0.0435\pm0.0215$ & $0.495\pm0.219$ & $0.049\pm0.043$ & $0.275\pm0.551$ & $4.1\pm11.3$ & $43.3\pm15.4$ & $22.2$ \\
 & AlphaQCM (FQF) & $0.0206\pm0.0187$ & $0.255\pm0.291$ & $0.020\pm0.035$ & $0.053\pm0.451$ & $1.7\pm12.2$ & $46.8\pm12.2$ & $19.3$ \\
 & AlphaSAGE (GNN) & $0.0552\pm0.0082$ & \snd{$0.555\pm0.142$} & $0.095\pm0.008$ & \trd{$0.730\pm0.050$} & $14.9\pm1.0$ & \snd{$27.4\pm3.9$} & $18.3$ \\
Ours & \textbf{AlphaQD} & \snd{$0.0673\pm0.0093$} & $0.522\pm0.078$ & \fst{$0.112\pm0.011$} & \snd{$0.911\pm0.185$} & \fst{$24.7\pm6.2$} & \trd{$31.5\pm1.7$} & $12.5$ \\
\midrule
\multicolumn{9}{c}{\textit{Panel B: U.S.\ Russell-3000-style universe}} \\
\midrule
Formulaic & Alpha101 & $0.0188\pm0.0000$ & $0.358\pm0.000$ & \snd{$-0.002\pm0.000$} & $1.511\pm0.000$ & $38.2\pm0.0$ & \snd{$15.9\pm0.0$} & $23.0$ \\
\multirow{3}{*}{Supervised} & MLP & $0.0446\pm0.0102$ & \fst{$0.604\pm0.147$} & \trd{$-0.002\pm0.007$} & $1.796\pm0.055$ & $99.8\pm5.0$ & \trd{$16.3\pm0.9$} & $9.1$ \\
 & LightGBM & $0.0404\pm0.0023$ & $0.454\pm0.025$ & $-0.009\pm0.001$ & $2.090\pm0.035$ & $119.1\pm2.3$ & $21.6\pm1.4$ & $8.7$ \\
 & XGBoost & $0.0335\pm0.0000$ & $0.395\pm0.000$ & $-0.006\pm0.000$ & $2.120\pm0.000$ & $129.2\pm0.0$ & $23.4\pm0.0$ & $9.6$ \\
\multirow{3}{*}{Genetic Prog.} & GP (predictor) & $-0.0104\pm0.0458$ & $-0.073\pm0.417$ & \fst{$0.004\pm0.010$} & $1.202\pm0.497$ & $39.4\pm45.8$ & \fst{$15.2\pm5.2$} & \fst{$5.7$} \\
 & GP w/ filter & \snd{$0.0691\pm0.0014$} & $0.512\pm0.069$ & $-0.005\pm0.003$ & \fst{$2.180\pm0.043$} & \fst{$141.6\pm6.4$} & $26.7\pm2.4$ & $9.5$ \\
 & GP w/o filter & \fst{$0.0717\pm0.0031$} & \snd{$0.568\pm0.013$} & $-0.013\pm0.001$ & \trd{$2.164\pm0.047$} & \snd{$141.3\pm5.5$} & $24.7\pm0.8$ & $10.2$ \\
\multirow{4}{*}{Generative / RL} & AlphaForge & $0.0481\pm0.0227$ & $0.430\pm0.181$ & $-0.008\pm0.002$ & $1.856\pm0.525$ & $107.0\pm41.3$ & $23.4\pm1.7$ & $11.6$ \\
 & AlphaGen-PPO & $0.0542\pm0.0082$ & $0.437\pm0.062$ & $-0.010\pm0.002$ & $2.056\pm0.223$ & $129.8\pm15.6$ & $23.3\pm1.1$ & $13.8$ \\
 & AlphaQCM (FQF) & $0.0486\pm0.0295$ & $0.453\pm0.286$ & $-0.002\pm0.009$ & $1.942\pm0.372$ & $103.8\pm46.4$ & $21.1\pm3.4$ & \trd{$7.9$} \\
 & AlphaSAGE (GNN) & $0.0614\pm0.0077$ & $0.539\pm0.068$ & $-0.007\pm0.003$ & $2.111\pm0.094$ & $131.3\pm7.4$ & $24.2\pm1.5$ & $10.5$ \\
Ours & \textbf{AlphaQD} & \trd{$0.0655\pm0.0022$} & \trd{$0.551\pm0.036$} & $-0.010\pm0.001$ & \snd{$2.177\pm0.020$} & \trd{$133.3\pm2.0$} & $22.5\pm0.2$ & \snd{$6.7$} \\
\bottomrule
\end{tabular}}
\end{table}

\paragraph{A-share findings.} On the A-share panel, AlphaQD attains the highest annualised return and the highest RankIC of any method. It is second on IC (behind GP-without-filter) and second on Sharpe (behind only the fixed Alpha101 catalogue). The leading signal-quality group comprises GP-without-filter, AlphaForge, AlphaSAGE, and AlphaQD. Within this group the portfolio columns differentiate more clearly: GP-without-filter and AlphaForge have substantially larger seed-to-seed variation in Sharpe and annualised return, while AlphaSAGE is stable but has an annualised return roughly ten percentage points below AlphaQD and turnover $46\%$ higher. Among the four generative/RL pool-style methods, AlphaQD has the lowest turnover and the smallest standard deviation in maximum drawdown. \fref{fig:alphaqd_nav} plots the cumulative net NAV of the five learned pool-style methods.

\begin{figure}[!htb]
\centering
\includegraphics[width=0.78\textwidth]{Chapter6/alphaqd_portfolio_nav}
\caption[Cumulative NAV on the A-share test window.]{Cumulative net NAV (start $=1.0$) on the A-share test window 2023-01-01 to 2026-04-30 for the five learned pool-style methods. Solid lines are the mean NAV across seeds $\{0,1,2,12345\}$ at each rebalance date; shaded bands span $\pm 1$ population standard deviation; the dashed line marks NAV $=1$. The portfolio is the long-only top-$50$ rebalanced every twenty trading days at $\mathrm{VWAP}_{t+1}$ with $12.5$ bps round-trip cost.}
\label{fig:alphaqd_nav}
\end{figure}

\paragraph{U.S.\ findings.} Under the shared U.S.\ protocol, AlphaQD ranks third on IC, ICIR, and annualised return, and second on net Sharpe and turnover. Among the four generative/RL pool-style methods it also has the smallest reported run-level deviation on every IC-family metric. These relative positions, rather than the absolute level of any one portfolio metric, are the intended cross-method comparison on this panel.

\subsection{Archive Diagnostics}
\label{sec:alphaqd_diagnostics}

The archive snapshot in \fref{fig:alphaqd_archive} contains $87$ occupied cells and $224$ stored elites. The largest representative ICIR annotated in the grid is $0.66$. These figure-derived values are used consistently as the descriptive archive result; the grid visualises broad coverage of the two-dimensional behaviour space, while it does not by itself establish a uniform correlation bound for the exported pool.

\begin{figure}[!htb]
\centering
\includegraphics[width=0.8\textwidth]{Chapter6/alphaqd_archive_heatmap}
\caption[The AlphaQD MAP-Elites archive.]{AlphaQD archive snapshot on the $10\times 10$ behaviour-descriptor grid (vertical axis PC1 bin, horizontal axis PC2 bin). \emph{(a)} number of stored elites per cell; \emph{(b)} representative ICIR per cell, with ten cells outlined. The plotted annotations contain $87$ occupied cells and $224$ elites; the largest annotated representative ICIR is $0.66$.}
\label{fig:alphaqd_archive}
\end{figure}

\subsection{Hyperparameter Sensitivity}
\label{sec:alphaqd_ablation}

A grid over behaviour-descriptor dimension $d_b\in\{1,2,3\}$, per-cell capacity $\kappa\in\{1,3,5\}$, and reward temperature $\beta_q\in\{0.5,1,2,5\}$ probes the sensitivity of the canonical configuration on the validation set. Marginal ICIR declines gently from $d_b=1$ to $d_b=3$, consistent with a higher-dimensional descriptor spreading candidates over more cells; the canonical $d_b=2$ preserves an interpretable two-axis archive while remaining in the favourable range. Capacity $\kappa\ge 3$ is preferred over $\kappa=1$. The reward-temperature response is non-monotone, reflecting the trade-off between broad exploration and concentration on high-quality formulas. The full $36$-configuration grid is reported in \tref{tab:alphaqd_ablation_full}; its purpose is to show the sensitivity of the method rather than to treat every grid cell as a statistically separated comparison.

This grid varies structural hyperparameters rather than removing complete modules. It should therefore be read as a sensitivity analysis, complementary to the archive diagnostics and main comparison, rather than as a component-by-component ablation. The descriptor dimension also changes the number of possible cells, while $K^\star=\min(10,5\kappa)$ links the cell capacity to the exported pool size; these coupled effects are considered when interpreting the marginal trends.

\begin{table}[!htb]
\centering
\caption[AlphaQD hyperparameter sensitivity.]{AlphaQD hyperparameter sensitivity on the A-share validation set under the unified-ridge protocol (single seed). The canonical configuration is $d_b=2$, $\kappa=3$, and $\beta_q=0.5$. Columns are descriptor dimension $d_b$, cell capacity $\kappa$, reward temperature $\beta_q$, signal IC, ICIR, Spearman rank-IC, and net Sharpe.}
\label{tab:alphaqd_ablation_full}
\renewcommand{\arraystretch}{1.05}
\footnotesize
\begin{tabular}{ccccccc}
\toprule
$d_b$ & $\kappa$ & $\beta_q$ & IC & ICIR & RankIC & Sharpe \\
\midrule
$1$ & $1$ & $0.5$ & $0.0565$ & $0.455$ & $0.074$ & $0.322$ \\
$1$ & $1$ & $1.0$ & $0.0584$ & $0.592$ & $0.097$ & $0.951$ \\
$1$ & $1$ & $2.0$ & $0.0772$ & $0.647$ & $0.127$ & $1.027$ \\
$1$ & $1$ & $5.0$ & $0.0328$ & $0.246$ & $0.045$ & $0.889$ \\
$1$ & $3$ & $0.5$ & $0.0740$ & $0.542$ & $0.123$ & $1.186$ \\
$1$ & $3$ & $1.0$ & $0.0577$ & $0.464$ & $0.094$ & $0.699$ \\
$1$ & $3$ & $2.0$ & $0.0757$ & $0.630$ & $0.106$ & $0.540$ \\
$1$ & $3$ & $5.0$ & $0.0760$ & $0.511$ & $0.127$ & $0.835$ \\
$1$ & $5$ & $0.5$ & $0.0669$ & $0.493$ & $0.096$ & $0.833$ \\
$1$ & $5$ & $1.0$ & $0.0658$ & $0.560$ & $0.100$ & $0.879$ \\
$1$ & $5$ & $2.0$ & $0.0683$ & $0.664$ & $0.107$ & $1.010$ \\
$1$ & $5$ & $5.0$ & $0.0486$ & $0.485$ & $0.077$ & $0.668$ \\
\midrule
$2$ & $1$ & $0.5$ & $0.0541$ & $0.461$ & $0.075$ & $0.874$ \\
$2$ & $1$ & $1.0$ & $0.0258$ & $0.218$ & $0.040$ & $0.909$ \\
$2$ & $1$ & $2.0$ & $0.0724$ & $0.662$ & $0.106$ & $1.175$ \\
$2$ & $1$ & $5.0$ & $0.0760$ & $0.646$ & $0.096$ & $0.738$ \\
$2$ & $3$ & $0.5$ & $0.0736$ & $0.537$ & $0.123$ & $0.955$ \\
$2$ & $3$ & $1.0$ & $0.0432$ & $0.366$ & $0.047$ & $0.907$ \\
$2$ & $3$ & $2.0$ & $0.0724$ & $0.541$ & $0.112$ & $0.718$ \\
$2$ & $3$ & $5.0$ & $0.0712$ & $0.565$ & $0.113$ & $0.869$ \\
$2$ & $5$ & $0.5$ & $0.0589$ & $0.529$ & $0.087$ & $0.976$ \\
$2$ & $5$ & $1.0$ & $0.0793$ & $0.613$ & $0.122$ & $0.892$ \\
$2$ & $5$ & $2.0$ & $0.0375$ & $0.511$ & $0.087$ & $1.005$ \\
$2$ & $5$ & $5.0$ & $0.0734$ & $0.578$ & $0.123$ & $1.141$ \\
\midrule
$3$ & $1$ & $0.5$ & $0.0699$ & $0.505$ & $0.109$ & $0.991$ \\
$3$ & $1$ & $1.0$ & $0.0318$ & $0.368$ & $0.058$ & $0.669$ \\
$3$ & $1$ & $2.0$ & $0.0502$ & $0.516$ & $0.080$ & $0.644$ \\
$3$ & $1$ & $5.0$ & $0.0669$ & $0.534$ & $0.083$ & $1.144$ \\
$3$ & $3$ & $0.5$ & $0.0756$ & $0.525$ & $0.127$ & $1.040$ \\
$3$ & $3$ & $1.0$ & $0.0603$ & $0.518$ & $0.078$ & $0.581$ \\
$3$ & $3$ & $2.0$ & $0.0554$ & $0.616$ & $0.085$ & $0.957$ \\
$3$ & $3$ & $5.0$ & $0.0355$ & $0.405$ & $0.048$ & $-0.008$ \\
$3$ & $5$ & $0.5$ & $0.0314$ & $0.352$ & $0.031$ & $0.422$ \\
$3$ & $5$ & $1.0$ & $0.0741$ & $0.562$ & $0.130$ & $1.025$ \\
$3$ & $5$ & $2.0$ & $0.0494$ & $0.414$ & $0.087$ & $0.584$ \\
$3$ & $5$ & $5.0$ & $0.0306$ & $0.342$ & $0.044$ & $0.847$ \\
\bottomrule
\end{tabular}
\end{table}

\paragraph{Computational cost.} A full training run of AlphaQD on the A-share panel---$2{,}000$ GFlowNet trajectory-balance episodes with mini-batch size $64$, archive maintenance, and PFS evaluation---completes in approximately five to six hours on a single NVIDIA RTX~4090 GPU. The dominant cost is the per-candidate backtest of the IC series on the training panel. With a fixed grid and per-cell quota, an ordinary cell lookup and insertion is $O(1)$; periodic PCA updates additionally reproject the stored elites. After training, the exported $K^\star=10$ alphas and validation-fitted ridge combiner are fixed.

%----------------------------------------------------------------------------------------
\section{Chapter Summary}
\label{sec:alphaqd_summary}

This chapter presented AlphaQD, a quality-diversity RL method for formulaic alpha mining. Its central choice is structural: the running pool is removed from the generator's reward, and diversity, robustness, and selection are handled by a MAP-Elites archive whose behaviour descriptor is the normalised daily IC series---an empirical proxy for a formula's return source. The training-window reward is thereby invariant to archive updates. Conditional on a realised archive, the final cell-selection class scales with the number of occupied cells rather than the cardinality of the formula space. On the A-share panel AlphaQD attains the highest annualised return and net Sharpe among learned formula-mining methods; among the four generative/RL pool-style methods it has the lowest turnover on both panels. On the U.S.\ panel it ranks third on IC, ICIR, and annualised return and second on net Sharpe and turnover. In the arc of the thesis, AlphaQD removes the aggregation mechanism from the learning reward entirely, complementing the learned router of EarnHFT and the VAE-based router of FineFT.
 % Conclusion and Future Work


%----------------------------------------------------------------------------------------
%	THESIS CONTENT - APPENDICES
%----------------------------------------------------------------------------------------

\addtocontents{toc}{\vspace{0.8em}} % Add a gap in the Contents, for aesthetics

\appendix % Cue to tell LaTeX that the following 'chapters' are Appendices

% Include the appendices of the thesis as separate files from the Appendices folder
% Uncomment the lines as you write the Appendices

%!TEX root=../mythesis.tex
% Appendix A

\chapter{Proofs and Derivations} % Main appendix title
\label{AppendixA}

This appendix collects the formal results referenced from the technical chapters: the convergence of the EarnHFT Q-teacher (\sref{app:earnhft_proof}), the convergence of the FineFT selective update (\sref{app:fineft_proof}), and the generalisation properties of the AlphaQD formulation (\sref{app:alphaqd_theory}).

%----------------------------------------------------------------------------------------
\section{Convergence of the EarnHFT Q-Teacher}
\label{app:earnhft_proof}

We prove Theorem~\ref{thm:earnhft_convergence} in its tabular idealisation: on the historical-replay MDP whose state is the pair $(t,p)$ of timestamp and position---the space on which the teacher table of Algorithm~\ref{alg:earnhft_qtable} is defined---with discount $\gamma\in(0,1)$, the action value learned with the idealised supervised update converges almost surely to $Q^\star$, provided the step sizes satisfy the Robbins--Monro conditions ($\sum_k a_k=\infty$, $\sum_k a_k^2<\infty$), every state--action pair is visited infinitely often, and the update noise has bounded variance. The analysis concerns the linear advantage-matching operator below, not the KL-form penalty of \eqref{eq:earnhft_loss} itself; no formal correspondence between the decaying coefficient $\alpha_{\mathrm{KL}}$ of \eqref{eq:earnhft_loss} and the fixed mixing weight $\lambda$ below is established. The method is based on $Q$-learning, so following the standard convergence argument~\cite{jaakkola1993convergence} it suffices, under the stated stochastic-approximation conditions, to show that the supervised update operator $\mathcal{T}$ is a contraction in the sup-norm whose fixed point is $Q^\star$. The operator combines the usual Bellman backup with the supervision term and is
\begin{equation}
\label{eq:app_earnhft_operator}
\mathcal{T}q(x,a) = \lambda\sum_{y\in X}T(y\mid x,a)\bigl(r(x,a,y)+\gamma\max_b q(y,b)\bigr) + (1-\lambda)\Bigl(\tfrac{1}{|A|}\sum_{a\in A}q(x,a)+Re(x,a)\Bigr),
\end{equation}
where the supervision residual is
\begin{equation}
\label{eq:app_earnhft_residual}
Re(x,a) = q^\star(x,a) - \tfrac{1}{|A|}\sum_{a\in A}q^\star(x,a),
\end{equation}
and $\lambda\in(0,1]$ weights the Bellman backup against the supervision term.

\begin{proposition}[$Q^\star$ is a fixed point]
\label{prop:app_earnhft_fixed}
The optimal action value $q^\star$ satisfies $\mathcal{T}q^\star = q^\star$.
\end{proposition}
\begin{proof}
Substituting $q^\star$ into \eqref{eq:app_earnhft_operator} and using the Bellman optimality equation $q^\star(x,a)=\sum_y T(y\mid x,a)(r(x,a,y)+\gamma\max_b q^\star(y,b))$ for the first term and \eqref{eq:app_earnhft_residual} for the second,
\begin{align*}
\mathcal{T}q^\star(x,a) &= \lambda\, q^\star(x,a) + (1-\lambda)\Bigl(\tfrac{1}{|A|}\textstyle\sum_a q^\star(x,a) + q^\star(x,a) - \tfrac{1}{|A|}\textstyle\sum_a q^\star(x,a)\Bigr) \\
&= \lambda\, q^\star(x,a) + (1-\lambda)\, q^\star(x,a) = q^\star(x,a). \qedhere
\end{align*}
\end{proof}

\begin{proposition}[Contraction in the sup-norm]
\label{prop:app_earnhft_contraction}
Let $\gamma\in(0,1)$. For any action values $q_1,q_2$, $\|\mathcal{T}q_1-\mathcal{T}q_2\|_\infty \le k\,\|q_1-q_2\|_\infty$ with $k=\lambda\gamma+(1-\lambda)=1-\lambda(1-\gamma)<1$.
\end{proposition}
\begin{proof}
The supervision residual $Re(x,a)$ depends only on $q^\star$, not on the argument of $\mathcal{T}$, so it cancels in the difference $\mathcal{T}q_1-\mathcal{T}q_2$. Writing the difference and bounding the $\max$ operator by $|\max_b q_1(y,b)-\max_b q_2(y,b)|\le\max_b|q_1(y,b)-q_2(y,b)|$,
\begin{align*}
\|\mathcal{T}q_1-\mathcal{T}q_2\|_\infty
&= \Bigl\| \lambda\textstyle\sum_y T(y\mid x,a)\gamma\bigl(\max_{b}q_1(y,b)-\max_{b}q_2(y,b)\bigr) \\
&\qquad + (1-\lambda)\tfrac{1}{|A|}\textstyle\sum_{a'}\bigl(q_1(x,a')-q_2(x,a')\bigr)\Bigr\|_\infty \\
&\le \lambda\gamma\,\|q_1-q_2\|_\infty + (1-\lambda)\,\|q_1-q_2\|_\infty \\
&= \bigl(\lambda\gamma+(1-\lambda)\bigr)\|q_1-q_2\|_\infty = \bigl(1-\lambda(1-\gamma)\bigr)\|q_1-q_2\|_\infty.
\end{align*}
Since $\gamma\in(0,1)$ and $\lambda\in(0,1]$, the factor $k=1-\lambda(1-\gamma)<1$, so $\mathcal{T}$ is a contraction. When $\gamma=1$ the factor equals $1$ and $\mathcal{T}$ is merely non-expansive; the undiscounted case is instead handled directly by the finite-horizon backward induction of Algorithm~\ref{alg:earnhft_qtable}.
\end{proof}

Combining the two propositions, $\mathcal{T}$ is a contraction with unique fixed point $q^\star$; under the Robbins--Monro step-size, infinite-visitation, and bounded-noise-variance conditions stated above, the stochastic-approximation argument of~\cite{jaakkola1993convergence} applies and the iterates converge almost surely to $Q^\star$, which proves Theorem~\ref{thm:earnhft_convergence}. The deployed implementation, which uses function approximation, lies outside this tabular idealisation; the result should be read as motivation rather than a guarantee for the deployed algorithm.

%----------------------------------------------------------------------------------------
\section{Convergence of the FineFT Selective Update}
\label{app:fineft_proof}

We argue that the selective-update rule of \sref{sec:fineft_stage1} preserves convergence of each learner on the dynamic it specialises in. The argument has two parts.

First, the pre-training loss \eqref{eq:fineft_pretrain} updates every learner with the same optimal-value supervisor used by EarnHFT (now constrained by the liquidation mask of Algorithm~\ref{alg:fineft_coav}). The per-learner update operator is exactly the operator $\mathcal{T}$ of \eqref{eq:app_earnhft_operator}, so by \sref{app:earnhft_proof} each learner's action value converges, in the same tabular idealisation and under the same stochastic-approximation conditions, to the constrained optimal action value on the demonstration distribution.

Second, consider the selective-update phase restricted to transitions from a single dynamic $z$. By the DP-MDP assumption (\sref{sec:fineft_dpmdp}), these transitions are generated by a stationary MDP $M_z$. The weight matrix of Algorithm~\ref{alg:fineft_weight} assigns, for each transition, a non-negative weight to the best-fitting learner $i^\star$ and its neighbours, and zero elsewhere; the diagonal weight of $i^\star$ is one. The update of learner $i^\star$ on transitions of dynamic $z$ is therefore a standard Huber-TD update under $M_z$ with a non-negative supervision term, whose update operator is a sup-norm contraction by the same argument as \sref{app:earnhft_proof}. The positive-feedback property---that specialising in $z$ lowers learner $i^\star$'s TD error on $z$, which raises its selection weight on subsequent $z$-transitions---ensures that, once a learner becomes the best fit for $z$, it continues to receive the $z$-transitions that drive its convergence on $M_z$. The cross-learner weights of neighbouring learners are bounded by the diagonal weights and decay with index distance, so they perturb the update by a contraction-preserving convex combination and do not break the contraction. Hence each specialist converges to the optimal action value of the dynamic it is selected for.

The full formal treatment, including the precise stochastic-approximation conditions on the selection weights, follows the convergence analysis of the original FineFT manuscript.

%----------------------------------------------------------------------------------------
\section{Generalisation Properties of AlphaQD}
\label{app:alphaqd_theory}

This section gives the full statements and proofs of the results summarised in \sref{sec:alphaqd_theory}. Two assumptions are used by name. \textbf{(A1)} \emph{Sign-symmetric warm-up shape distribution}: the warm-up sample $\{\mathbf{r}(\alpha_i)\}_{i=1}^M$ has zero mean and a second moment equal to that of its sign-augmented version $\{\mathbf{r}(\alpha_i),-\mathbf{r}(\alpha_i)\}$, which holds because a formula and its sign-flip are both legal under the same mask and produce IC series differing only by a global sign. \textbf{(A2)} \emph{Positive ridge weights on retained orbit pairs}: the shared ridge combiner assigns strictly positive weight to both representatives of an orbit when both are exported, which holds in the empirical regime where $\ell_2$-regularised regression on positively-edged factors places weight on both inputs.

\subsection{Decoupling Quality from Shape}
\label{app:alphaqd_decouple}

\begin{lemma}[Affine invariance of the shape descriptor; restates Lemma~\ref{lem:alphaqd_affine}]
\label{app:lem_affine}
For any candidate $\alpha$ and any positive affine transform $\mathbf{y}=a\,\mathbf{x}(\alpha)+b\mathbf{1}$ with $a>0$, the normalised shape vector built from $\mathbf{y}$ equals $\mathbf{r}(\alpha)$.
\end{lemma}
\begin{proof}
The mean satisfies $\overline{\mathbf{y}}=a\,\overline{\mathbf{x}(\alpha)}+b$, so $\mathbf{y}-\overline{\mathbf{y}}\mathbf{1}=a(\mathbf{x}(\alpha)-\overline{\mathbf{x}(\alpha)}\mathbf{1})$, and normalising by the $\ell_2$ norm cancels $a>0$, giving $\mathbf{r}(\alpha)$.
\end{proof}

The consequence is that the descriptor encodes only \emph{which} dates were relatively strong, not the overall edge: mean and scale are absorbed into the ICIR, removing the leakage that would let two alphas with the same temporal pattern but different edge occupy different cells.

\subsection{Reward Stationarity}
\label{app:alphaqd_stationary}

\begin{theorem}[Stationary conditional reward; restates Theorem~\ref{thm:alphaqd_stationary}]
\label{app:thm_stationary}
Let $R_q(\alpha)=\exp(\beta_q\,\mathrm{ICIR}(\alpha))$. For any fixed target orbit $\mathcal{O}^\star$ and any kernel $K:\mathcal{O}\times\mathcal{O}\to\mathbb{R}_{\ge 0}$ independent of the archive state and the training-step index, the conditional reward $R_{\mathcal{O}^\star}(\alpha)=R_q(\alpha)\,K(\mathcal{O}(c(\alpha)),\mathcal{O}^\star)$ is a stationary terminal reward.
\end{theorem}
\begin{proof}
For each fixed $\mathcal{O}^\star$, $R_{\mathcal{O}^\star}(\cdot)$ depends on $\alpha$ only through $\mathrm{ICIR}(\alpha)$ on the fixed training window and through the fixed PCA basis and bin edges entering $c(\cdot)$. The kernel $K$ is fixed and $\mathcal{O}^\star$ is a parameter, not a function of training-time state, so $R_{\mathcal{O}^\star}$ has no time-dependent component. A scheduler that varies $\mathcal{O}^\star$ across episodes changes only the per-episode prior over conditions, never the conditional reward.
\end{proof}

The implementation uses the unconditional reward ($K\equiv 1$), so the property holds a fortiori; the conditional form documents compatibility with future descriptor-biased exploration extensions.

\subsection{Search Complexity and the Boosting Counter-Example}
\label{app:alphaqd_complexity}

\begin{theorem}[Residual-edge noise floor]
\label{app:thm_boost}
Let $\mathcal{F}=\{\alpha_1,\dots,\alpha_M\}$ be a finite candidate set. After $t-1$ boosting rounds with residual $r_{t-1}$, the empirical residual-edge maximiser is $\alpha_t=\argmax_{\alpha\in\mathcal{F}}\widehat{\langle\alpha,r_{t-1}\rangle}$. Let $\gamma_t(\alpha)=\mathbb{E}[\langle\alpha,r_{t-1}\rangle]$ be the population residual edge, $\Delta_t=\sup_\alpha\gamma_t(\alpha)$, and $\eta_M=\sigma\sqrt{2\log(2M/\delta)/n_{\mathrm{eff}}}$ a sub-Gaussian deviation. Then with probability at least $1-\delta$, $\gamma_t(\alpha_t)\ge\Delta_t-2\eta_M$; when $\Delta_t\le 2\eta_M$, the bound is non-positive.
\end{theorem}
\begin{proof}
Uniform sub-Gaussian concentration over $\mathcal{F}$ gives $\sup_\alpha|\hat\gamma_t(\alpha)-\gamma_t(\alpha)|\le\eta_M$ with probability $\ge 1-\delta$. Writing $\alpha_t^\star=\argmax_\alpha\gamma_t(\alpha)$, $\gamma_t(\alpha_t)\ge\hat\gamma_t(\alpha_t)-\eta_M\ge\hat\gamma_t(\alpha_t^\star)-\eta_M\ge\gamma_t(\alpha_t^\star)-2\eta_M=\Delta_t-2\eta_M$, where the middle step uses the empirical optimality of $\alpha_t$.
\end{proof}

With $M>10^{20}$ and $n_{\mathrm{eff}}$ bounded by training-window days, the late-round residual winner is governed by the noise floor $\eta_M$ rather than by population signal.

\begin{theorem}[Search-complexity contrast; restates Theorem~\ref{thm:alphaqd_complexity}]
\label{app:thm_complexity}
A $K$-step boosting procedure over $\mathcal{F}$ with $|\mathcal{F}|=M$ has an ordered-sequence hypothesis class of size at most $M^K$, with finite-class complexity $\mathcal{C}_{\mathrm{boost},K}=\sqrt{K\log M/n_{\mathrm{eff}}}$. A quality-diversity rule mapping candidates to at most $G$ occupied cells and exporting one representative per cell has a $K$-subset class of size at most $\binom{G}{K}$, with complexity $\mathcal{C}_{\mathrm{QD},K}=\sqrt{K\log(G/K)/n_{\mathrm{eff}}}$. For $G\ll M$, $\mathcal{C}_{\mathrm{QD},K}\ll\mathcal{C}_{\mathrm{boost},K}$.
\end{theorem}
\begin{proof}
The finite-class generalisation bound scales as $\sqrt{\log|\mathcal{H}|/n_{\mathrm{eff}}}$. Substituting $|\mathcal{H}_{\mathrm{boost},K}|\le M^K$ gives $\log|\mathcal{H}_{\mathrm{boost},K}|\le K\log M$. Substituting $|\mathcal{H}_{\mathrm{QD},K}|\le\binom{G}{K}$ and $\log\binom{G}{K}\le K\log(eG/K)$ gives the QD scaling.
\end{proof}

In the reported grid $G\le 100$ while $|\mathcal{A}|>10^{20}$, so the gap spans several orders of magnitude.

\subsection{Per-Cell Selection Bias}
\label{app:alphaqd_selection}

\begin{theorem}[Per-cell selection bias under a bounded quota]
\label{app:thm_selection}
Fix an occupied cell holding at most $\kappa$ admitted elites at selection time, each with population quality $\theta_i$ and training-window estimate $\hat\theta_i=\theta_i+\varepsilon_i$ with sub-Gaussian $\varepsilon_i$ of effective sample size $n_{\mathrm{eff}}$. With probability at least $1-\delta$, the selection bias of the per-cell representative on raw ICIR satisfies $\max_{i\le\kappa}(\hat\theta_i-\theta_i)\le\sigma\sqrt{2\log(\kappa/\delta)/n_{\mathrm{eff}}}$. If the ranking uses the PFS-augmented score averaging the ICIR with $m$ independent perturbations, $\sigma$ is replaced by $\sigma/\sqrt{1+m}$.
\end{theorem}
\begin{proof}
A union bound on the $\kappa$ sub-Gaussian tails gives $\mathbb{P}(\max_{i\le\kappa}\varepsilon_i>t)\le\kappa\exp(-n_{\mathrm{eff}}t^2/(2\sigma^2))$; setting the right-hand side to $\delta$ yields $t=\sigma\sqrt{2\log(\kappa/\delta)/n_{\mathrm{eff}}}$. Averaging $1+m$ independent sub-Gaussian estimates with parameter $\sigma$ yields parameter $\sigma/\sqrt{1+m}$.
\end{proof}

The bounded quota replaces the candidate-count exponent that would govern an unbounded per-cell maximum by the constant $\kappa$; under $\kappa=3$, $\delta=0.05$ the bound is $\approx 2.86\,\sigma/\sqrt{n_{\mathrm{eff}}}$, and the $m=5$ perturbations tighten it by $\sqrt{6}$ when the PFS-augmented score decides the ranking.

\subsection{Ensemble ICIR Lower Bound}
\label{app:alphaqd_ensemble}

\begin{theorem}[Ensemble ICIR lower bound; restates Theorem~\ref{thm:alphaqd_ensemble}]
\label{app:thm_ensemble}
If the export retains $K$ positive-ICIR elites with daily IC variables $X_1,\dots,X_K$ satisfying $\mathbb{E}[X_i]=\mu_i\ge\mu_{\min}>0$, $\Var(X_i)\le\sigma^2$, and $\corr(X_i,X_j)\le\rho_{\max}$, then the equal-weight ensemble $X_{\mathrm{ens}}=\frac{1}{K}\sum_i X_i$ obeys $\mathrm{ICIR}(X_{\mathrm{ens}})\ge\mu_{\min}\sqrt{K}/(\sigma\sqrt{1+(K-1)\rho_{\max}})$.
\end{theorem}
\begin{proof}
By linearity, $\mathbb{E}[X_{\mathrm{ens}}]\ge\mu_{\min}$. For the variance,
\[
\Var(X_{\mathrm{ens}}) = \tfrac{1}{K^2}\textstyle\sum_i\Var(X_i) + \tfrac{1}{K^2}\sum_{i\ne j}\Cov(X_i,X_j) \le \tfrac{\sigma^2}{K^2}\bigl[K+K(K-1)\rho_{\max}\bigr] = \tfrac{\sigma^2(1+(K-1)\rho_{\max})}{K},
\]
so $\mathrm{ICIR}(X_{\mathrm{ens}})=\mathbb{E}[X_{\mathrm{ens}}]/\sqrt{\Var(X_{\mathrm{ens}})}\ge\mu_{\min}\sqrt{K}/(\sigma\sqrt{1+(K-1)\rho_{\max}})$.
\end{proof}

The bound rests on positive standalone edges, enforced by the positive-ICIR admission rule, and on controlled redundancy $\rho_{\max}$, enforced by the bounded per-cell quota and the PCA descriptor; \sref{sec:alphaqd_diagnostics} reports an empirical pairwise correlation of $0.2478$. Theorem~\ref{app:thm_boost} shows that boosting-style fitting preserves neither premise, so no analogous bound is available for it.

%----------------------------------------------------------------------------------------
\section{Full AlphaQD Ablation Grid}
\label{app:alphaqd_ablation}

\tref{tab:app_ablation_full} reports the full $36$-configuration ablation grid summarised by the marginal means in \sref{sec:alphaqd_ablation}. All entries are evaluated on the A-share test window under the unified-ridge protocol with a single training seed and the canonical training budget ($2{,}000$ episodes). The only varied hyperparameters are descriptor dimension $d_b$, cell capacity $\kappa$, and reward temperature $\beta_q$. Within-grid noise (ICIR roughly $0.22$--$0.66$) is larger than the marginal-mean effects reported in \sref{sec:alphaqd_ablation}, so individual grid cells should be read as samples from the AlphaQD performance distribution rather than as separated comparisons.

\begin{table}[!htb]
\centering
\caption[Full AlphaQD ablation grid.]{AlphaQD ablation: per-configuration metrics on the A-share panel under the unified-ridge protocol (single seed, test window). Columns are descriptor dimension $d_b$, cell capacity $\kappa$, reward temperature $\beta_q$, signal IC, ICIR, Spearman rank-IC, and net Sharpe.}
\label{tab:app_ablation_full}
\renewcommand{\arraystretch}{1.05}
\footnotesize
\begin{tabular}{ccccccc}
\toprule
$d_b$ & $\kappa$ & $\beta_q$ & IC & ICIR & RankIC & Sharpe \\
\midrule
$1$ & $1$ & $0.5$ & $0.0565$ & $0.455$ & $0.074$ & $0.322$ \\
$1$ & $1$ & $1.0$ & $0.0584$ & $0.592$ & $0.097$ & $0.951$ \\
$1$ & $1$ & $2.0$ & $0.0772$ & $0.647$ & $0.127$ & $1.027$ \\
$1$ & $1$ & $5.0$ & $0.0328$ & $0.246$ & $0.045$ & $0.889$ \\
$1$ & $3$ & $0.5$ & $0.0740$ & $0.542$ & $0.123$ & $1.186$ \\
$1$ & $3$ & $1.0$ & $0.0577$ & $0.464$ & $0.094$ & $0.699$ \\
$1$ & $3$ & $2.0$ & $0.0757$ & $0.630$ & $0.106$ & $0.540$ \\
$1$ & $3$ & $5.0$ & $0.0760$ & $0.511$ & $0.127$ & $0.835$ \\
$1$ & $5$ & $0.5$ & $0.0669$ & $0.493$ & $0.096$ & $0.833$ \\
$1$ & $5$ & $1.0$ & $0.0658$ & $0.560$ & $0.100$ & $0.879$ \\
$1$ & $5$ & $2.0$ & $0.0683$ & $0.664$ & $0.107$ & $1.010$ \\
$1$ & $5$ & $5.0$ & $0.0486$ & $0.485$ & $0.077$ & $0.668$ \\
\midrule
$2$ & $1$ & $0.5$ & $0.0541$ & $0.461$ & $0.075$ & $0.874$ \\
$2$ & $1$ & $1.0$ & $0.0258$ & $0.218$ & $0.040$ & $0.909$ \\
$2$ & $1$ & $2.0$ & $0.0724$ & $0.662$ & $0.106$ & $1.175$ \\
$2$ & $1$ & $5.0$ & $0.0760$ & $0.646$ & $0.096$ & $0.738$ \\
$2$ & $3$ & $0.5$ & $0.0736$ & $0.537$ & $0.123$ & $0.955$ \\
$2$ & $3$ & $1.0$ & $0.0432$ & $0.366$ & $0.047$ & $0.907$ \\
$2$ & $3$ & $2.0$ & $0.0724$ & $0.541$ & $0.112$ & $0.718$ \\
$2$ & $3$ & $5.0$ & $0.0712$ & $0.565$ & $0.113$ & $0.869$ \\
$2$ & $5$ & $0.5$ & $0.0589$ & $0.529$ & $0.087$ & $0.976$ \\
$2$ & $5$ & $1.0$ & $0.0793$ & $0.613$ & $0.122$ & $0.892$ \\
$2$ & $5$ & $2.0$ & $0.0375$ & $0.511$ & $0.087$ & $1.005$ \\
$2$ & $5$ & $5.0$ & $0.0734$ & $0.578$ & $0.123$ & $1.141$ \\
\midrule
$3$ & $1$ & $0.5$ & $0.0699$ & $0.505$ & $0.109$ & $0.991$ \\
$3$ & $1$ & $1.0$ & $0.0318$ & $0.368$ & $0.058$ & $0.669$ \\
$3$ & $1$ & $2.0$ & $0.0502$ & $0.516$ & $0.080$ & $0.644$ \\
$3$ & $1$ & $5.0$ & $0.0669$ & $0.534$ & $0.083$ & $1.144$ \\
$3$ & $3$ & $0.5$ & $0.0756$ & $0.525$ & $0.127$ & $1.040$ \\
$3$ & $3$ & $1.0$ & $0.0603$ & $0.518$ & $0.078$ & $0.581$ \\
$3$ & $3$ & $2.0$ & $0.0554$ & $0.616$ & $0.085$ & $0.957$ \\
$3$ & $3$ & $5.0$ & $0.0355$ & $0.405$ & $0.048$ & $-0.008$ \\
$3$ & $5$ & $0.5$ & $0.0314$ & $0.352$ & $0.031$ & $0.422$ \\
$3$ & $5$ & $1.0$ & $0.0741$ & $0.562$ & $0.130$ & $1.025$ \\
$3$ & $5$ & $2.0$ & $0.0494$ & $0.414$ & $0.087$ & $0.584$ \\
$3$ & $5$ & $5.0$ & $0.0306$ & $0.342$ & $0.044$ & $0.847$ \\
\bottomrule
\end{tabular}
\end{table}

%\input{./Appendices/AppendixB}

\cleardoublepage    % start a new page

%----------------------------------------------------------------------------------------
%	AUTHOR'S PUBLICATIONS
%----------------------------------------------------------------------------------------

\authorpublications{

\noindent The following publications and manuscripts arose from the research conducted during this doctoral candidature. Author lists are abbreviated; the candidate is the leading author of each. (Co-author names and venues are to be finalised by the candidate.)

\section*{Conference Proceedings}

\begin{itemize}
  \item \textbf{Molei Qin} et al., ``EarnHFT: Efficient Hierarchical Reinforcement Learning for High Frequency Trading,'' in \emph{Proceedings of the AAAI Conference on Artificial Intelligence (AAAI)}, 2024. (Chapter~\ref{ch:earnhft})
\end{itemize}

\section*{Manuscripts Under Review}

\begin{itemize}
  \item \textbf{Molei Qin} et al., ``FineFT: Efficient and Risk-Aware Ensemble Reinforcement Learning for Futures Trading,'' submitted to \emph{ACM SIGKDD Conference on Knowledge Discovery and Data Mining (KDD)}, 2026. (Chapter~\ref{ch:fineft})
  \item \textbf{Molei Qin} et al., ``AlphaQD: Quality-Diversity Reinforcement Learning for Formulaic Alpha Mining,'' submitted to \emph{ACM SIGKDD Conference on Knowledge Discovery and Data Mining (KDD)}. (Chapter~\ref{ch:alphaqd})
\end{itemize}

}

\backmatter

%----------------------------------------------------------------------------------------
%	BIBLIOGRAPHY
%----------------------------------------------------------------------------------------

\label{Bibliography}
\setstretch{1}
\bibliographystyle{unsrtnat}
\bibliography{References/thesis}

\end{document}
